% ===========================================================================
%  PLATAFORMA UNOORIA — Comparativo de Infraestrutura em Nuvem (GCP x AWS x Azure)
%  Documento autonomo. Compilar a partir desta pasta:
%      tectonic -k comparativo-infraestrutura-nuvem.tex
% ===========================================================================
\documentclass[11pt, a4paper, oneside, openany]{report}

% ============================================================================
%  estilo.tex — identidade visual, tipografia e componentes do documento
%
%  Tipografia: Roboto (texto e titulos) + Roboto Mono (codigo). As duas vem do
%  pacote `roboto`, entao o documento compila com pdfLaTeX, XeLaTeX ou Tectonic
%  sem fontspec explicito.
%
%  Paleta: os tokens `--unooria-*` do @unooria/design-system, os MESMOS valores
%  que `packages/ui/styles/plataforma.css` usa como fallback. Cor nova neste
%  documento so entra se existir como token no pacote.
% ============================================================================

% --- Codificacao e idioma ---------------------------------------------------
\usepackage[T1]{fontenc}
\usepackage[utf8]{inputenc}
\usepackage[shorthands=off,brazil]{babel}   % shorthands OFF: com babel-pt o `"` fica ativo e quebra codigo em linha
\usepackage{csquotes}

% --- Tipografia -------------------------------------------------------------
\usepackage[sfdefault]{roboto}            % Roboto como familia padrao
\usepackage[scaled=0.90]{roboto-mono}     % Roboto Mono para codigo
\usepackage{microtype}
\usepackage{textcomp}

% --- Pagina -----------------------------------------------------------------
\usepackage[
  a4paper,
  top=2.6cm, bottom=2.5cm, left=2.6cm, right=2.4cm,
  headheight=15pt, headsep=18pt, footskip=26pt
]{geometry}
\usepackage{pdflscape}
\usepackage{setspace}
\onehalfspacing

% --- Cores canônicas do Brandbook Unooria ("BLUNOORIA") -----------------------
% Regra semântica: "A cor não decora. Ela organiza."
%   - Azul Profundo (#0E1039): Estrutura, corpo, regras e títulos
%   - Lavanda (#8B85AC): Inteligência, decisões arquiteturais (ADR), IA
%   - Azul Claro (#7FA5E2): Interface, navegação, referências e notas
\usepackage{xcolor}
\definecolor{unooNavy}        {HTML}{0E1039}  % Azul Profundo — Estrutura (C100 M96 Y44 K56)
\definecolor{unooLavender}    {HTML}{8B85AC}  % Lavanda — Inteligência (C52 M48 Y15 K2)
\definecolor{unooBlue}        {HTML}{7FA5E2}  % Azul Claro — Interface (C54 M29 Y0 K0)
\definecolor{unooBlueLight}   {HTML}{B8D0F5}  % Variação clara de interface
\definecolor{unooBlueWash}    {HTML}{F0F4FC}  % Fundo suave de interface
\definecolor{unooLavenderWash}{HTML}{F3F2F8}  % Fundo suave de inteligência
\definecolor{unooLavenderSoft}{HTML}{5D577D}  % Lavanda escurecido p/ alto contraste tipográfico
\definecolor{unooInk}         {HTML}{0E1039}  % neutral-950 / texto primário ancorado na marca
\definecolor{unooInkSoft}     {HTML}{3A3C59}  % neutral-700 / texto secundário
\definecolor{unooMuted}       {HTML}{6D6E82}  % neutral-500 / legendas e metadados
\definecolor{unooGrid}        {HTML}{E3E5EE}  % Linhas de grade e bordas sutis
\definecolor{unooAxis}        {HTML}{CACDDC}  % Réguas e separadores
\definecolor{unooSurface}     {HTML}{FFFFFF}  % Fundo puro
\definecolor{unooPaper}       {HTML}{F7F8FC}  % Fundo alternado / zebra
% Tons funcionais (alertas de sistema e status):
\definecolor{unooRed}         {HTML}{C92A2A}  % danger-600
\definecolor{unooRedWash}     {HTML}{FDF2F2}
\definecolor{unooOrange}      {HTML}{D97706}  % warning-600
\definecolor{unooOrangeWash}  {HTML}{FFFBEB}
\definecolor{unooGreen}       {HTML}{2F855A}  % Status funcional aprovado
\definecolor{unooGreenWash}   {HTML}{F0FDF4}
\definecolor{unooYellow}      {HTML}{B45309}
\definecolor{unooYellowWash}  {HTML}{FEF3C7}
\definecolor{unooPink}        {HTML}{9D4EDD}

% --- Tabelas ----------------------------------------------------------------
\usepackage{booktabs}
\usepackage{tabularx}
\usepackage{longtable}
\usepackage{array}
\usepackage{multirow}
\usepackage{makecell}
\usepackage{colortbl}
\usepackage{ltablex}          % tabularx com quebra de pagina
\keepXColumns
\renewcommand{\arraystretch}{1.22}
\arrayrulecolor{unooAxis}
\setlength{\aboverulesep}{0pt}
\setlength{\belowrulesep}{0pt}
\setlength{\extrarowheight}{2pt}

% Tabelas em \small por decisao de composicao: as celulas deste documento
% misturam prosa e identificadores monoespacados longos (caminhos de arquivo,
% nomes de coluna), que NAO hifenizam. Em corpo 11 elas transbordavam a caixa;
% em \small cabem, e a tabela continua legivel porque a densidade e alta de
% proposito. `\emergencystretch` cobre o residuo.
\usepackage{etoolbox}
\AtBeginEnvironment{tabularx}{\small\setlength{\emergencystretch}{1.2em}}
\AtBeginEnvironment{longtable}{\small\setlength{\emergencystretch}{1.2em}}

\newcolumntype{L}[1]{>{\raggedright\arraybackslash}p{#1}}
\newcolumntype{C}[1]{>{\centering\arraybackslash}p{#1}}
\newcolumntype{R}[1]{>{\raggedleft\arraybackslash}p{#1}}
\newcolumntype{Y}{>{\raggedright\arraybackslash}X}

% Colunas X sem justificacao: as celulas misturam prosa e codigo monoespacado,
% que nao hifeniza — justificar produziria espacos internos muito largos.
\renewcommand{\tabularxcolumn}[1]{>{\raggedright\arraybackslash}p{#1}}

\newcommand{\thd}[1]{\textcolor{white}{\textbf{#1}}}
\newcommand{\headrow}{\rowcolor{unooNavy}}
\newcommand{\zebra}{\rowcolor{unooPaper}}

% --- Listas -----------------------------------------------------------------
\usepackage{enumitem}
\setlist[itemize]{leftmargin=1.35em, itemsep=2pt, topsep=4pt, parsep=0pt,
                  label=\textcolor{unooBlue}{\raisebox{0.18ex}{\rule{0.40em}{0.40em}}}}
\setlist[enumerate]{leftmargin=1.7em, itemsep=2pt, topsep=4pt, parsep=0pt,
                    font=\bfseries\color{unooNavy}}
\setlist[description]{leftmargin=1.4em, style=nextline,
                      font=\bfseries\color{unooNavy}}

% --- Matematica e unidades --------------------------------------------------
\usepackage{amsmath}
\usepackage{amssymb}
\usepackage{siunitx}
\sisetup{output-decimal-marker={,}, group-separator={.}, group-minimum-digits=4,
         per-mode=symbol, detect-weight=true, detect-family=true}
\DeclareSIUnit{\tokens}{tokens}
\DeclareSIUnit{\req}{req}

% --- Figuras e legendas -----------------------------------------------------
\usepackage{graphicx}
\usepackage{float}
\usepackage{adjustbox}
\usepackage[font=small, labelfont={bf,color=unooNavy}, textfont=normalfont,
            labelsep=endash, justification=raggedright, singlelinecheck=false,
            skip=8pt]{caption}
\usepackage{subcaption}
\captionsetup[table]{position=top, skip=6pt}

% Ajusta a figura a largura do texto apenas quando ela excede a caixa.
\newcommand{\fitwidth}[1]{\adjustbox{max width=\linewidth, center}{#1}}
\newcommand{\fitpage}[1]{\adjustbox{max width=\linewidth, max height=0.82\textheight, center}{#1}}

% --- Blocos de destaque (tcolorbox) ----------------------------------------
\usepackage[most]{tcolorbox}
\tcbuselibrary{skins,breakable}

\tcbset{
  unoobase/.style={
    enhanced, breakable, sharp corners,
    boxrule=0pt, leftrule=2.6pt, arc=0pt,
    left=10pt, right=10pt, top=7pt, bottom=7pt,
    fonttitle=\bfseries\small,
    before skip=10pt, after skip=10pt
  }
}

% `colbacktitle` e OBRIGATORIO aqui: sem ele o tcolorbox pinta a faixa do
% titulo com a cor de `colframe`, e o titulo (que usa a MESMA cor em
% `coltitle`) fica invisivel — texto vermelho sobre faixa vermelha. O sintoma
% e uma barra colorida solida no lugar do rotulo do bloco.
\newtcolorbox{nota}[1][]{unoobase, colback=unooBlueWash, colframe=unooBlue,
  colbacktitle=unooBlueWash, coltitle=unooNavy, titlerule=0pt, title={NOTA}, #1}
\newtcolorbox{atencao}[1][]{unoobase, colback=unooOrangeWash, colframe=unooOrange,
  colbacktitle=unooOrangeWash, coltitle=unooOrange, titlerule=0pt,
  title={ATENÇÃO}, #1}
\newtcolorbox{decisao}[1][]{unoobase, colback=unooLavenderWash, colframe=unooLavender,
  colbacktitle=unooLavenderWash, coltitle=unooNavy, titlerule=0pt,
  title={DECISÃO DE PROJETO}, #1}
\newtcolorbox{limitacao}[1][]{unoobase, colback=unooRedWash, colframe=unooRed,
  colbacktitle=unooRedWash, coltitle=unooRed, titlerule=0pt,
  title={LIMITAÇÃO CONHECIDA}, #1}
\newtcolorbox{regra}[1][]{unoobase, colback=unooRedWash, colframe=unooRed,
  colbacktitle=unooRedWash, coltitle=unooRed, titlerule=0pt,
  title={REGRA INVIOLÁVEL}, #1}
\newtcolorbox{destaque}[1][]{unoobase, colback=unooPaper, colframe=unooLavender, #1}
\newtcolorbox{especificacao}[1]{unoobase, colback=unooSurface, colframe=unooNavy,
  colbacktitle=unooSurface, coltitle=unooNavy, titlerule=0pt, title={#1}}

% --- Codigo-fonte -----------------------------------------------------------
\usepackage{listings}

\definecolor{lstBg}     {HTML}{FAFBFB}
\definecolor{lstFrame}  {HTML}{E4E9EB}
\definecolor{lstKeyword}{HTML}{2F4F7F}
\definecolor{lstString} {HTML}{0F6B52}
\definecolor{lstComment}{HTML}{617379}
\definecolor{lstEmph}   {HTML}{A9490F}

\lstdefinestyle{unoo}{
  backgroundcolor=\color{lstBg},
  basicstyle=\ttfamily\scriptsize\color{unooInk},
  keywordstyle=\color{lstKeyword}\bfseries,
  keywordstyle=[2]\color{lstEmph}\bfseries,
  keywordstyle=[3]\color{unooBlue},
  stringstyle=\color{lstString},
  commentstyle=\color{lstComment}\itshape,
  numberstyle=\ttfamily\tiny\color{unooMuted},
  identifierstyle=\color{unooInk},
  numbers=left, numbersep=9pt, stepnumber=1,
  frame=single, framesep=6pt, rulecolor=\color{lstFrame},
  xleftmargin=17pt, xrightmargin=2pt,
  breaklines=true, breakatwhitespace=false, breakindent=0pt,
  postbreak=\mbox{\textcolor{unooBlue}{$\hookrightarrow$}\space},
  showstringspaces=false, tabsize=2,
  columns=fullflexible, keepspaces=true, upquote=true,
  aboveskip=10pt, belowskip=10pt,
  literate=
    {á}{{\'a}}1 {à}{{\`a}}1 {ã}{{\~a}}1 {â}{{\^a}}1
    {é}{{\'e}}1 {ê}{{\^e}}1 {í}{{\'i}}1
    {ó}{{\'o}}1 {õ}{{\~o}}1 {ô}{{\^o}}1
    {ú}{{\'u}}1 {ü}{{\"u}}1 {ç}{{\c c}}1
    {Á}{{\'A}}1 {À}{{\`A}}1 {Ã}{{\~A}}1 {Â}{{\^A}}1
    {É}{{\'E}}1 {Ê}{{\^E}}1 {Í}{{\'I}}1
    {Ó}{{\'O}}1 {Õ}{{\~O}}1 {Ô}{{\^O}}1
    {Ú}{{\'U}}1 {Ç}{{\c C}}1
    {°}{{\textdegree}}1 {º}{{\textordmasculine}}1
    {µ}{{$\mu$}}1 {×}{{$\times$}}1 {≈}{{$\approx$}}1
    {≥}{{$\geq$}}1 {≤}{{$\leq$}}1 {≠}{{$\neq$}}1
    {→}{{$\rightarrow$}}1 {←}{{$\leftarrow$}}1 {⇒}{{$\Rightarrow$}}1
    {—}{{---}}3 {–}{{--}}2 {…}{{\ldots}}3
    {“}{{\textquotedblleft}}1 {”}{{\textquotedblright}}1
    {‘}{{\textquoteleft}}1 {’}{{\textquoteright}}1
    {ª}{{\textordfeminine}}1 {§}{{\S}}1
}
\lstset{style=unoo}

\lstdefinelanguage{tsunoo}[]{Java}{
  morekeywords={const,let,export,default,import,from,as,interface,type,await,
                async,function,return,string,number,boolean,readonly,unknown,
                never,satisfies,enum,declare,namespace,implements},
  morekeywords=[2]{z,Prisma,useEffect,postMessage},
  sensitive=true,
  morecomment=[l]{//},
  morecomment=[s]{/*}{*/},
  morestring=[b]",
  morestring=[b]',
  morestring=[b]`
}
\lstdefinelanguage{sqlunoo}[]{SQL}{
  morekeywords={POLICY,ENABLE,FORCE,ROW,LEVEL,SECURITY,USING,PARTITION,RANGE,
                CONFLICT,DO,NOTHING,JSONB,TIMESTAMPTZ,current_setting,
                set_config,ATTACH,DEFAULT,GENERATED,ALWAYS},
  sensitive=false,
  morecomment=[l]{--},
  morestring=[b]'
}
\lstdefinelanguage{yamlunoo}{
  morekeywords={true,false,null,services,image,build,command,environment,
                volumes,ports,depends_on,healthcheck,condition,runs-on,steps,
                uses,needs,jobs,name,on,with,env,if,run,concurrency,
                permissions,workflow_call,workflow_dispatch},
  sensitive=true,
  morecomment=[l]{\#},
  morestring=[b]",
  morestring=[b]'
}
\lstdefinelanguage{jsonunoo}{
  morekeywords={true,false,null},
  keywordstyle=\color{lstEmph}\bfseries,
  sensitive=true,
  morestring=[b]"
}
\lstdefinelanguage{mermaid}{
  morekeywords={flowchart,graph,sequenceDiagram,erDiagram,classDiagram,gantt,
            subgraph,end,participant,actor,alt,else,loop,par,opt,note,activate,
            deactivate,direction,TB,LR,RL,BT,Person,Person_Ext,System,System_Ext,
            SystemDb,Container,ContainerDb,ContainerQueue,Component,ComponentDb,
            Rel,Rel_D,Rel_U,Rel_L,Rel_R,System_Boundary,Container_Boundary,
            Enterprise_Boundary,Boundary,Deployment_Node,UpdateLayoutConfig,
            UpdateRelStyle,classDef,class,style,state,C4Context,C4Container,
            C4Component,C4Deployment,stateDiagram,stateDiagram-v2},
  sensitive=true,
  morecomment=[l]{\%\%},
  morestring=[b]"
}
\lstdefinelanguage{bashunoo}[]{bash}{
  morekeywords={pnpm,docker,gcloud,node,psql,tectonic,turbo,prisma,gsutil},
  sensitive=true
}

% Codigo em linha.
%
% Uso obrigatorio de \texttt e \textcolor (e nao de {\ttfamily\color{...}}):
% os comandos \text... emitem \leavevmode antes de trocar a fonte, o que faz o
% paragrafo comecar na fonte do texto corrente. Sem isso, uma celula de tabela
% que inicia com codigo teria a primeira linha de base calculada com as metricas
% da fonte monoespacada e ficaria visivelmente desalinhada das demais colunas.
\newcommand{\cd}[1]{\texttt{\textcolor{unooNavy}{#1}}}
\newcommand{\arq}[1]{\texttt{\textcolor{unooInkSoft}{#1}}}
\newcommand{\rota}[1]{\texttt{\textcolor{unooGreen}{#1}}}
\newcommand{\env}[1]{\texttt{\textcolor{lstEmph}{#1}}}
% Ponto de quebra de linha sem hifen, para caminho longo em celula estreita:
% identificador monoespacado nao hifeniza, e sem isto ele transborda a caixa.
\newcommand{\B}{\discretionary{}{}{}}

% --- Titulos ----------------------------------------------------------------
\usepackage{titlesec}
\usepackage{titletoc}

\titleformat{\chapter}[display]
  {\normalfont}
  {\raggedleft\fontsize{62}{62}\selectfont\color{unooBlueLight}\textbf{\thechapter}}
  {-10pt}
  {\raggedright\fontsize{23}{27}\selectfont\bfseries\color{unooNavy}}
  [\vspace{5pt}{\color{unooBlue}\rule{\linewidth}{2.2pt}}]
\titlespacing*{\chapter}{0pt}{-30pt}{26pt}

\titleformat{name=\chapter,numberless}[display]
  {\normalfont}
  {}
  {0pt}
  {\raggedright\fontsize{23}{27}\selectfont\bfseries\color{unooNavy}}
  [\vspace{5pt}{\color{unooBlue}\rule{\linewidth}{2.2pt}}]

\titleformat{\section}
  {\normalfont\Large\bfseries\color{unooNavy}}{\thesection}{0.7em}{}
\titleformat{\subsection}
  {\normalfont\large\bfseries\color{unooInk}}{\thesubsection}{0.7em}{}
\titleformat{\subsubsection}
  {\normalfont\normalsize\bfseries\color{unooInkSoft}}{\thesubsubsection}{0.7em}{}
\titleformat{\paragraph}[runin]
  {\normalfont\normalsize\bfseries\color{unooNavy}}{}{0pt}{}[\;---\;]

\titlespacing*{\section}{0pt}{20pt}{8pt}
\titlespacing*{\subsection}{0pt}{15pt}{6pt}
\titlespacing*{\subsubsection}{0pt}{12pt}{5pt}
\titlespacing*{\paragraph}{0pt}{9pt}{0.4em}

\setcounter{secnumdepth}{3}
\setcounter{tocdepth}{2}

% --- Sumario ----------------------------------------------------------------
\titlecontents{chapter}[0em]
  {\addvspace{9pt}\bfseries\color{unooNavy}}
  {\contentslabel[\thecontentslabel]{1.5em}}
  {}
  {\hfill\contentspage}
\titlecontents{section}[1.7em]
  {\addvspace{1.5pt}}
  {\contentslabel{2.2em}}
  {}
  {\titlerule*[0.7pc]{\color{unooAxis}.}\contentspage}
\titlecontents{subsection}[4.0em]
  {\addvspace{0.5pt}\small\color{unooInkSoft}}
  {\contentslabel{2.8em}}
  {}
  {\titlerule*[0.7pc]{\color{unooGrid}.}\contentspage}

% --- Cabecalho e rodape -----------------------------------------------------
\usepackage{fancyhdr}
\usepackage{lastpage}

\newcommand{\docid}{PU-DOC-001}
\newcommand{\docversao}{1.2}

\renewcommand{\headrule}{{\color{unooGrid}\hrule height \headrulewidth}}
\renewcommand{\footrule}{{\color{unooGrid}\hrule height \footrulewidth}}

\fancypagestyle{unoomain}{%
  \fancyhf{}
  \fancyhead[L]{\footnotesize\color{unooInkSoft}\nouppercase{\leftmark}}
  \fancyhead[R]{\raisebox{-0.15ex}{\includegraphics[height=1.8ex]{figuras/marca/simbolo_azul_profundo.png}}\hspace{4pt}\footnotesize\bfseries\color{unooNavy}Unooria}
  \fancyfoot[L]{\footnotesize\color{unooMuted}\docid{} \textperiodcentered{} v\docversao}
  \fancyfoot[R]{\footnotesize\color{unooInkSoft}\thepage\,/\,\pageref{LastPage}}
  \renewcommand{\headrulewidth}{0.5pt}
  \renewcommand{\footrulewidth}{0.5pt}
}
% Pre-textual: numeracao romana, sem total de paginas (o total do
% \pageref{LastPage} refere-se a numeracao arabe do corpo do documento).
\fancypagestyle{unoofront}{%
  \fancyhf{}
  \fancyhead[L]{\footnotesize\color{unooInkSoft}\nouppercase{\leftmark}}
  \fancyhead[R]{\raisebox{-0.15ex}{\includegraphics[height=1.8ex]{figuras/marca/simbolo_azul_profundo.png}}\hspace{4pt}\footnotesize\bfseries\color{unooNavy}Unooria}
  \fancyfoot[L]{\footnotesize\color{unooMuted}\docid{} \textperiodcentered{} v\docversao}
  \fancyfoot[R]{\footnotesize\color{unooInkSoft}\thepage}
  \renewcommand{\headrulewidth}{0.5pt}
  \renewcommand{\footrulewidth}{0.5pt}
}
\fancypagestyle{plain}{%
  \fancyhf{}
  \fancyfoot[R]{\footnotesize\color{unooInkSoft}\thepage}
  \renewcommand{\headrulewidth}{0pt}
  \renewcommand{\footrulewidth}{0pt}
}
\pagestyle{unoofront}
\renewcommand{\chaptermark}[1]{\markboth{\thechapter\quad #1}{}}

% --- Referencias cruzadas e PDF --------------------------------------------
\usepackage[
  colorlinks=true,
  linkcolor=unooNavy,
  citecolor=unooBlue,
  urlcolor=unooBlue,
  filecolor=unooBlue,
  bookmarksnumbered=true,
  bookmarksopen=true,
  bookmarksopenlevel=1,
  pdfdisplaydoctitle=true
]{hyperref}
\usepackage{bookmark}
\usepackage{url}
\urlstyle{sf}

% --- Bibliografia (ABNT NBR 6023) ------------------------------------------
%
% `abntex2cite` + BibTeX, e nao biblatex+biber: o pipeline de compilacao
% (Tectonic) chama o BibTeX classico automaticamente, sem ferramenta externa.
% A saida sai formatada em ABNT do mesmo jeito.
\usepackage[alf]{abntex2cite}

% --- Semantica de referencia do projeto ------------------------------------
\newcommand{\rf}[1]{\textbf{\textcolor{unooNavy}{#1}}}
\newcommand{\adrr}[1]{\textbf{\textcolor{unooLavenderSoft}{#1}}}
\newcommand{\ucr}[1]{\textbf{\textcolor{unooBlue}{#1}}}
\newcommand{\drr}[1]{\textbf{\textcolor{unooOrange}{#1}}}
\newcommand{\yes}{\textcolor{unooGreen}{\textbf{\checkmark}}}
\newcommand{\no}{\textcolor{unooOrange}{\textbf{--}}}

% Setas: a Roboto NAO tem U+2192 nem U+21D2 (medido — o XeTeX avisa
% "Missing character" e o glifo simplesmente nao e desenhado). Em prosa,
% sempre estes comandos; nunca a seta literal.
\newcommand{\seta}{\ensuremath{\rightarrow}}
\newcommand{\implica}{\ensuremath{\Rightarrow}}
\newcommand{\setabi}{\ensuremath{\leftrightarrow}}

% Espacamento de paragrafo
\setlength{\parindent}{0pt}
\setlength{\parskip}{7pt}

\widowpenalty=10000
\clubpenalty=10000
\flushbottom

% ============================================================================
%  diagramas.tex — biblioteca TikZ do documento
%
%  Todos os diagramas do documento (C4, UML de casos de uso, atividades,
%  sequencia, entidade-relacionamento, classes e infraestrutura) sao vetoriais
%  e desenhados nativamente em TikZ. Isso garante tipografia unica (Roboto),
%  nitidez em qualquer zoom e ausencia de dependencia externa de renderizacao.
%  O codigo Mermaid equivalente de cada diagrama esta no Apendice A.
% ============================================================================

\usepackage{tikz}
\usetikzlibrary{
  arrows.meta, positioning, fit, backgrounds, calc, shapes.geometric,
  shapes.symbols, shapes.multipart, shapes.misc, decorations.pathreplacing,
  decorations.pathmorphing, decorations.markings, patterns, patterns.meta,
  matrix, chains, scopes, quotes, shadows.blur, math, intersections
}

\usepackage{pgfplots}
\pgfplotsset{compat=1.18}
\usepgfplotslibrary{groupplots}

% Numeracao dos graficos na convencao brasileira (virgula decimal), coerente
% com o restante do documento.
\pgfkeys{/pgf/number format/.cd, use comma, 1000 sep={.}}

% Estilo comum dos graficos quantitativos: eixos recessivos, grade sutil,
% tipografia do documento. Mesma paleta validada da interface.
\pgfplotsset{
  unooplot/.style={
    width=\linewidth, height=6cm,
    axis line style={draw=unooAxis, line width=0.6pt},
    tick label style={font=\scriptsize\color{unooInkSoft}},
    label style={font=\small\color{unooInkSoft}},
    title style={font=\small\bfseries\color{unooNavy}, align=left},
    legend style={font=\scriptsize, draw=unooGrid, fill=unooSurface,
                  /tikz/every even column/.append style={column sep=6pt}},
    grid=major, grid style={draw=unooGrid, line width=0.5pt},
    major tick style={draw=unooAxis},
    axis background/.style={fill=unooSurface},
    every axis plot post/.append style={line width=0.9pt},
    enlarge x limits=0.06,
  }
}

% ---------------------------------------------------------------------------
%  Base comum
% ---------------------------------------------------------------------------
\tikzset{
  >={Stealth[length=5.4pt, width=4.2pt]},
  every node/.append style={align=center, inner sep=4pt},
  cxlabel/.style={font=\scriptsize\color{unooInkSoft}, fill=unooSurface,
                  inner sep=1.6pt, rounded corners=1pt},
  cxtitle/.style={font=\small\bfseries\color{unooNavy}},
  legenda/.style={font=\scriptsize\color{unooMuted}},
}

% ---------------------------------------------------------------------------
%  C4 model — notacao de Simon Brown adaptada a paleta do projeto
% ---------------------------------------------------------------------------
\tikzset{
  % --- elementos -----------------------------------------------------------
  c4box/.style={
    rectangle, rounded corners=2.5pt, draw, line width=0.8pt,
    text width=32mm, minimum height=15mm, font=\footnotesize,
    inner sep=5pt
  },
  c4person/.style={
    c4box, fill=unooNavy, draw=unooNavy!85!black, text=white,
    rounded corners=7pt, text width=28mm, minimum height=16mm
  },
  c4system/.style={
    c4box, fill=unooBlue, draw=unooNavy, text=white
  },
  c4systemext/.style={
    c4box, fill=unooMuted, draw=unooInkSoft, text=white
  },
  c4container/.style={
    c4box, fill=unooBlue!88, draw=unooNavy, text=white
  },
  c4containerext/.style={
    c4box, fill=unooMuted!85, draw=unooInkSoft, text=white
  },
  c4component/.style={
    c4box, fill=unooBlueLight!45, draw=unooBlue, text=unooInk, text width=30mm,
    minimum height=13mm
  },
  c4db/.style={
    cylinder, shape border rotate=90, aspect=0.28, draw, line width=0.8pt,
    fill=unooNavy, draw=unooNavy!80!black, text=white,
    text width=25mm, minimum height=17mm, font=\footnotesize, inner sep=3pt
  },
  c4queue/.style={
    c4box, fill=unooYellow!85, draw=unooYellow!60!black, text=unooInk,
    rounded corners=8pt
  },
  c4artifact/.style={
    c4box, fill=unooGreen!80, draw=unooGreen!60!black, text=white
  },
  % --- fronteiras ----------------------------------------------------------
  c4boundary/.style={
    draw=unooAxis, dashed, dash pattern=on 3pt off 2.4pt, line width=0.9pt,
    rounded corners=4pt, inner sep=9pt
  },
  c4boundarylbl/.style={
    font=\scriptsize\bfseries\color{unooInkSoft}, anchor=north west,
    inner sep=2pt
  },
  c4zone/.style={
    c4boundary, draw=unooBlue!55, fill=unooBlue!4
  },
  c4zoneext/.style={
    c4boundary, draw=unooMuted, fill=unooPaper
  },
  % --- relacoes ------------------------------------------------------------
  c4rel/.style={-{Stealth[length=5pt,width=4pt]}, line width=0.85pt,
                draw=unooInkSoft},
  c4relbi/.style={c4rel, {Stealth[length=5pt,width=4pt]}-{Stealth[length=5pt,width=4pt]}},
  c4reldash/.style={c4rel, dashed, dash pattern=on 3pt off 2pt},
  c4relhi/.style={c4rel, draw=unooOrange, line width=1.1pt},
}

% Rotulo tecnologico dentro de um elemento C4
\newcommand{\tech}[1]{\\[1.5pt]{\scriptsize\itshape [#1]}}
% Descricao curta dentro de um elemento C4
\newcommand{\desc}[1]{\\[2pt]{\scriptsize #1}}

% ---------------------------------------------------------------------------
%  UML — casos de uso
% ---------------------------------------------------------------------------
\tikzset{
  ucase/.style={
    ellipse, draw=unooBlue, line width=0.85pt, fill=unooBlueWash,
    text=unooInk, font=\scriptsize, minimum width=34mm, minimum height=11mm,
    inner sep=2.5pt, text width=31mm
  },
  ucasesec/.style={ucase, fill=unooSurface, draw=unooAxis},
  ucsystem/.style={
    draw=unooNavy, line width=1pt, rounded corners=3pt, fill=unooBlue!3,
    inner sep=11pt
  },
  ucrel/.style={line width=0.8pt, draw=unooInkSoft},
  ucinclude/.style={
    -{Stealth[length=5pt,width=4pt]}, dashed, dash pattern=on 3pt off 2pt,
    line width=0.75pt, draw=unooGreen
  },
  ucextend/.style={
    -{Stealth[length=5pt,width=4pt]}, dashed, dash pattern=on 3pt off 2pt,
    line width=0.75pt, draw=unooOrange
  },
  ucgen/.style={-{Triangle[open, length=6pt, width=6pt]}, line width=0.8pt,
                draw=unooInkSoft},
  ucnote/.style={font=\scriptsize\color{unooInkSoft}},
}

% Ator UML (boneco de palito) — #1 nome do no, #2 posicao, #3 rotulo
\newcommand{\ator}[3]{%
  \begin{scope}[shift={#2}]
    \draw[line width=0.9pt, draw=unooNavy] (0,0.30) circle (0.155);
    \draw[line width=0.9pt, draw=unooNavy] (0,0.145) -- (0,-0.24);
    \draw[line width=0.9pt, draw=unooNavy] (-0.245,-0.01) -- (0.245,-0.01);
    \draw[line width=0.9pt, draw=unooNavy] (0,-0.24) -- (-0.20,-0.60);
    \draw[line width=0.9pt, draw=unooNavy] (0,-0.24) -- (0.20,-0.60);
    \node[font=\scriptsize\bfseries\color{unooNavy}, text width=24mm,
          anchor=north] at (0,-0.68) {#3};
    \node[inner sep=0pt, minimum width=13mm, minimum height=19mm] (#1) at (0,-0.15) {};
  \end{scope}
}

% Ator secundario (sistema externo) — retangulo com estereotipo
\newcommand{\atorsys}[3]{%
  \node[draw=unooInkSoft, line width=0.85pt, fill=unooPaper, rounded corners=2pt,
        font=\scriptsize, text width=22mm, minimum height=12mm] (#1) at #2
        {{\scriptsize$\ll$ator$\gg$}\\\textbf{#3}};
}

% ---------------------------------------------------------------------------
%  UML — diagramas de atividades
% ---------------------------------------------------------------------------
\tikzset{
  actstart/.style={
    circle, fill=unooInk, draw=unooInk, minimum size=4.6mm, inner sep=0pt
  },
  actend/.style={
    circle, fill=unooInk, draw=unooInk, minimum size=4.6mm, inner sep=0pt,
    double, double distance=1.1pt, outer sep=0pt
  },
  action/.style={
    rectangle, rounded corners=4pt, draw=unooBlue, line width=0.85pt,
    fill=unooBlueWash, text=unooInk, font=\scriptsize, text width=36mm,
    minimum height=10mm, inner sep=4pt
  },
  actionsys/.style={action, fill=unooSurface, draw=unooAxis},
  actionerr/.style={action, fill=unooOrangeWash, draw=unooOrange},
  actionok/.style={action, fill=unooGreenWash, draw=unooGreen},
  decision/.style={
    diamond, draw=unooNavy, line width=0.85pt, fill=unooYellowWash,
    aspect=1.9, font=\scriptsize, text width=26mm, inner sep=1pt
  },
  bar/.style={
    rectangle, fill=unooInk, minimum width=34mm, minimum height=1.5mm,
    inner sep=0pt
  },
  barv/.style={
    rectangle, fill=unooInk, minimum width=1.5mm, minimum height=30mm,
    inner sep=0pt
  },
  flow/.style={-{Stealth[length=5pt,width=4pt]}, line width=0.85pt,
               draw=unooInkSoft},
  flowlbl/.style={font=\scriptsize\color{unooInkSoft}, fill=unooSurface,
                  inner sep=1.5pt},
  lane/.style={
    draw=unooAxis, line width=0.7pt, fill=unooSurface
  },
  lanehdr/.style={
    draw=unooAxis, line width=0.7pt, fill=unooNavy, text=white,
    font=\scriptsize\bfseries
  },
  actnote/.style={
    draw=unooYellow!70!black, fill=unooYellowWash, line width=0.6pt,
    font=\scriptsize\color{unooInkSoft}, text width=34mm, align=left,
    rounded corners=0pt, inner sep=4pt
  },
  actnoteline/.style={dashed, dash pattern=on 2pt off 1.6pt, draw=unooYellow!70!black},
  objflow/.style={
    rectangle, draw=unooGreen, line width=0.7pt, fill=unooGreenWash,
    font=\scriptsize, text width=30mm, inner sep=3.5pt
  },
}

% ---------------------------------------------------------------------------
%  UML — diagramas de sequencia
% ---------------------------------------------------------------------------
\tikzset{
  seqhead/.style={
    rectangle, rounded corners=2pt, draw=unooNavy, line width=0.8pt,
    fill=unooNavy, text=white, font=\scriptsize, text width=26mm,
    minimum height=9mm, inner sep=3pt
  },
  seqheadext/.style={seqhead, fill=unooMuted, draw=unooInkSoft},
  seqheadact/.style={seqhead, fill=unooBlue, draw=unooNavy},
  seqline/.style={dashed, dash pattern=on 2.6pt off 2.2pt, draw=unooAxis,
                  line width=0.7pt},
  seqmsg/.style={-{Stealth[length=5pt,width=4pt]}, line width=0.85pt,
                 draw=unooInkSoft},
  seqret/.style={seqmsg, dashed, dash pattern=on 2.6pt off 2pt,
                 draw=unooInkSoft},
  seqself/.style={seqmsg, rounded corners=2pt},
  seqlbl/.style={font=\scriptsize\color{unooInk}, fill=unooSurface,
                 inner sep=1.6pt},
  seqact/.style={fill=unooBlueLight!60, draw=unooBlue, line width=0.5pt,
                 minimum width=2.2mm, inner sep=0pt},
  seqframe/.style={draw=unooAxis, line width=0.7pt, rounded corners=0pt},
  seqframelbl/.style={fill=unooNavy, text=white, font=\scriptsize\bfseries,
                      inner sep=2.2pt, anchor=north west},
  seqcond/.style={font=\scriptsize\itshape\color{unooInkSoft}, anchor=west},
}

% Raia de sequencia: #1 = x, #2 = profundidade (y final, negativo),
% #3 = titulo, #4 = tecnologia
\newcommand{\raia}[4]{%
  \node[seqhead] at (#1,0) {\textbf{#3}\\[0.5pt]{\scriptsize\itshape #4}};
  \draw[seqline] (#1,-0.52) -- (#1,#2);
}
\newcommand{\raiaext}[4]{%
  \node[seqheadext] at (#1,0) {\textbf{#3}\\[0.5pt]{\scriptsize\itshape #4}};
  \draw[seqline] (#1,-0.52) -- (#1,#2);
}
% Mensagem: #1 = x origem, #2 = x destino, #3 = y, #4 = rotulo
\newcommand{\msgs}[4]{%
  \draw[seqmsg] (#1,#3) -- (#2,#3) node[seqlbl, midway, above=0.4pt]{#4};
}
% Retorno (tracejado)
\newcommand{\msgr}[4]{%
  \draw[seqret] (#1,#3) -- (#2,#3) node[seqlbl, midway, above=0.4pt]{#4};
}
% Auto-mensagem: #1 = x, #2 = y, #3 = rotulo
\newcommand{\msgself}[3]{%
  \draw[seqself] (#1,#2) -- ++(0.62,0) -- ++(0,-0.42) -- ++(-0.62,0);
  \node[seqlbl, anchor=west] at (#1+0.72,#2-0.21) {#3};
}

% ---------------------------------------------------------------------------
%  Entidade-relacionamento e classes
% ---------------------------------------------------------------------------
\tikzset{
  entidade/.style={
    rectangle split, rectangle split parts=2, draw=unooNavy,
    line width=0.85pt, rectangle split part fill={unooNavy, unooSurface},
    font=\scriptsize, text width=37mm, inner sep=4pt, rounded corners=1pt,
    every text node part/.style={align=center},
    every two node part/.style={align=left}
  },
  classe/.style={
    rectangle split, rectangle split parts=3, draw=unooNavy,
    line width=0.85pt, rectangle split part fill={unooNavy, unooSurface, unooPaper},
    font=\scriptsize, text width=44mm, inner sep=4pt, rounded corners=1pt,
    every text node part/.style={align=center},
    every two node part/.style={align=left},
    every three node part/.style={align=left}
  },
  erel/.style={line width=0.85pt, draw=unooInkSoft},
  ercard/.style={font=\scriptsize\color{unooNavy}, fill=unooSurface,
                 inner sep=1.4pt},
  hered/.style={-{Triangle[open, length=6.5pt, width=6.5pt]}, line width=0.85pt,
                draw=unooInkSoft},
  compos/.style={-{Diamond[open, length=6pt, width=4.5pt]}, line width=0.85pt,
                 draw=unooInkSoft},
  depend/.style={-{Stealth[length=5pt,width=4pt]}, dashed,
                 dash pattern=on 3pt off 2pt, line width=0.75pt,
                 draw=unooInkSoft},
}

% Titulo branco em negrito para a primeira particao de entidade/classe
\newcommand{\etitulo}[1]{\textcolor{white}{\textbf{#1}}}
\newcommand{\pk}[1]{\textcolor{unooNavy}{\textbf{#1}}}
\newcommand{\fk}[1]{\textcolor{unooOrange}{#1}}

% ---------------------------------------------------------------------------
%  Infraestrutura / implantacao
% ---------------------------------------------------------------------------
\tikzset{
  nuvem/.style={
    draw=unooBlue, line width=1pt, dash pattern=on 4pt off 2.6pt,
    rounded corners=6pt, fill=unooBlue!3, inner sep=11pt
  },
  nuvemlbl/.style={
    font=\scriptsize\bfseries\color{unooNavy}, anchor=north west,
    fill=unooSurface, inner sep=2.4pt
  },
  regiao/.style={
    draw=unooAxis, line width=0.8pt, dash pattern=on 2.6pt off 2pt,
    rounded corners=4pt, fill=unooSurface, inner sep=8pt
  },
  servico/.style={
    rectangle, rounded corners=2.5pt, draw=unooNavy, line width=0.85pt,
    fill=unooBlue!88, text=white, font=\scriptsize, text width=30mm,
    minimum height=13mm, inner sep=4pt
  },
  servicoger/.style={servico, fill=unooNavy},
  servicobroker/.style={servico, fill=unooYellow!85, text=unooInk,
                        draw=unooYellow!55!black},
  volume/.style={
    cylinder, shape border rotate=90, aspect=0.30, draw=unooGreen!60!black,
    line width=0.85pt, fill=unooGreen!80, text=white, font=\scriptsize,
    text width=22mm, minimum height=15mm, inner sep=3pt
  },
  externo/.style={
    rectangle, rounded corners=2.5pt, draw=unooInkSoft, line width=0.85pt,
    fill=unooMuted, text=white, font=\scriptsize, text width=30mm,
    minimum height=13mm, inner sep=4pt
  },
  internet/.style={
    cloud, cloud puffs=13, cloud puff arc=110, draw=unooAxis, line width=0.85pt,
    fill=unooPaper, text=unooInkSoft, font=\scriptsize, minimum width=26mm,
    minimum height=15mm, aspect=1.6
  },
  netpriv/.style={-{Stealth[length=5pt,width=4pt]}, line width=0.9pt,
                  draw=unooGreen!70!black},
  netpub/.style={-{Stealth[length=5pt,width=4pt]}, line width=0.9pt,
                 draw=unooOrange},
  netbi/.style={{Stealth[length=5pt,width=4pt]}-{Stealth[length=5pt,width=4pt]},
                line width=0.9pt, draw=unooInkSoft},
  portlbl/.style={font=\scriptsize\ttfamily\color{unooInkSoft},
                  fill=unooSurface, inner sep=1.5pt},
}

% ---------------------------------------------------------------------------
%  Pipeline / fluxo de dados
% ---------------------------------------------------------------------------
\tikzset{
  etapa/.style={
    rectangle, rounded corners=2.5pt, draw=unooBlue, line width=0.85pt,
    fill=unooBlueWash, font=\scriptsize, text width=27mm, minimum height=13mm,
    inner sep=4pt
  },
  etapaio/.style={
    etapa, fill=unooPaper, draw=unooAxis
  },
  etapaout/.style={etapa, fill=unooGreenWash, draw=unooGreen},
  etapaerr/.style={etapa, fill=unooOrangeWash, draw=unooOrange},
  fita/.style={
    signal, signal to=east, signal from=west, draw=unooBlue,
    line width=0.85pt, fill=unooBlueWash, font=\scriptsize, text width=24mm,
    minimum height=12mm
  },
  pipe/.style={-{Stealth[length=5.4pt,width=4.4pt]}, line width=1pt,
               draw=unooBlue},
  pipelbl/.style={font=\scriptsize\color{unooInkSoft}, fill=unooSurface,
                  inner sep=1.5pt},
}

% ---------------------------------------------------------------------------
%  Legenda padronizada de diagrama
% ---------------------------------------------------------------------------
\newcommand{\amostra}[1]{%
  \tikz[baseline=-0.55ex]{\node[#1, text width=0pt, minimum width=3.6mm,
    minimum height=3.0mm, inner sep=0pt, font=\tiny]{};}%
}
\newenvironment{legendadiag}
  {\par\vspace{2pt}\begingroup\footnotesize\color{unooInkSoft}%
   \setlength{\parskip}{2pt}\textbf{Legenda:}\space}
  {\endgroup\par\vspace{2pt}}

% ---------------------------------------------------------------------------
%  Acrescimos deste documento (nao vieram da biblioteca de origem)
% ---------------------------------------------------------------------------
\tikzset{
  % Bundle estatico servido pela CDN — nao e um servico, e um artefato.
  bundle/.style={
    rectangle, draw=unooGreen!60!black, line width=0.85pt, fill=unooGreen!80,
    text=white, font=\scriptsize, text width=28mm, minimum height=13mm,
    inner sep=4pt, rounded corners=1pt
  },
  % Fronteira do iframe: a linha que separa o chrome do hub da ficcao da
  % experiencia. Tracejada e grossa de proposito — e fronteira de seguranca.
  iframe/.style={
    draw=unooOrange, line width=1.1pt, dash pattern=on 4pt off 2.4pt,
    rounded corners=3pt, fill=unooOrangeWash!45, inner sep=9pt
  },
  iframelbl/.style={
    font=\scriptsize\bfseries\color{unooOrange}, anchor=north west,
    fill=unooSurface, inner sep=2.2pt
  },
  % Maquina de estados
  estado/.style={
    rectangle, rounded corners=6pt, draw=unooBlue, line width=0.85pt,
    fill=unooBlueWash, font=\scriptsize, text width=24mm, minimum height=9mm,
    inner sep=3pt
  },
  estadoini/.style={estado, fill=unooNavy, text=white, draw=unooNavy},
  estadofim/.style={estado, fill=unooGreenWash, draw=unooGreen},
  estadoerr/.style={estado, fill=unooRedWash, draw=unooRed},
  trans/.style={-{Stealth[length=5pt,width=4pt]}, line width=0.8pt,
                draw=unooInkSoft},
  translbl/.style={font=\scriptsize\color{unooInkSoft}, fill=unooSurface,
                   inner sep=1.4pt},
  % Camada horizontal (diagramas de camadas)
  camada/.style={
    rectangle, draw=unooAxis, line width=0.7pt, fill=unooSurface,
    minimum height=13mm, inner sep=5pt, font=\scriptsize
  },
  camadalbl/.style={font=\scriptsize\bfseries\color{unooNavy}},
  % Tabela do banco, com prefixo por experiencia
  tabela/.style={
    rectangle split, rectangle split parts=2, draw=unooNavy,
    line width=0.8pt, rectangle split part fill={unooNavy, unooSurface},
    font=\scriptsize, text width=34mm, inner sep=3.5pt, rounded corners=1pt,
    every text node part/.style={align=center},
    every two node part/.style={align=left}
  },
  tabelaplat/.style={tabela, rectangle split part fill={unooBlue, unooSurface}},
  tabelacat/.style={tabela, rectangle split part fill={unooMuted, unooPaper}},
}

% Rotulo de prefixo de tabela por experiencia
\newcommand{\pref}[1]{{\scriptsize\ttfamily #1}}

% ---------------------------------------------------------------------------
%  UML 2.5.1 --- os tipos que a biblioteca de origem nao tinha (Secao
%  "Modelagem UML", secoes/05b-modelagem-uml.tex; ADR-122)
%
%  Pacote, componente com interfaces fornecidas e requeridas, no de
%  implantacao, artefato, objeto, estrutura composta (parte e porta), perfil,
%  fragmento de interacao, comunicacao e tempo. A notacao segue a
%  especificacao da OMG (UML 2.5.1, formal/17-12-05); a cor segue os tokens
%  unoo* como o resto do documento.
% ---------------------------------------------------------------------------
\newcommand{\estereo}[1]{{\guillemotleft}#1{\guillemotright}}

\tikzset{
  % --- relacoes de classe que faltavam ------------------------------------
  associa/.style={line width=0.85pt, draw=unooInkSoft},
  navega/.style={associa, -{Stealth[length=5pt,width=4pt]}},
  agrega/.style={associa, {Diamond[open, length=7.5pt, width=5pt]}-},
  compoe/.style={associa, {Diamond[length=7.5pt, width=5pt, fill=unooInkSoft]}-},
  realiza/.style={-{Triangle[open, length=6.5pt, width=6.5pt]}, dashed,
                  dash pattern=on 3pt off 2pt, line width=0.8pt, draw=unooInkSoft},
  % extensao de perfil: seta cheia, que nao se confunde com a heranca (aberta)
  estende/.style={-{Triangle[length=6.5pt, width=6.5pt, fill=unooInkSoft]},
                  line width=0.85pt, draw=unooInkSoft},
  multip/.style={font=\scriptsize\color{unooNavy}, inner sep=1.2pt},
  papelassoc/.style={font=\scriptsize\itshape\color{unooInkSoft}, inner sep=1.2pt},
  % --- classificadores ----------------------------------------------------
  interface/.style={classe, draw=unooLavenderSoft,
    rectangle split part fill={unooLavenderSoft, unooSurface, unooPaper}},
  enumeracao/.style={entidade, draw=unooInkSoft, text width=30mm,
    rectangle split part fill={unooMuted, unooSurface}},
  metaclasse/.style={rectangle, draw=unooInkSoft, line width=0.8pt,
    fill=unooPaper, font=\scriptsize, text width=28mm, minimum height=11mm,
    inner sep=3pt},
  % --- objeto: nome sublinhado, instancia e nao tipo -----------------------
  objeto/.style={rectangle split, rectangle split parts=2, draw=unooBlue,
    line width=0.85pt, rectangle split part fill={unooBlueLight!55, unooSurface},
    font=\scriptsize, text width=40mm, inner sep=3.5pt,
    every text node part/.style={align=center},
    every two node part/.style={align=left}},
  objetoum/.style={rectangle, draw=unooBlue, line width=0.85pt,
    fill=unooBlueLight!45, font=\scriptsize, text width=30mm,
    minimum height=10mm, inner sep=3pt},
  objetolink/.style={line width=0.8pt, draw=unooInkSoft},
  % --- pacote --------------------------------------------------------------
  pacotecorpo/.style={draw=unooNavy, line width=0.85pt, fill=unooBlueWash,
    inner sep=0pt},
  pacoteaba/.style={draw=unooNavy, line width=0.85pt, fill=unooNavy,
    text=white, font=\scriptsize\bfseries, inner sep=2.4pt,
    minimum height=4.6mm},
  pacotedep/.style={-{Stealth[length=5pt,width=4pt]}, dashed,
    dash pattern=on 3pt off 2pt, line width=0.8pt, draw=unooInkSoft},
  % --- componente: retangulo com o icone de componente no canto -------------
  componente/.style={rectangle, draw=unooNavy, line width=0.85pt,
    fill=unooBlueWash, font=\scriptsize, text width=30mm, minimum height=15mm,
    inner sep=4pt,
    path picture={
      \coordinate (icc) at ([xshift=-3.0mm, yshift=-3.1mm]path picture bounding box.north east);
      \draw[draw=unooNavy, line width=0.45pt, fill=unooSurface]
        ([xshift=-1.3mm, yshift=-1.8mm]icc) rectangle ([xshift=1.6mm, yshift=1.8mm]icc);
      \draw[draw=unooNavy, line width=0.45pt, fill=unooSurface]
        ([xshift=-2.1mm, yshift=0.5mm]icc) rectangle ([xshift=-0.5mm, yshift=1.2mm]icc);
      \draw[draw=unooNavy, line width=0.45pt, fill=unooSurface]
        ([xshift=-2.1mm, yshift=-1.2mm]icc) rectangle ([xshift=-0.5mm, yshift=-0.5mm]icc);
    }},
  componenteext/.style={componente, fill=unooPaper, draw=unooInkSoft},
  ifacelbl/.style={font=\scriptsize\color{unooInkSoft}, inner sep=1.4pt},
  % --- implantacao ---------------------------------------------------------
  nodocorpo/.style={draw=unooNavy, line width=0.85pt, fill=unooSurface,
    inner sep=0pt},
  nodolbl/.style={font=\scriptsize\color{unooNavy}, align=left, inner sep=3pt},
  artefato/.style={rectangle, draw=unooGreen!60!black, line width=0.8pt,
    fill=unooGreenWash, font=\scriptsize, text width=30mm, minimum height=11mm,
    inner sep=3.5pt,
    path picture={
      \coordinate (ica) at ([xshift=-2.9mm, yshift=-2.9mm]path picture bounding box.north east);
      \draw[draw=unooGreen!60!black, line width=0.45pt, fill=unooSurface]
        ([xshift=-1.3mm, yshift=1.7mm]ica) -- ([xshift=-1.3mm, yshift=-1.7mm]ica)
        -- ([xshift=1.3mm, yshift=-1.7mm]ica) -- ([xshift=1.3mm, yshift=0.8mm]ica)
        -- ([xshift=0.4mm, yshift=1.7mm]ica) -- cycle;
      \draw[draw=unooGreen!60!black, line width=0.45pt]
        ([xshift=0.4mm, yshift=1.7mm]ica) -- ([xshift=0.4mm, yshift=0.8mm]ica)
        -- ([xshift=1.3mm, yshift=0.8mm]ica);
    }},
  implanta/.style={-{Stealth[length=5pt,width=4pt]}, dashed,
    dash pattern=on 3pt off 2pt, line width=0.75pt, draw=unooGreen!60!black},
  caminhocom/.style={line width=0.9pt, draw=unooInkSoft},
  % --- estrutura composta: parte, porta e conector -------------------------
  parte/.style={rectangle, draw=unooNavy, line width=0.85pt, fill=unooSurface,
    font=\scriptsize, text width=28mm, minimum height=11mm, inner sep=3pt},
  partemult/.style={parte, fill=unooBlueWash},
  porta/.style={rectangle, draw=unooNavy, line width=0.8pt, fill=unooSurface,
    minimum size=2.6mm, inner sep=0pt},
  portaborda/.style={porta, fill=unooOrange!85, draw=unooOrange!70!black},
  conector/.style={line width=0.85pt, draw=unooInkSoft},
  % --- interacao: operador de fragmento e uso de interacao (ref) -----------
  frameop/.style={chamfered rectangle, chamfered rectangle corners=south east,
    chamfered rectangle xsep=1.4mm, chamfered rectangle ysep=1.2mm,
    fill=unooNavy, draw=unooNavy, text=white, font=\scriptsize\bfseries,
    inner sep=1.8pt, anchor=north west},
  guarda/.style={font=\scriptsize\itshape\color{unooInkSoft}, anchor=west,
    inner sep=1.4pt},
  refint/.style={rectangle, draw=unooNavy, line width=0.85pt, fill=unooSurface,
    font=\scriptsize, text width=40mm, minimum height=11mm, inner sep=3pt},
  % --- comunicacao ---------------------------------------------------------
  msgcomlbl/.style={font=\scriptsize\color{unooInk}, fill=unooSurface,
    inner sep=1.4pt},
  % --- maquina de estados: pseudoestados e estado composto -----------------
  escolha/.style={diamond, draw=unooNavy, line width=0.85pt, fill=unooSurface,
    minimum size=6mm, inner sep=0pt},
  estadocomposto/.style={draw=unooBlue, line width=0.9pt, rounded corners=8pt,
    fill=unooBlue!4},
  % --- tempo ---------------------------------------------------------------
  tempofaixa/.style={draw=unooAxis, line width=0.7pt, fill=unooSurface},
  tempolinha/.style={line width=1.2pt, draw=unooBlue!70!unooNavy},
  tempoestado/.style={font=\scriptsize\color{unooInkSoft}, anchor=east,
    inner sep=1.5pt},
  tempotick/.style={font=\scriptsize\color{unooMuted}, inner sep=1.5pt},
  tempomarca/.style={dashed, dash pattern=on 2pt off 1.6pt, draw=unooAxis,
    line width=0.6pt},
  restricao/.style={font=\scriptsize\color{unooNavy}, fill=unooSurface,
    inner sep=1.4pt},
}

% Titulo de classificador com estereotipo: #1 estereotipo, #2 nome (branco,
% para a primeira particao escura de classe/entidade/interface)
\newcommand{\ctitulo}[2]{\textcolor{white}{{\scriptsize\estereo{#1}}\\[-0.5pt]\textbf{#2}}}

% Pacote: #1 opcoes do corpo, #2 nome do no, #3 e #4 cantos opostos (fit),
% #5 rotulo da aba
\newcommand{\pacote}[5][]{%
  \node[pacotecorpo, fit={#3 #4}, #1] (#2) {};
  \node[pacoteaba, anchor=south west] at (#2.north west) {#5};
}

% No de implantacao (caixa em perspectiva): #1 opcoes do corpo, #2 nome do
% no, #3 e #4 cantos opostos (fit), #5 rotulo (estereotipo + nome)
\newcommand{\nodo}[5][]{%
  \node[nodocorpo, fit={#3 #4}, #1] (#2) {};
  \filldraw[draw=unooNavy, line width=0.85pt, fill=unooNavy!22]
    (#2.north west) -- ++(2.4mm,2.4mm) -- ([shift={(2.4mm,2.4mm)}]#2.north east)
    -- (#2.north east) -- cycle;
  \filldraw[draw=unooNavy, line width=0.85pt, fill=unooNavy!40]
    (#2.north east) -- ++(2.4mm,2.4mm) -- ([shift={(2.4mm,2.4mm)}]#2.south east)
    -- (#2.south east) -- cycle;
  \node[nodolbl, anchor=north west] at (#2.north west) {#5};
}

% Interface fornecida (pirulito): #1 centro da bola
\newcommand{\bola}[1]{%
  \draw[line width=0.85pt, draw=unooInkSoft, fill=unooSurface] #1 circle (1.3mm);}
% Interface requerida (soquete): #1 centro, #2 angulo para onde o soquete ABRE
\newcommand{\soquete}[2]{%
  \draw[line width=0.85pt, draw=unooInkSoft]
    #1 ++(#2+90:2.2mm) arc[start angle=#2+90, end angle=#2+270, radius=2.2mm];}

% Mensagem de diagrama de comunicacao: seta curta ao lado do link.
% #1 posicao do rotulo, #2 ponto medio, #3 angulo da seta, #4 rotulo
\newcommand{\msgcom}[4][above]{%
  \draw[seqmsg] #2 ++(#3+180:4.5mm) -- ++(#3:9mm)
    node[msgcomlbl, midway, #1=1.5pt]{#4};}

\graphicspath{{../../}}

% --- Selos de aderencia arquitetural e operacional (somente tokens unoo*) ----
\newcommand{\selo}[3]{%
  \tikz[baseline=-0.55ex]\node[
    fill=#1, draw=#2, text=#2, line width=0.5pt, rounded corners=2pt,
    inner xsep=4.5pt, inner ysep=1.6pt,
    font=\sffamily\bfseries\fontsize{6.8}{8}\selectfont
  ]{#3};%
}
\newcommand{\nativo}{\selo{unooGreenWash}{unooGreen!75!black}{NATIVO / DIRETO}}
\newcommand{\equivalente}{\selo{unooBlueWash}{unooBlue}{EQUIVALENTE}}
\newcommand{\adaptacao}{\selo{unooOrangeWash}{unooOrange!85!black}{EXIGE ADAPTAÇÃO}}
\newcommand{\incompativel}{\selo{unooRedWash}{unooRed}{INCOMPATÍVEL}}

% Selos de nuvem para leitura rapida nas tabelas e fichas
\newcommand{\sgcp}{\selo{unooBlueWash}{unooBlue}{GCP (ATUAL)}}
\newcommand{\saws}{\selo{unooOrangeWash}{unooOrange!85!black}{AWS}}
\newcommand{\sazure}{\selo{unooLavenderWash}{unooNavy}{AZURE}}

% --- Ficha-resumo de eixo arquitetural (mesma linguagem visual do mapa) ------
\newtcolorbox{fichaeixo}[1]{
  enhanced, breakable,
  colback=unooSurface, colframe=unooLavender,
  boxrule=0.6pt, arc=2pt,
  left=10pt, right=10pt, top=7pt, bottom=7pt,
  before skip=8pt, after skip=12pt,
  title={#1},
  fonttitle=\sffamily\bfseries\small,
  coltitle=white, colbacktitle=unooNavy,
  toptitle=4pt, bottomtitle=4pt, lefttitle=10pt
}

% --- Caixa de componente / comparativo detalhado (estilo historia do mapa) ---
\newtcolorbox{pecaarq}[3]{
  enhanced, breakable,
  colback=white, colframe=unooGrid,
  borderline west={3pt}{0pt}{unooNavy},
  boxrule=0.5pt, arc=2pt,
  left=10pt, right=10pt, top=6pt, bottom=6pt,
  before skip=10pt, after skip=10pt,
  colbacktitle=unooLavenderWash, coltitle=unooNavy,
  fonttitle=\sffamily\bfseries\small,
  toptitle=4pt, bottomtitle=4pt, lefttitle=10pt,
  title={\texttt{#1}\hspace{6pt}#2\hfill#3}
}

% Lista compacta para os campos de uma ficha ou peca arquitetural
\newlist{campos}{description}{1}
\setlist[campos]{
  font=\sffamily\bfseries\footnotesize\color{unooNavy},
  leftmargin=34mm, labelwidth=31mm, labelsep=3mm,
  align=left, itemsep=2.5pt, topsep=3pt, parsep=0pt
}

\begin{document}

% ===========================================================================
%  CAPA
% ===========================================================================
\begin{titlepage}
\begin{tikzpicture}[remember picture, overlay]
  \fill[unooNavy] (current page.north west) rectangle ([yshift=-26mm]current page.north east);
  \fill[unooBlue] ([yshift=-26mm]current page.north west) rectangle ([yshift=-28.5mm]current page.north east);
  \fill[unooLavenderWash] (current page.south west) rectangle ([yshift=22mm]current page.south east);
  \draw[unooLavender, line width=0.6pt]
    ([yshift=22mm]current page.south west) -- ([yshift=22mm]current page.south east);
  \node[anchor=west, text=white, font=\sffamily\bfseries\small]
    at ([xshift=22mm, yshift=-13mm]current page.north west)
    {UNOORIA \enspace\textperiodcentered\enspace ARQUITETURA E ENGENHARIA DE PLATAFORMA};
  \node[anchor=east, text=unooLavenderSoft, font=\sffamily\footnotesize]
    at ([xshift=-22mm, yshift=-13mm]current page.north east)
    {CÓDIGO: PU-EST-002 \enspace|\enspace VERSÃO 1.0};
\end{tikzpicture}

\vspace*{28mm}
\begin{flushleft}
  \includegraphics[width=50mm]{figuras/marca/logo_azul_profundo.png}\\[14mm]
  {\sffamily\bfseries\fontsize{10}{12}\selectfont\color{unooBlue}
    ESTUDO COMPARATIVO DE INFRAESTRUTURA EM NUVEM}\\[4mm]
  {\sffamily\bfseries\fontsize{25}{30}\selectfont\color{unooNavy}
    Plataforma Unooria}\\[3mm]
  {\sffamily\fontsize{14}{19}\selectfont\color{unooInkSoft}
    Análise técnica, operacional e financeira entre \textbf{GCP}, \textbf{AWS} e
    \textbf{Azure} para a Plataforma, as Experiências Imersivas, a API e o Uso de
    Inteligência Artificial}\\[10mm]
  {\color{unooBlue}\rule{36mm}{2pt}}\\[10mm]
\end{flushleft}

\begin{tcolorbox}[
  enhanced, colback=unooSurface, colframe=unooGrid,
  boxrule=0.5pt, arc=2pt, left=10pt, right=10pt, top=8pt, bottom=8pt
]
\footnotesize\renewcommand{\arraystretch}{1.35}
\begin{tabularx}{\linewidth}{@{}L{40mm}X@{}}
  \textbf{\sffamily\color{unooNavy}Escopo do estudo} &
    Comparação exaustiva da pilha de nuvem por contrato B2G (\num{10000}
    estudantes, picos de \num{2000} a \num{4000} simultâneos às 8h da manhã,
    100+ escolas) nos quatro eixos do produto: \textbf{Plataforma} (\arq{apps/hub},
    borda, banco, mensageria, identidade e observabilidade),
    \textbf{Experiências} (\arq{apps/experiencias/*}, bundles estáticos, CDN e
    mídia), \textbf{API} (\arq{apps/api}, computação \emph{serverless}, jobs em
    segundo plano, SSE e WebSockets) e \textbf{Uso de IA} (\arq{src/ai},
    \arq{apps/agente-professor}, busca semântica/estruturada e voz ao vivo).\\
  \textbf{\sffamily\color{unooNavy}Nuvem de referência} &
    Google Cloud Platform (\textbf{GCP}), provisionada via Terraform no
    repositório de fundação (\cd{1 contrato = 1 stack isolada}).\\
  \textbf{\sffamily\color{unooNavy}Nuvens comparadas} &
    Amazon Web Services (\textbf{AWS}) e Microsoft Azure (\textbf{Azure}).\\
  \textbf{\sffamily\color{unooNavy}Restrição de operação} &
    Plataforma desenvolvida e operada por \textbf{uma pessoa} (\emph{simplicidade
    $>$ sofisticação}), com meta de produção em novembro de 2026.\\
  \textbf{\sffamily\color{unooNavy}Data de referência} &
    13 de outubro de 2026\\
\end{tabularx}
\end{tcolorbox}

\vfill
\begin{flushleft}
  \footnotesize\sffamily\color{unooMuted}
  Documento complementar à Documentação Técnica da Plataforma Unooria
  (\arq{documentacao-latex/main.pdf}, Capítulos 6, 11 e 17) e ao Mapa de Produto
  (\cd{PU-PRD-001}).
\end{flushleft}
\vspace*{10mm}
\end{titlepage}

\pagenumbering{roman}
\tableofcontents
\cleardoublepage
\pagenumbering{arabic}

% ===========================================================================
%  CAPITULO 1..7 + APENDICE
% ===========================================================================
% ===========================================================================
%  CAPITULO 1 — Como ler este comparativo e premissas inegociaveis
% ===========================================================================
\chapter{Como ler este comparativo e premissas da plataforma}
\label{cap:como-ler}

Este documento analisa, peça por peça, a infraestrutura em nuvem da
\textbf{Plataforma Unooria} e compara a arquitetura atualmente implantada no
\textbf{Google Cloud Platform (GCP)} com as suas traduções arquiteturais diretas
em \textbf{Amazon Web Services (AWS)} e \textbf{Microsoft Azure}.

Diferente de um comparativo genérico de mercado, este estudo parte das
\textbf{invariantes reais do código e do modelo de negócio} da Unooria ---
documentadas nas ADRs do monorepo e nos módulos Terraform do repositório de
fundação. O objetivo é responder a três perguntas práticas:

\begin{enumerate}
  \item \textbf{Equivalência técnica:} Qual serviço de cada nuvem cumpre
    exatamente o papel das peças atuais da Plataforma, das Experiências, da API
    e do Uso de IA sem violar as regras invioláveis de segurança, \emph{tenancy}
    e tempo de resposta?
  \item \textbf{Custo de operação para uma pessoa:} Onde a troca de nuvem reduz
    ou multiplica a carga cognitiva e operacional de um único desenvolvedor
    mantendo múltiplos contratos B2G?
  \item \textbf{Estrutura de custos (TCO) e acoplamento:} Como se comportam o
    piso fixo mensal por contrato e o custo variável de IA nas três nuvens, e
    quais componentes têm portabilidade imediata versus acoplamento profundo?
\end{enumerate}

\section{Estrutura do documento e convenções visuais}

O estudo está organizado nos quatro eixos solicitados, precedidos por uma visão
executiva de equivalência e seguidos pela análise financeira e de migração:

\begin{itemize}
  \item \textbf{Capítulo~\ref{cap:visao-geral} --- Visão geral e matriz de equivalência:}
    diagrama comparativo da topologia completa e tabela sintética de mapeamento
    serviço a serviço (\cd{GCP} $\times$ \cd{AWS} $\times$ \cd{Azure}).
  \item \textbf{Capítulo~\ref{cap:eixo-plataforma} --- Eixo 1: Plataforma:}
    borda única (Load Balancer, TLS e WAF), hospedagem do \arq{apps/hub}
    (Next.js App Router), banco de dados PostgreSQL com \emph{Row-Level Security}
    (RLS), mensageria de telemetria (\emph{push} autenticado), identidade
    federada (\emph{keyless}) e observabilidade (logs, métricas, alertas e OTLP).
  \item \textbf{Capítulo~\ref{cap:eixo-experiencias} --- Eixo 2: Experiências:}
    entrega de \emph{bundles} estáticos imutáveis por versão (\arq{apps/experiencias/*}),
    estratégia de \emph{cache} e \emph{rollback} em segundos via CDN, isolamento
    de origem para \cd{postMessage}, mídia pesada (\cd{-assets}), guarda
    versionada de planos de aula (\cd{-planos}) e \emph{sourcemaps} privados.
  \item \textbf{Capítulo~\ref{cap:eixo-api} --- Eixo 3: API:}
    runtime \emph{serverless} do \arq{apps/api} (NestJS), alocação contínua de
    CPU para \emph{workers} em segundo plano pós-\cd{HTTP 202}, concorrência por
    instância, \emph{timeouts} longos para \emph{streams} SSE e WebSockets de
    voz ao vivo, saída privada de rede (VPC) e execução de \emph{migrations} via
    \emph{job} isolado.
  \item \textbf{Capítulo~\ref{cap:eixo-ia} --- Eixo 4: Uso de Inteligência Artificial:}
    a camada única (\cd{AiProvider} com texto curto, saída longa de 64k
    \emph{tokens}, JSON validado, geração de imagem com proporção nativa e TTS),
    o Agente de IA do Professor (\arq{apps/agente-professor} em Python/ADK com
    sessões gerenciadas e dupla busca híbrida estruturada/layout com
    \cd{pageSpan}) e a Ajuda por voz bidirecional (\cd{Gemini Live}) e vídeo
    (\cd{LiveAvatar}).
  \item \textbf{Capítulo~\ref{cap:custos-migracao} --- Custos, operação solo e matriz de acoplamento:}
    estimativa comparada do piso fixo mensal por contrato e consumo variável de
    IA, impacto na rotina de operação de uma única pessoa e esforço real de
    portabilidade de cada peça.
\end{itemize}

\subsection{Selos de aderência arquitetural}

Cada componente analisado nos Capítulos~\ref{cap:eixo-plataforma}
a~\ref{cap:eixo-ia} traz uma ficha técnica e selos que classificam o grau de
aderência de cada nuvem aos requisitos da Unooria:

\begin{tabularx}{\linewidth}{@{}L{42mm}X@{}}
\toprule
\headrow \thd{Selo} & \thd{Significado técnico neste estudo}\\
\midrule
\nativo &
  O serviço da nuvem atende ao requisito de forma nativa e direta, com
  configuração declarativa em Terraform e sem exigir peças intermediárias nem
  mudanças no código da aplicação.\\
\zebra \equivalente &
  A nuvem oferece capacidade funcional e de segurança equivalente, mas exige a
  combinação de dois recursos (ex.: fila + função consumidora, ou proxy de
  conexões separado) ou pequenos ajustes em variáveis/adaptadores.\\
\adaptacao &
  A nuvem atende ao caso de uso, porém impõe mudança relevante de arquitetura,
  reescrita de cliente/SDK (ex.: refatorar o agente ADK ou o pipeline de
  \emph{chunking} de PDFs com páginas) ou salto expressivo no piso fixo de custo.\\
\zebra \incompativel &
  O serviço candidato viola uma premissa técnica da plataforma (ex.: limite
  rígido de 15 minutos ou congelamento de CPU no AWS Lambda para conexões
  WebSocket de 60 minutos e \emph{jobs} \emph{in-process} pós-202).\\
\bottomrule
\end{tabularx}

\section{Premissas de negócio e arquitetura que a nuvem precisa respeitar}

Qualquer comparação honesta entre GCP, AWS e Azure precisa julgar os serviços
contra as restrições reais do produto. Uma solução que exige um time de
SRE dedicado ou que encarece o piso ocioso de cada contrato em milhares de
dólares é inviável para o modelo da Unooria, ainda que seja tecnicamente
elegante.

\begin{fichaeixo}{Invariantes de negócio e operação da Plataforma Unooria}
\begin{campos}
  \item[Modelo B2G] \textbf{1 contrato = 1 stack isolada de infraestrutura},
    provisionada por Terraform em projeto/conta própria. Dentro do contrato,
    100+ escolas compartilham a stack com isolamento lógico estrito por
    \cd{escola\_id} (\emph{Row-Level Security} no PostgreSQL).
  \item[Escala por contrato] Cerca de \textbf{\num{10000} estudantes} ativos,
    com picos concentrados de \textbf{\num{2000} a \num{4000} usuários
    simultâneos} no início do turno escolar (8h da manhã) e vales profundos à
    noite e nos fins de semana.
  \item[Equipe de operação] \textbf{Uma pessoa} desenvolve o código, mantém os
    módulos Terraform, responde a incidentes e publica versões. Todo serviço que
    exige manutenção de cluster (como Kubernetes/EKS/AKS/GKE) ou ajustes manuais
    de rede/certificados é descartado em favor de serviços gerenciados
    \emph{serverless}.
  \item[Zero credenciais longas] Nenhuma chave estática de API ou JSON de conta
    de serviço em código, \arq{.env} de produção ou segredos do GitHub. O CI/CD
    autentica por federação OIDC (\emph{Workload Identity}) e os serviços em
    execução usam a identidade nativa do ambiente (ADC / Task Role / Managed
    Identity).
\end{campos}
\end{fichaeixo}

\subsection{Requisitos técnicos determinantes por eixo}

A Tabela~\ref{tab:requisitos-determinantes} sintetiza os requisitos técnicos que
funcionam como \enquote{filtro eliminatório} ao escolher o serviço correspondente
em AWS e Azure.

\begin{table}[H]
\centering
\caption{Requisitos técnicos da Unooria que determinam a escolha dos serviços em nuvem}
\label{tab:requisitos-determinantes}
\begin{tabularx}{\linewidth}{@{}L{28mm}L{44mm}X@{}}
\toprule
\headrow \thd{Eixo} & \thd{Requisito crítico} & \thd{Por que elimina soluções comuns de mercado}\\
\midrule
\textbf{Plataforma} &
  Borda única com roteamento por \emph{hostname}, WAF OWASP e certificados
  validados por DNS sem janela sem TLS &
  O Hub, a API, as Experiências e os Assets respondem sob domínios próprios no
  mesmo balanceador. Adicionar um domínio nunca pode recriar o certificado de
  quem já está no ar (\adrr{ADR-57}, \arq{loadbalancer.tf}).\\
\zebra \textbf{Plataforma} &
  PostgreSQL 16 privado com \cd{FORCE ROW LEVEL SECURITY}, particionamento
  mensal e duas \emph{roles} sem \cd{BYPASSRLS} &
  Descarta bancos NoSQL e exige PostgreSQL gerenciado real, sem IP público, com
  acesso administrativo exclusivamente via túnel autenticado por IAM (IAP/SSM) e
  \emph{connection pooling} para absorver escalonamento de contêineres.\\
\textbf{Plataforma} &
  Fila de telemetria com entrega \emph{push} HTTP autenticada por OIDC e DLQ &
  A API centraliza a ingestão em \rota{POST /telemetry/ingest}. A nuvem deve
  empurrar a mensagem autenticada para a própria API sem exigir um serviço
  \emph{worker} separado rodando 24/7 só para fazer \emph{polling} de fila.\\
\zebra \textbf{Experiências} &
  \emph{Bundles} estáticos imutáveis em \cd{/\{expId\}/\{versao\}/} com origem
  própria no CDN &
  Experiências nunca rodam em contêiner (\emph{regra inviolável}). Precisam de
  domínio dedicado (\cd{exp\_domain}) para que o \cd{targetOrigin} do
  \cd{postMessage} do Hub não compartilhe origem com o restante da nuvem.\\
\textbf{API} &
  Contêiner \emph{serverless} com alta concorrência (80 req/instância), CPU
  sempre alocada e \emph{timeout} de até \num{3600}\,s &
  A geração de planos e pacotes responde \cd{HTTP 202} e segue processando em
  segundo plano no mesmo processo; o Agente do Professor mantém SSE por até
  300\,s; o mascote da Ajuda mantém WebSocket bidirecional por até 60 minutos.\\
\zebra \textbf{Uso de IA} &
  Saída longa (64k \emph{tokens}), controle de orçamento de raciocínio, geração
  de imagem com \emph{aspect ratio} nativo (\cd{4:5}) e voz bidirecional &
  Planos de aula anuais cortam em 32k \emph{tokens}; chamadas rápidas em sala de
  aula (\cd{diagramar}) não podem estourar latência nem orçamento oculto de
  \emph{thinking tokens} (\adrr{ADR-148}).\\
\textbf{Uso de IA} &
  Runtime gerenciado para agente Python (ADK) com sessões persistentes + busca
  estruturada (BNCC) e por \emph{layout} de PDF com \cd{pageSpan} &
  O Agente do Professor precisa listar 100\% das habilidades de um ano por
  filtro exato e citar livro + intervalo exato de páginas (\cd{p. X--Y}) sem
  manter banco vetorial provisionado por hora a custo proibitivo.\\
\bottomrule
\end{tabularx}
\end{table}

\begin{regra}
\textbf{Simplicidade operacional é requisito funcional.} Quando uma nuvem oferece
dois caminhos para o mesmo problema --- por exemplo, rodar a API em Kubernetes
(EKS/AKS) versus rodá-la em um serviço de contêineres \emph{serverless} (ECS
Fargate / Azure Container Apps) ---, este estudo seleciona sempre o caminho
\textbf{gerenciado e sem cluster}, preservando a regra de negócio de que toda a
plataforma é operada por uma única pessoa.
\end{regra}

% ===========================================================================
%  CAPITULO 2 — Visao geral e matriz de equivalencia (GCP x AWS x Azure)
% ===========================================================================
\chapter{Visão geral e matriz de equivalência}
\label{cap:visao-geral}

Antes de mergulhar nas particularidades de cada um dos quatro eixos, este
capítulo apresenta o mapa completo de correspondência arquitetural entre a pilha
atual da Unooria no \textbf{GCP} e as suas arquiteturas equivalentes na
\textbf{AWS} e no \textbf{Azure}.

\begin{fichaeixo}{Síntese executiva da comparação entre as três nuvens}
\begin{campos}
  \item[GCP (Atual)] \nativo\enspace É a nuvem onde o desenho atual atinge a
    \textbf{menor fricção operacional por contrato}: um único Load Balancer
    global agrega tanto contêineres (\emph{Serverless NEGs} do Cloud Run) quanto
    buckets estáticos (\emph{Backend Buckets} do GCS) sob o mesmo IP e WAF; o
    Cloud Run combina alta concorrência (80 req/instância), CPU sempre alocada,
    WebSockets de 60 minutos e \emph{Jobs} de migração no mesmo modelo; o
    Pub/Sub faz \emph{push} HTTP autenticado por OIDC direto para a API; e o
    ecossistema Vertex AI concentra Gemini (64k saída, imagem \cd{4:5}, Live
    voz), Agent Engine (ADK nativo com sessões) e Vertex AI Search (busca
    estruturada + \emph{layout parser} com \cd{pageSpan} cobrado por uso).
  \item[AWS] \adaptacao\enspace Oferece maturidade exemplar em rede, banco de
    dados (RDS Multi-AZ + RDS Proxy) e entrega estática (CloudFront + S3 com
    OAC), mas \textbf{fragmenta peças que no GCP são unificadas}: o ALB não
    aponta para buckets S3 (exigindo CloudFront na frente de tudo para unificar
    borda e WAF); o AWS Lambda é incompatível com os requisitos da API (limite
    de 15 min e modelo de execução), levando ao \textbf{ECS Fargate} (ou App
    Runner com ressalvas de VPC/Jobs); o SQS não faz \emph{push} HTTP nativo
    sem peça intermediária (EventBridge Pipe / Lambda); e no eixo de IA, as
    \emph{Bedrock Knowledge Bases} exigem um banco vetorial subjacente (como
    OpenSearch Serverless, com piso fixo elevado em OCUs, ou Aurora pgvector) e
    reescrita do agente ADK.
  \item[Azure] \equivalente\enspace É a alternativa \textbf{mais próxima do
    modelo operacional do Cloud Run} para computação e borda: o \textbf{Azure
    Front Door} roteia nativamente para contêineres e Blob Storage com WAF
    integrado; o \textbf{Azure Container Apps} (ACA) entrega experiência muito
    similar ao Cloud Run (concorrência HTTP, \emph{scale-to-zero} ou réplicas
    mínimas, sessões WebSocket longas e \emph{Container Apps Jobs} para o
    \cd{prisma migrate deploy}); e o PostgreSQL Flexible Server já traz
    \emph{PgBouncer} embutido sem custo extra. No eixo de IA, combina Azure
    OpenAI (\cd{GPT-4o} + \cd{Realtime API}) e \textbf{Azure AI Search} com
    \emph{Document Intelligence}, mas ainda exige portar o código do agente ADK
    e gerenciar índices cobrados por hora.
\end{campos}
\end{fichaeixo}

\section{Topologia comparada nas três nuvens}

A Figura~\ref{fig:topologia-comparada} ilustra como os quatro eixos da
Plataforma Unooria (Borda/Plataforma, Experiências, API/Dados e Uso de IA) se
organizam em cada provedor de nuvem. O código Mermaid correspondente encontra-se
no Apêndice~\ref{ap:mermaid} (Listagem~\ref{lst:mermaid-topologia}).

\begin{figure}[H]
\centering
\fitwidth{%
\begin{tikzpicture}[x=1cm, y=1cm,
  every node/.append style={execute at begin node={\hyphenpenalty=10000\exhyphenpenalty=10000}},
  basebox/.style={rectangle, rounded corners=2pt, line width=0.7pt, inner sep=4pt, align=center},
  colhdr/.style={basebox, draw=none, font=\sffamily\bfseries\footnotesize, text=white, minimum height=9mm, text width=50mm},
  rowlbl/.style={basebox, fill=unooLavenderWash, draw=unooLavender, text=unooNavy, font=\sffamily\bfseries\scriptsize, text width=24mm, minimum height=19mm},
  cellgcp/.style={basebox, fill=unooBlueWash, draw=unooBlue, font=\scriptsize, text width=50mm, minimum height=19mm},
  cellaws/.style={basebox, fill=unooOrangeWash, draw=unooOrange!85!black, font=\scriptsize, text width=50mm, minimum height=19mm},
  cellaz/.style={basebox, fill=unooSurface, draw=unooNavy, font=\scriptsize, text width=50mm, minimum height=19mm}]

  % Cabecalhos das colunas
  \node[colhdr, fill=unooBlue]            (hgcp) at (4.5,  0) {GCP (Arquitetura Atual)};
  \node[colhdr, fill=unooOrange!85!black] (haws) at (10.0, 0) {AWS (Equivalente)};
  \node[colhdr, fill=unooNavy]            (haz)  at (15.5, 0) {Azure (Equivalente)};

  % Linha 1: Borda e Entrega Estatica (Eixos 1 e 2)
  \node[rowlbl]  (r1)   at (0.4, -1.65) {1. Borda, WAF\\e Experiências\\(Eixos 1 e 2)};
  \node[cellgcp] (g1)   at (4.5, -1.65) {\textbf{Global External ALB + Cloud Armor}\\Cert Manager DNS + Cloud CDN\\\texttt{Backend Buckets}: \texttt{-exp}, \texttt{-assets}};
  \node[cellaws] (w1)   at (10.0,-1.65) {\textbf{CloudFront + AWS WAF + ACM}\\Roteia estático p/ \textbf{S3 (OAC)} e\\dinâmico p/ \textbf{ALB} (\texttt{hub} e \texttt{api})};
  \node[cellaz]  (z1)   at (15.5,-1.65) {\textbf{Azure Front Door + WAF}\\Certificados gerenciados + CDN\\Origens: \textbf{Blob Storage} e \textbf{ACA}};

  % Linha 2: Computacao Serverless (Eixos 1 e 3)
  \node[rowlbl]  (r2)   at (0.4, -3.95) {2. Hub, API\\e Migrations\\(Eixos 1 e 3)};
  \node[cellgcp] (g2)   at (4.5, -3.95) {\textbf{Cloud Run Services + Jobs}\\\texttt{hub} e \texttt{api} (CPU 24/7, conc. 80,\\\texttt{timeout} 900--3600s) + Job \texttt{migrate}};
  \node[cellaws] (w2)   at (10.0,-3.95) {\textbf{ECS Fargate} (Service + Task)\\\texttt{hub} e \texttt{api} em Fargate atrás do ALB\\+ Task avulsa p/ \texttt{prisma migrate}};
  \node[cellaz]  (z2)   at (15.5,-3.95) {\textbf{Azure Container Apps (ACA)}\\Apps p/ \texttt{hub} e \texttt{api} (KEDA, Envoy\\timeout longo) + \textbf{ACA Job} \texttt{migrate}};

  % Linha 3: Dados, Mensageria e Rede Privada (Eixo 1)
  \node[rowlbl]  (r3)   at (0.4, -6.25) {3. Banco, Fila\\e Rede Privada\\(Eixo 1)};
  \node[cellgcp] (g3)   at (4.5, -6.25) {\textbf{Cloud SQL PG 16} (HA, IP privado)\\Direct VPC Egress + Bastion \textbf{IAP}\\\textbf{Pub/Sub} Push OIDC \seta{} \texttt{/ingest}};
  \node[cellaws] (w3)   at (10.0,-6.25) {\textbf{RDS PG 16 Multi-AZ + RDS Proxy}\\Subnets privadas + Bastion \textbf{SSM}\\\textbf{SQS + EventBridge Pipe} \seta{} API};
  \node[cellaz]  (z3)   at (15.5,-6.25) {\textbf{Flexible Server PG 16} (PgBouncer)\\VNet Integration + \textbf{Azure Bastion}\\\textbf{Event Grid / Service Bus} \seta{} API};

  % Linha 4: Camada de IA, Agente e Voz (Eixo 4)
  \node[rowlbl]  (r4)   at (0.4, -8.55) {4. Modelos, RAG,\\Agente e Voz\\(Eixo 4)};
  \node[cellgcp] (g4)   at (4.5, -8.55) {\textbf{Vertex AI}: Gemini 3.6/3.7, Imagen,\\Cloud TTS e \textbf{Gemini Live} (WSS)\\\textbf{Agent Engine} + \textbf{Vertex AI Search}};
  \node[cellaws] (w4)   at (10.0,-8.55) {\textbf{Amazon Bedrock}: Claude/Nova, Titan,\\Polly e \textbf{Nova Sonic} (voz)\\\textbf{AgentCore} + \textbf{Knowledge Bases}};
  \node[cellaz]  (z4)   at (15.5,-8.55) {\textbf{Azure AI Foundry}: GPT-4o, DALL-E,\\AI Speech e \textbf{GPT-4o Realtime}\\\textbf{AI Agent Service} + \textbf{AI Search}};

  % Setas verticais entre camadas por coluna
  \foreach \c in {g, w, z} {
    \draw[flow] (\c1.south) -- (\c2.north);
    \draw[flow] (\c2.south) -- (\c3.north);
    \draw[flow] (\c2.south) -- (\c4.north);
  }
\end{tikzpicture}}
\caption[Topologia comparada GCP, AWS e Azure]{Correspondência estrutural das
  quatro camadas entre GCP, AWS e Azure para um contrato da Plataforma Unooria.}
\label{fig:topologia-comparada}
\end{figure}

\section{Matriz completa de equivalência serviço a serviço}

A Tabela~\ref{tab:matriz-equivalencia} mapeia todos os recursos declarados nos
arquivos \arq{.tf} do repositório \arq{Fundacao Plataforma Unooria} para os seus
equivalentes diretos em AWS e Azure.

\begin{table}[H]
\centering
\caption{Mapeamento completo dos serviços da Plataforma Unooria entre GCP, AWS e Azure}
\label{tab:matriz-equivalencia}
\footnotesize
\begin{tabularx}{\linewidth}{@{}L{29mm}L{38mm}L{38mm}X@{}}
\toprule
\headrow \thd{Função na Unooria} & \thd{GCP (Atual)} & \thd{AWS} & \thd{Azure}\\
\midrule
\multicolumn{4}{@{}l}{\textbf{\sffamily\color{unooNavy}Eixo 1 --- Plataforma (Borda, Hub, Banco, Rede, Mensageria, Segurança e Observabilidade)}}\\
\midrule
Isolamento por contrato &
  1 GCP Project por ambiente/contrato &
  1 AWS Account (via AWS Organizations) &
  1 Subscription ou Resource Group isolado\\
\zebra Entrada única HTTP/S &
  Global External Application Load Balancer &
  Amazon CloudFront + Application Load Balancer (ALB) &
  Azure Front Door (Standard/Premium)\\
WAF e \emph{rate limit} &
  Cloud Armor Standard (OWASP CRS 3.3 + RPM/IP) &
  AWS WAF v2 (\emph{Managed Rules} + \emph{Rate-based rule}) &
  Azure WAF no Front Door (OWASP DRS + \emph{Rate limit})\\
\zebra Certificados TLS &
  Certificate Manager (validação DNS por host) &
  AWS Certificate Manager --- ACM (validação DNS) &
  Front Door Managed Certificates / Key Vault\\
Hospedagem do \cd{hub} &
  Cloud Run Service (\cd{1--2 vCPU}, \cd{min=1}) &
  ECS Fargate Service (ou AWS App Runner) &
  Azure Container Apps --- ACA (\cd{minReplicas=1})\\
\zebra Banco relacional + RLS &
  Cloud SQL for PostgreSQL 16 (Regional HA, IP privado) &
  Amazon RDS for PostgreSQL 16 (Multi-AZ) + RDS Proxy &
  Azure Database for PostgreSQL Flexible Server (HA + PgBouncer)\\
Rede privada e saída &
  VPC + Direct VPC Egress + Cloud NAT + PSA &
  VPC (Subnets privadas) + NAT Gateway &
  Virtual Network (VNet) + VNet Integration + NAT Gateway\\
\zebra Acesso dev ao banco &
  VM Bastion (\cd{e2-micro}) sem IP público + túnel IAP &
  Instância \cd{t4g.nano} privada + SSM Session Manager &
  Azure Bastion (\emph{Developer} / \emph{Basic}) ou ACA túnel\\
Fila de telemetria (\emph{push}) &
  Cloud Pub/Sub (tópico + \emph{Push} OIDC + DLQ) &
  Amazon SQS + EventBridge Pipe (\emph{API Destination}) + DLQ &
  Azure Event Grid (\emph{Push} c/ Entra ID) ou Service Bus\\
\zebra Segredos de runtime &
  Secret Manager (injetado como \emph{env} no boot) &
  AWS Secrets Manager / SSM Parameter Store &
  Azure Key Vault (referenciado no Container App)\\
Registro de imagens &
  Artifact Registry (limpeza em 7 dias, 20 tags) &
  Amazon ECR (\emph{Lifecycle Policy} equivalente) &
  Azure Container Registry --- ACR (retenção por política)\\
\zebra Federação CI/CD &
  Workload Identity Federation (OIDC GitHub) &
  IAM OIDC Identity Provider (\cd{sts:\B{}AssumeRoleWithWebIdentity}) &
  Entra ID Workload Identity Federation (\emph{Federated Creds})\\
Observabilidade e OTLP &
  Cloud Logging, Monitoring, Trace, Error Reporting e API OTLP &
  CloudWatch (Logs, Metrics, Synthetics, Alarms) + X-Ray / ADOT &
  Azure Monitor + Log Analytics + Application Insights\\
\midrule
\multicolumn{4}{@{}l}{\textbf{\sffamily\color{unooNavy}Eixo 2 --- Experiências (Bundles Estáticos, CDN, Assets e Planos)}}\\
\midrule
Bundles imutáveis (\cd{-exp}) &
  Cloud Storage + Backend Bucket no LB + Cloud CDN &
  Amazon S3 + CloudFront (\emph{Origin Access Control}) &
  Azure Blob Storage + Azure Front Door CDN\\
\zebra Mídia pública (\cd{-assets}) &
  Cloud Storage público + Cloud CDN (\cd{assets\_domain}) &
  Amazon S3 + CloudFront (\cd{assets\_domain}) &
  Azure Blob Storage + Front Door CDN\\
Planos e áudios (\cd{-planos}) &
  Cloud Storage privado com \emph{Object Versioning} (180d) &
  Amazon S3 privado com \emph{S3 Versioning} + \emph{Lifecycle} &
  Azure Blob Storage privado com \emph{Blob Versioning}\\
\zebra Mapas (\cd{-exp-mapas}) &
  Cloud Storage privado (leitura restrita) &
  Amazon S3 privado (leitura restrita via IAM) &
  Azure Blob Storage privado\\
\midrule
\multicolumn{4}{@{}l}{\textbf{\sffamily\color{unooNavy}Eixo 3 --- API (Computação Serverless, Jobs em Fundo, SSE e WebSockets)}}\\
\midrule
Hospedagem da \cd{api} &
  Cloud Run (\cd{cpu\_idle=false}, conc. 80, \cd{timeout} 900s) &
  ECS Fargate (\cd{2 vCPU/4 GB}, ALB \cd{idle\_timeout} alto) &
  Azure Container Apps (\emph{Dedicated/Consumption}, Envoy 900s+)\\
\zebra WebSockets longos (60m) &
  Cloud Run (\cd{timeout} \num{3600}\,s nativo p/ WS) &
  ECS Fargate + ALB (ajuste de \cd{idle\_timeout} até \num{4000}\,s) &
  Azure Container Apps (suporte nativo a WS longos)\\
Job de migração SQL &
  Cloud Run Job (\cd{\{contrato\}-migrate} na VPC) &
  ECS Fargate Standalone Task (\cd{RunTask} no CI) &
  Azure Container Apps Job (\emph{Manual trigger} no CI)\\
\midrule
\multicolumn{4}{@{}l}{\textbf{\sffamily\color{unooNavy}Eixo 4 --- Uso de Inteligência Artificial (Modelos, Agente, RAG e Voz)}}\\
\midrule
LLM texto e saída longa &
  Vertex AI (\cd{gemini-3.6-flash} / \cd{3.7-flash}, 64k saída) &
  Amazon Bedrock (\cd{Claude 3.7/4 Sonnet}, \cd{Amazon Nova Pro/Lite}) &
  Azure OpenAI / AI Foundry (\cd{GPT-4o}, \cd{GPT-4.1}, \cd{o3-mini})\\
\zebra Geração de imagem &
  Vertex AI (\cd{gemini-3.1-flash-image}, ratio \cd{4:5} nativo) &
  Amazon Bedrock (\cd{Amazon Nova Canvas} / \cd{Titan Image}) &
  Azure OpenAI (\cd{gpt-image-1} / \cd{DALL-E 3})\\
Síntese de voz (TTS) &
  Cloud Text-to-Speech (\cd{gemini-2.5-flash-tts}, multi-voz) &
  Amazon Polly (\emph{Generative / Neural Voices}) &
  Azure AI Speech (\emph{Neural HD Voices} + SSML multi-voz)\\
\zebra Voz bidirecional (Mascote) &
  Vertex AI \textbf{Gemini Live} (\cd{gemini-3.8-live} via WSS) &
  Amazon Bedrock \textbf{Nova Sonic} (stream bidirecional) &
  Azure OpenAI \textbf{GPT-4o Realtime API} (WSS/WebRTC)\\
Runtime do Agente &
  Vertex AI \textbf{Agent Engine} (\cd{ReasoningEngine} BYOC + sessões) &
  \textbf{Bedrock AgentCore} (ou agente em ECS Fargate + DynamoDB) &
  \textbf{Azure AI Agent Service} (ou ACA + Cosmos DB/Postgres)\\
\zebra Busca RAG (BNCC + PDFs) &
  \textbf{Vertex AI Search} (1 estruturado + 1 \emph{layout chunking} c/ página) &
  \textbf{Bedrock Knowledge Bases} + OpenSearch Serverless / Aurora &
  \textbf{Azure AI Search} + \textbf{Document Intelligence} (\emph{Layout})\\
\bottomrule
\end{tabularx}
\end{table}

% ===========================================================================
%  CAPITULO 3 — Eixo 1: Plataforma (Borda, Hub, Banco, Rede, Mensageria, Obs)
% ===========================================================================
\chapter{Eixo 1 --- Plataforma (Borda, Hub, Banco, Mensageria e Operação)}
\label{cap:eixo-plataforma}

O eixo \textbf{Plataforma} compreende todas as fundações compartilhadas que
mantêm um contrato B2G no ar: a borda pública com proteção WAF e gestão de
certificados TLS, a hospedagem do \arq{apps/hub} (Next.js App Router), o banco
de dados relacional com isolamento por escola (\emph{Row-Level Security}), a
rede privada, o barramento assíncrono de telemetria, a gestão de identidades sem
chaves estáticas e a pilha de observabilidade.

\begin{fichaeixo}{Ficha do Eixo 1 --- Fundação da Plataforma}
\begin{campos}
  \item[Arquivos Terraform] \arq{loadbalancer.tf}, \arq{cloudarmor.tf},
    \arq{cloudrun.tf}, \arq{database.tf}, \arq{network.tf}, \arq{bastion.tf},
    \arq{pubsub.tf}, \arq{iam\_github.tf}, \arq{secrets.tf}, \arq{registry.tf}
    e \arq{observabilidade.tf}.
  \item[Aderência por nuvem] \sgcp~\nativo \quad \saws~\equivalente \quad
    \sazure~\nativo
  \item[Ponto focal] No GCP e no Azure (via Front Door), um único recurso de
    borda unifica o roteamento por \emph{hostname} para contêineres e buckets
    estáticos sob a mesma política de WAF. Na AWS, o ALB não roteia para S3, o
    que exige colocar o \textbf{Amazon CloudFront} como camada de entrada única
    diante tanto do S3 quanto do ALB.
\end{campos}
\end{fichaeixo}

\section{Borda única, WAF e gestão de certificados TLS}

\begin{pecaarq}{PLT-01}{Load Balancer Global, Certificados por DNS e WAF OWASP}{\sgcp~\nativo\enspace\saws~\equivalente\enspace\sazure~\nativo}
\begin{campos}
  \item[Como é hoje (GCP)] Um \textbf{Global External Application Load Balancer}
    detém um único IP anycast estático e um \cd{url\_map} que distribui o
    tráfego por \emph{hostname}: \cd{hub\_domain} \seta{} Cloud Run \cd{frontend},
    \cd{api\_domain} \seta{} Cloud Run \cd{api}, \cd{exp\_domain} \seta{}
    Backend Bucket \cd{-exp} com Cloud CDN, e \cd{assets\_domain} \seta{}
    Backend Bucket \cd{-assets} com Cloud CDN. A proteção de borda usa o
    \textbf{Cloud Armor Standard} anexado aos backends de computação, aplicando
    regras pré-configuradas \textbf{OWASP CRS 3.3} (\cd{sqli-v33-stable},
    \cd{xss-v33-stable}, \cd{lfi-v33-stable}, \cd{rfi-v33-stable},
    \cd{rce-v33-stable}, \cd{scannerdetection-v33-stable},
    \cd{protocolattack-v33-stable}, \cd{sessionfixation-v33-stable}) em nível de
    sensibilidade 2, além de \emph{rate limit} por IP (\cd{300} a \cd{2000} RPM
    conforme o ambiente, retornando HTTP 429). Os certificados TLS usam o
    \textbf{Certificate Manager} com validação por DNS (\cd{\_acme-challenge}):
    cada \emph{hostname} tem certificado próprio em um \cd{Certificate Map}, o
    que permite adicionar novos domínios sem tocar nos certificados já ativos e
    provisionar o TLS \textbf{antes} de virar o apontamento DNS (\arq{loadbalancer.tf}).
  \item[Equivalente AWS] \textbf{Amazon CloudFront} + \textbf{AWS WAF v2} +
    \textbf{AWS Certificate Manager (ACM)} + \textbf{Application Load Balancer
    (ALB)}. Como o ALB da AWS não aceita um bucket S3 como \emph{Target Group}
    para servir arquivos estáticos diretamente, a borda unificada precisa ser o
    CloudFront (ou múltiplos ALBs/distribuições separadas). No CloudFront,
    associa-se uma WebACL do AWS WAF (\emph{AWSManagedRulesCommonRuleSet},
    \emph{SQLiRuleSet}, \emph{KnownBadInputsRuleSet} e regra \emph{Rate-based}
    por IP). O ACM emite certificados gratuitos validados por CNAME no DNS
    (\cd{us-east-1} para CloudFront), também sem janela de indisponibilidade.
    Para impedir que alguém contorne o WAF acessando o ALB diretamente, o ALB
    exige cabeçalho secreto rotacionado ou \emph{prefix list} gerenciada do
    CloudFront no Security Group.
  \item[Equivalente Azure] \textbf{Azure Front Door (Standard ou Premium)} +
    \textbf{Azure WAF Policy} + \textbf{Managed Certificates}. O modelo do Front
    Door é arquiteturalmente idêntico ao do Global External LB do GCP: um ponto
    global anycast que aceita múltiplos domínios validados por CNAME de DNS
    (\emph{pre-validated custom domains} sem tempo fora do ar) e roteia, por
    \emph{hostname}, tanto para origens de contêiner (Azure Container Apps)
    quanto para origens de armazenamento (Azure Blob Storage), com cache CDN e
    regras OWASP DRS + \emph{rate limit} na mesma borda.
\end{campos}
\end{pecaarq}

\begin{tabularx}{\linewidth}{@{}L{32mm}L{40mm}L{40mm}X@{}}
\toprule
\headrow \thd{Critério de borda} & \thd{GCP (LB + Cloud Armor)} & \thd{AWS (CloudFront + ALB + WAF)} & \thd{Azure (Front Door + WAF)}\\
\midrule
Contêiner + Bucket no mesmo recurso &
  \textbf{Sim} (\emph{Serverless NEG} + \emph{Backend Bucket}) &
  \textbf{Só via CloudFront} (exige ALB adicional para Fargate) &
  \textbf{Sim} (origens ACA e Blob no mesmo perfil)\\
\zebra Fechamento do endereço direto do contêiner &
  Nativo (\cd{INGRESS\_\B{}TRAFFIC\_\B{}INTERNAL\_\B{}LOAD\_\B{}BALANCER}) &
  Exige Security Group restrito à \emph{prefix list} do CloudFront + ALB interno &
  Nativo no ACA com VNet + \emph{Private Link} no Front Door Premium\\
Adição de domínio sem recriar cert ativo &
  \textbf{Sim} (\cd{Certificate Map} com 1 cert por host via DNS) &
  \textbf{Sim} (certificados ACM independentes por domínio/SAN DNS) &
  \textbf{Sim} (\emph{Managed Certificate} individual por domínio no Front Door)\\
\zebra Modo \emph{preview} de regras WAF &
  Nativo (\cd{preview = true} por regra ou na política inteira) &
  Nativo (ação \cd{Count} em cada regra gerenciada do AWS WAF) &
  Nativo (modo \cd{Detection} vs \cd{Prevention} na WAF Policy)\\
\bottomrule
\end{tabularx}

\section{Hospedagem do Hub (\texttt{apps/hub}) e Identidade (OIDC)}

\begin{pecaarq}{PLT-02}{Hub Next.js (SSR + BFF) e Federação OIDC / JWKS}{\sgcp~\nativo\enspace\saws~\equivalente\enspace\sazure~\nativo}
\begin{campos}
  \item[Como é hoje (GCP)] O \arq{apps/hub} (Next.js App Router, saída
    \cd{standalone}) roda como serviço no \textbf{Cloud Run} (\cd{1} a \cd{2
    vCPU}, \cd{1} a \cd{2 GiB} de RAM, concorrência 80, \cd{min\_instances = 1}
    para eliminar \emph{cold start} no horário de entrada escolar e manter o
    cache de sessão quente). Ele atua como BFF, autentica o usuário contra o
    \textbf{WSO2 Identity Server} via OIDC padrão, assina em memória os tokens
    de curta duração (15 min, RS256) das experiências com a chave privada injetada
    do Secret Manager e publica a chave pública em \rota{/.well-known/jwks.json}.
    A API busca esse JWKS pelo endereço do Hub (\env{OIDC\_JWKS\_URI}).
  \item[Equivalente AWS] \textbf{Amazon ECS on AWS Fargate} (ou \textbf{AWS App
    Runner}). Embora existam soluções específicas para Next.js (como AWS
    Amplify Hosting), o Hub da Unooria funciona como emissor OIDC (\cd{issuer})
    com chaves montadas via segredo e comunicação controlada com a API; mantê-lo
    como contêiner OCI padrão (`output: 'standalone'`) no ECS Fargate preserva
    100\% de paridade com o Dockerfile existente. Uma tarefa Fargate (\cd{1
    vCPU / 2 GB}, \cd{desired\_count >= 1}) atrás do ALB cumpre exatamente o
    papel do Cloud Run.
  \item[Equivalente Azure] \textbf{Azure Container Apps (ACA)}. Recebe a mesma
    imagem Docker do Hub sem nenhuma alteração, mantendo \cd{minReplicas = 1},
    escalonamento automático por concorrência HTTP (via KEDA embutido) e injeção
    de segredos do Azure Key Vault como variáveis de ambiente.
\end{campos}
\end{pecaarq}

\begin{atencao}
\textbf{Armadilha de roteamento entre Hub e API na nuvem.} No GCP, quando o Hub
tem saída de rede configurada como \cd{PRIVATE\_RANGES\_ONLY} e tenta chamar a
API pelo endereço \cd{*.run.app}, a borda do Google responde 404 (\arq{variables.tf},
linha 240). Por isso o Hub chama a API pelo domínio público do Load Balancer (ou
exige \cd{ALL\_TRAFFIC} via Cloud NAT). Na \textbf{AWS} (ECS Fargate) e no
\textbf{Azure} (ACA), ambos os serviços vivem dentro da mesma VPC/VNet e podem se
comunicar diretamente por DNS interno (\emph{ECS Service Connect} / \emph{Cloud
Map} na AWS, ou o FQDN interno \cd{*.internal.<env>.azurecontainerapps.io} no
Azure), economizando tráfego de NAT Gateway e evitando que as chamadas BFF
contem no \emph{rate limit} por IP do WAF externo.
\end{atencao}

\section{Banco de Dados Relacional, RLS, Rede Privada e Acesso Dev}

O coração da segurança multi-escola da Unooria mora no PostgreSQL 16: todas as
tabelas com dados de estudante ou escola possuem \cd{escola\_id NOT NULL} e
políticas \cd{FORCE ROW LEVEL SECURITY}, com duas \emph{roles} distintas em
execução (\cd{app\_runtime} para o tráfego isolado por \cd{SET LOCAL
app.escola\_id} dentro de \cd{comTenancy}, e \cd{app\_auth} restrita à resolução
de login antes de saber a escola), nenhuma das duas com atributo \cd{BYPASSRLS}
nem \cd{SUPERUSER}. Além disso, a tabela \cd{telemetria\_eventos} usa
particionamento declarativo mensal.

\begin{pecaarq}{PLT-03}{PostgreSQL 16 Privado, Alta Disponibilidade, Pooler e Bastion}{\sgcp~\nativo\enspace\saws~\equivalente\enspace\sazure~\nativo}
\begin{campos}
  \item[Como é hoje (GCP)] \textbf{Cloud SQL for PostgreSQL 16} (\cd{db-custom-2-7680}
    em produção, \cd{availability\_type = REGIONAL} com \emph{failover}
    automático entre zonas, backups diários + \emph{Point-in-Time Recovery} via
    WAL, sem IPv4 público). A conectividade usa \emph{Private Service Access}
    (peering VPC) e os serviços Cloud Run chegam ao banco via \textbf{Direct VPC
    Egress}. Para acesso de desenvolvimento e aplicação do script de
    endurecimento de \emph{roles}, há uma VM \textbf{Bastion} (\cd{e2-micro},
    incluída no \emph{free tier} em \cd{us-east1}) sem IP público, acessada via
    túnel TCP do \textbf{Identity-Aware Proxy (IAP)} autenticado pelo IAM.
  \item[Equivalente AWS] \textbf{Amazon RDS for PostgreSQL 16 (Multi-AZ)} em
    subnets privadas (\cd{db.t4g.large} ou \cd{db.m7g.large}). Para evitar o
    esgotamento de conexões quando tarefas Fargate escalam, a AWS utiliza o
    \textbf{Amazon RDS Proxy} (cobrado por vCPU do banco). O acesso seguro da
    máquina do desenvolvedor ao RDS privado é feito de forma idêntica ao IAP por
    meio de uma instância \cd{t4g.nano} sem IP público usando o \textbf{AWS
    Systems Manager (SSM) Session Manager} (\cd{AWS-StartPortForwardingSessionToRemoteHost}),
    autenticado via IAM e sem abrir a porta 22 para a internet.
  \item[Equivalente Azure] \textbf{Azure Database for PostgreSQL --- Flexible
    Server} (versão 16, modo \emph{Zone-redundant HA}, integrado à VNet por
    delegação de subnet ou \emph{Private Endpoint}). Grande vantagem operacional
    e de custo do Azure: o Flexible Server possui \textbf{PgBouncer embutido} na
    própria instância (porta \cd{6432}), ativado por uma flag no Terraform sem
    custo adicional de proxy separado. Para o túnel de desenvolvimento, pode-se
    usar o \textbf{Azure Bastion} (SKU \emph{Standard} para túnel nativo via
    \cd{az network bastion tunnel}, ou uma VM \cd{B1ls} com acesso JIT/Entra ID).
\end{campos}
\end{pecaarq}

\begin{nota}
\textbf{Particularidade de \emph{roles} gerenciadas nos três provedores.} No
Cloud SQL (GCP), todo usuário criado pela Admin API recebe \cd{cloudsqlsuperuser};
no RDS (AWS), recebe \cd{rds\_superuser}; no Flexible Server (Azure), recebe
\cd{azure\_pg\_admin}. Nos três provedores, o passo SQL de endurecimento das
contas de execução (\cd{NOSUPERUSER NOBYPASSRLS NOCREATEDB NOCREATEROLE})
permanece idêntico e portável, pois opera dentro do motor padrão do PostgreSQL 16.
\end{nota}

\section{Mensageria de Telemetria Analítica (Push Autenticado)}

As experiências enviam lotes de eventos pedagógicos para \rota{POST
/telemetry/batch}. A API valida o envelope com Zod, enriquece com dados do JWT
(\cd{escola\_id}, \cd{estudante\_id} opaco) e publica no barramento em menos de
15\,ms, respondendo \cd{HTTP 202}. O armazenamento em lote no banco particionado
acontece em \rota{POST /telemetry/ingest}, que é uma rota da \textbf{própria API}
protegida por token OIDC da conta de serviço do barramento --- evitando criar e
monitorar um serviço \emph{worker} separado (\adrr{ADR-05}).

\begin{pecaarq}{PLT-04}{Fila de Telemetria com Entrega Push OIDC e Fila Morta (DLQ)}{\sgcp~\nativo\enspace\saws~\adaptacao\enspace\sazure~\equivalente}
\begin{campos}
  \item[Como é hoje (GCP)] \textbf{Cloud Pub/Sub}: um tópico principal
    (\cd{\{contrato\}-telemetria}), uma assinatura \textbf{Push} que faz
    \cd{POST} direto em \cd{https://\{api\_domain\}/telemetry/ingest} assinando
    um JWT OIDC da conta de serviço \cd{pubsub-push} (validado pelo
    \cd{PubsubPushGuard} da API via JWKS do Google), política de \emph{backoff}
    exponencial (10\,s a 600\,s) e redirecionamento para o tópico de fila morta
    (\cd{\{contrato\}-telemetria-dlq}) após 5 tentativas falhas (\arq{pubsub.tf}).
  \item[Equivalente AWS] O \textbf{Amazon SQS} padrão é puramente baseado em
    \emph{pull} (não faz \emph{push} HTTP), e o \textbf{Amazon SNS} faz
    \emph{push} HTTP mas não possui fila persistente longa nem assina tokens OIDC
    dinâmicos nativamente (e exige um passo de confirmação de assinatura
    \cd{SubscriptionConfirmation} no endpoint). Para reproduzir o mesmo desenho
    na AWS sem manter um \emph{poller} no código da API, há dois caminhos:
    \begin{enumerate}
      \item \textbf{SQS + EventBridge Pipes + API Destinations:} a API publica no
        SQS (com DLQ nativa após 5 tentativas); um \emph{EventBridge Pipe} lê da
        fila e faz o \cd{POST} na rota \rota{/telemetry/ingest} autenticando via
        \emph{EventBridge Connection} (OAuth/OIDC ou chave simétrica no
        Secrets Manager).
      \item \textbf{SQS Consumer embutido na API:} como a API na AWS já roda com
        no mínimo 1 tarefa Fargate ligada 24/7, adicionar um consumidor SQS em
        segundo plano dentro do módulo \arq{telemetry/} da API elimina o
        EventBridge e mantém o tráfego 100\% interno na VPC.
    \end{enumerate}
  \item[Equivalente Azure] \textbf{Azure Event Grid} (modelo \emph{Push} nativo
    com \emph{Dead-Lettering} em Blob Storage e entrega HTTP autenticada com
    token \textbf{Microsoft Entra ID / OIDC} gerado por \emph{Managed Identity})
    ou \textbf{Azure Service Bus} (filas mais robustas, porém baseadas em
    \emph{pull} via SDK ou acionadas via Dapr/ACA). O Event Grid Push com
    autenticação Entra ID é o espelho funcional exato da assinatura Push OIDC do
    Pub/Sub.
\end{campos}
\end{pecaarq}

\section{Identidade Federada (CI/CD), Segredos e Observabilidade}

\begin{pecaarq}{PLT-05}{CI/CD Keyless, Segredos e Observabilidade OTLP}{\sgcp~\nativo\enspace\saws~\nativo\enspace\sazure~\nativo}
\begin{campos}
  \item[Identidade e Segredos] Nas três nuvens, o padrão de \textbf{zero chaves
    de longa duração} é suportado de primeira classe:
    \begin{itemize}
      \item \textbf{GCP:} \emph{Workload Identity Federation} valida o token OIDC
        de \cd{token.actions.githubusercontent.com} restrito a
        \cd{assertion.repository}; segredos ficam no \textbf{Secret Manager};
        imagens no \textbf{Artifact Registry}.
      \item \textbf{AWS:} \emph{IAM OIDC Identity Provider} confia no GitHub
        Actions via condição \cd{token.actions.githubusercontent.com:sub};
        segredos no \textbf{AWS Secrets Manager} (injetados no bloco \cd{secrets}
        da Task Definition do ECS); imagens no \textbf{Amazon ECR}.
      \item \textbf{Azure:} \emph{Federated Identity Credentials} em uma
        \emph{User-Assigned Managed Identity} no Entra ID; segredos no
        \textbf{Azure Key Vault}; imagens no \textbf{Azure Container Registry
        (ACR)}.
    \end{itemize}
  \item[Observabilidade e OTLP] A plataforma instrumenta logs estruturados
    (\cd{pino}), métricas customizadas, 13 políticas de alerta, SLO de 99,5\% e
    reexporta os \emph{spans} OpenTelemetry do navegador (enviados pelo SDK das
    experiências para a API) para o coletor da nuvem (\arq{observabilidade.tf}):
    \begin{itemize}
      \item \textbf{GCP:} \textbf{Cloud Logging} (com exclusão das sondas no sink
        \cd{\_Default} para economizar franquia), \textbf{Cloud Monitoring} (3
        \emph{uptime checks}, alertas e Service Monitoring SLO), \textbf{Cloud
        Error Reporting} e ingestão direta na API \textbf{OTLP} oficial
        (\cd{telemetry.googleapis.com}).
      \item \textbf{AWS:} \textbf{CloudWatch Logs}, \textbf{CloudWatch Metrics
        \& Alarms}, \textbf{CloudWatch Synthetics (Canaries)} para sondas de
        saúde (atenção: Canaries na AWS têm custo por execução significativamente
        maior que os \emph{uptime checks} do GCP) e \textbf{AWS X-Ray} (recebendo
        OTLP via \emph{sidecar} ADOT --- AWS Distro for OpenTelemetry --- na
        tarefa Fargate).
      \item \textbf{Azure:} \textbf{Azure Monitor} + \textbf{Log Analytics
        Workspace} + \textbf{Application Insights}. O Application Insights aceita
        OpenTelemetry nativamente, agrupa exceções equivalentes ao Error
        Reporting e oferece \emph{Standard Availability Tests} equivalentes aos
        \emph{uptime checks} a custo muito baixo.
    \end{itemize}
\end{campos}
\end{pecaarq}

% ===========================================================================
%  CAPITULO 4 — Eixo 2: Experiencias (Bundles Estaticos, CDN, Assets e Planos)
% ===========================================================================
\chapter{Eixo 2 --- Experiências (Bundles Estáticos, CDN e Armazenamento)}
\label{cap:eixo-experiencias}

Na Plataforma Unooria, uma regra arquitetural inviolável determina que
\textbf{experiências pedagógicas (\arq{apps/experiencias/*}) nunca viram
serviços de contêiner no Cloud Run}. Elas são aplicações cliente puras (Vite +
React + \cd{@edu/sdk}), compiladas para arquivos estáticos e publicadas em
armazenamento de objetos atrás de CDN sob caminhos versionados e imutáveis:
\cd{/\{expId\}/\{versao\}/}. Toda inteligência, correção de nota, estado de
sessão e chamada de IA acontece no servidor (\arq{apps/api}).

Além do bucket dos \emph{bundles} das experiências (\cd{-exp}), este eixo
abrange os três outros buckets de objetos que sustentam o funcionamento das
missões e das ferramentas do professor (\arq{storage.tf}): \cd{-assets} (mídia
pesada pública e imagens geradas por IA), \cd{-planos} (planos de aula em
Markdown e áudios DUA dos professores, privados e versionados) e
\cd{-exp-mapas} (\emph{sourcemaps} privados para desofuscação de pilhas de erro).

\begin{fichaeixo}{Ficha do Eixo 2 --- Experiências e Armazenamento de Objetos}
\begin{campos}
  \item[Arquivos de referência] \arq{storage.tf}, \arq{loadbalancer.tf},
    \arq{scripts/publicar-manifesto.mjs} e \arq{apps/experiencias/AGENTS.md}.
  \item[Aderência por nuvem] \sgcp~\nativo \quad \saws~\nativo \quad
    \sazure~\nativo
  \item[Ponto focal] Este é o eixo de \textbf{maior paridade e menor atrito de
    migração} entre as três nuvens. GCS + Cloud CDN, Amazon S3 + CloudFront e
    Azure Blob Storage + Azure Front Door oferecem garantias de durabilidade
    (11 noves), versionamento de objetos, políticas de ciclo de vida e controle
    de cache por cabeçalho \cd{Cache-Control} virtualmente idênticas.
\end{campos}
\end{fichaeixo}

\section{Bundles estáticos imutáveis, CDN e isolamento de origem (\texttt{-exp})}

O contrato entre o Hub e a experiência embutida no \cd{<iframe>} depende de duas
propriedades da infraestrutura de entrega estática:

\begin{enumerate}
  \item \textbf{Imutabilidade e \emph{rollback} instantâneo:} publicar uma
    experiência copia os arquivos com \emph{hash} para
    \cd{/\{expId\}/\{versao\}/assets/*} com \cd{Cache-Control: public,
    max-age=31536000, immutable}, e o \arq{index.html} + \arq{experiencia.json}
    com \cd{Cache-Control: public, max-age=60}. Ativar ou reverter uma versão
    não apaga nem move arquivos no bucket --- apenas aponta o manifesto para a
    versão desejada (\arq{publicar-manifesto.mjs}), fazendo o \emph{rollback}
    propagar em até 60 segundos sem invalidar o cache global da CDN.
  \item \textbf{Origem dedicada para o \cd{postMessage}:} o Hub envia o token
    JWT para o \cd{<iframe>} via \cd{postMessage} com \cd{targetOrigin}
    explícito (nunca \cd{*}). Se as experiências fossem servidas por
    \cd{https://storage.googleapis.com/\{bucket\}} (ou \cd{s3.amazonaws.com}),
    elas dividiriam a mesma origem web (\emph{origin}) com qualquer outro bucket
    público daquela nuvem. Por isso o Load Balancer/CDN expõe o bucket sob
    domínio próprio do contrato (\cd{exp\_domain}, ex.:
    \cd{https://exp-entreredes.unooria.com}).
\end{enumerate}

\begin{pecaarq}{EXP-01}{Entrega de Bundles Estáticos e Manifesto por CDN (\texttt{\{contrato\}-exp})}{\sgcp~\nativo\enspace\saws~\nativo\enspace\sazure~\nativo}
\begin{campos}
  \item[Como é hoje (GCP)] Bucket \textbf{Cloud Storage}
    (\cd{\{project\_id\}-exp}) com \emph{Uniform Bucket-Level Access}, conectado
    ao Global External Load Balancer como \cd{google\_compute\_backend\_bucket}
    com \textbf{Cloud CDN} ativado no modo \cd{USE\_ORIGIN\_HEADERS}
    (\arq{loadbalancer.tf}). A publicação no CI usa \cd{gcloud storage rsync}
    com cabeçalhos diferenciados para \arq{assets/} versus \arq{index.html} e
    \arq{experiencia.json}.
  \item[Equivalente AWS] Bucket \textbf{Amazon S3} privado (com \emph{S3 Block
    Public Access} ligado em 100\% do bucket) servido exclusivamente pelo
    \textbf{Amazon CloudFront} usando \textbf{Origin Access Control (OAC)}. É um
    pequeno ganho de segurança sobre o Backend Bucket clássico do GCP: na AWS, o
    bucket S3 nem precisa ter leitura pública (\cd{allUsers}), pois apenas a
    distribuição CloudFront assinada via SigV4 consegue ler os objetos. A
    Política de Cache do CloudFront (\cd{UseOriginCacheControlHeaders}) respeita
    exatamente o \cd{max-age=60} do \arq{index.html} e o \cd{immutable} dos
    \emph{chunks} JS/CSS, e o deploy usa \cd{aws s3 sync}.
  \item[Equivalente Azure] Conta de armazenamento \textbf{Azure Blob Storage}
    (contêiner \cd{exp}) atrás do \textbf{Azure Front Door} com CDN habilitado e
    respeito aos cabeçalhos \cd{Cache-Control} da origem. No Front Door
    \emph{Premium}, o contêiner Blob também pode permanecer 100\% privado via
    \emph{Private Link} entre a borda da Microsoft e o Storage Account; no
    \emph{Standard}, usa acesso de leitura de blobs anônimos restrito ao
    contêiner, igual ao GCS. O deploy usa \cd{az storage blob sync} ou \cd{azcopy
    sync}.
\end{campos}
\end{pecaarq}

\section{Mídia pesada e imagens geradas por IA (\texttt{-assets})}

Para manter o repositório Git enxuto (regra inviolável: nenhum binário $>2\,\text{MB}$
no Git), vídeos de abertura das missões (como os clipes MP4 da \cd{exp02em1b}),
modelos 3D (\arq{mascote-unooria.glb}), áudios longos de podcasts e todas as
imagens \arq{.webp} geradas pela IA no preparo do professor (deduplicadas por
SHA-256 da descrição) moram no bucket \cd{\{contrato\}-assets}.

\begin{pecaarq}{EXP-02}{Bucket Público de Assets e Escrita pela API (\texttt{\{contrato\}-assets})}{\sgcp~\nativo\enspace\saws~\nativo\enspace\sazure~\nativo}
\begin{campos}
  \item[Como é hoje (GCP)] Bucket \textbf{Cloud Storage}
    (\cd{\{project\_id\}-assets}) exposto sob \cd{assets\_domain} no mesmo Load
    Balancer + Cloud CDN, com política CORS permitindo \cd{GET}/\cd{HEAD} a
    partir das origens do Hub e das Experiências (necessário para carregamento
    de texturas WebGL/Three.js no mascote 3D e \cd{<canvas>}). A conta de
    serviço da API possui \cd{roles/storage.objectAdmin} apenas neste bucket e
    no de planos, gravando as imagens convertidas em \arq{.webp} via \cd{sharp}
    com \cd{Cache-Control: public, max-age=31536000, immutable}.
  \item[Equivalente AWS] Bucket \textbf{Amazon S3} (\cd{\{contrato\}-assets})
    atrás do CloudFront (\cd{assets\_domain}) com \emph{Response Headers Policy}
    para CORS (\cd{Access-Control-Allow-Origin} dinâmico por lista de origens
    do contrato). A \emph{Task Role} da API no ECS Fargate recebe permissão IAM
    \cd{s3:PutObject} e \cd{s3:HeadObject} restrita ao ARN do bucket, usando o
    \cd{@aws-sdk/client-s3} (ou a interface compatível S3 do próprio GCS, se
    desejado unificar código).
  \item[Equivalente Azure] Contêiner \cd{assets} no \textbf{Azure Blob Storage}
    atrás do Azure Front Door, com regras nativas de CORS no serviço de Blob e
    permissão RBAC \emph{Storage Blob Data Contributor} atribuída à
    \emph{Managed Identity} do Container App da API.
\end{campos}
\end{pecaarq}

\section{Guarda versionada de planos de aula e sourcemaps privados}

Os dois últimos buckets do Eixo 2 são estritamente privados (\cd{public\_read =
false}) e nunca são expostos no Load Balancer nem na CDN:

\begin{pecaarq}{EXP-03}{Buckets Privados: Planos Versionados (\texttt{-planos}) e Mapas (\texttt{-exp-mapas})}{\sgcp~\nativo\enspace\saws~\nativo\enspace\sazure~\nativo}
\begin{campos}
  \item[Planos e anexos DUA (\texttt{-planos})] Guarda os planos de aula anuais
    gerados pela IA e editados pelos professores, além dos áudios e materiais
    de Desenho Universal para a Aprendizagem (DUA). Possui \textbf{Object
    Versioning} ativo e regra de ciclo de vida que retém versões não atuais por
    \cd{planos\_historico\_dias} (180 dias em produção), garantindo que uma
    edição acidental de plano nunca perca o texto anterior. Só a API lê e
    escreve; o navegador nunca acessa este bucket diretamente.
  \item[Sourcemaps (\texttt{-exp-mapas})] Guarda os arquivos \arq{*.js.map}
    gerados no build das experiências (\cd{sourcemap: 'hidden'}), permitindo
    desofuscar um \emph{stack trace} capturado em produção sem expor o código
    original aos estudantes no bucket público \cd{-exp}.
  \item[Tradução AWS e Azure] Tanto o \textbf{Amazon S3} (\emph{S3 Versioning} +
    regra \cd{NoncurrentVersionExpiration} de 180 dias) quanto o \textbf{Azure
    Blob Storage} (\emph{Blob Versioning} + política de gerenciamento de ciclo
    de vida \cd{daysAfterVersionCreationGreaterThan: 180}) implementam esse
    comportamento de maneira declarativa no Terraform, com paridade funcional de
    100\%.
\end{campos}
\end{pecaarq}

\begin{table}[H]
\centering
\caption{Comparativo técnico dos 4 buckets do Eixo de Experiências nas três nuvens}
\label{tab:comparativo-buckets-exp}
\footnotesize
\begin{tabularx}{\linewidth}{@{}L{28mm}L{38mm}L{38mm}X@{}}
\toprule
\headrow \thd{Recurso / Requisito} & \thd{GCP (Cloud Storage)} & \thd{AWS (Amazon S3)} & \thd{Azure (Blob Storage)}\\
\midrule
SDK na API (\arq{apps/api}) &
  \cd{@google-cloud/storage} (ADC nativo) &
  \cd{@aws-sdk/client-s3} (IAM Task Role) &
  \cd{@azure/storage-blob} (Managed Identity)\\
\zebra Sincronização no CI &
  \cd{gcloud storage rsync} &
  \cd{aws s3 sync} &
  \cd{azcopy sync} / \cd{az storage blob sync}\\
Proteção da origem pública &
  Bucket público lido pelo Cloud CDN &
  Bucket privado lido via \textbf{CloudFront OAC} &
  Contêiner público ou privado via \textbf{Private Link}\\
\zebra Versionamento (\cd{-planos}) &
  \cd{versioning = true} + \cd{num\_newer\_versions} / \cd{days\_since\_noncurrent} &
  \cd{aws\_s3\_bucket\_versioning} + \cd{noncurrent\_version\_expiration} &
  \cd{versioning\_enabled = true} + \cd{ManagementPolicy} por idade da versão\\
Custo base (armazenamento Standard) &
  $\approx$ US\$\,0,020 / GB-mês (\cd{us-east1}) &
  $\approx$ US\$\,0,023 / GB-mês (\cd{us-east-1}) &
  $\approx$ US\$\,0,018 a 0,021 / GB-mês (\emph{Hot tier})\\
\bottomrule
\end{tabularx}
\end{table}

% =============================================================================
% CAPÍTULO 5 — Eixo 3: API (Backend Único, Jobs, Filas e Tempo Real)
% =============================================================================
\chapter{Eixo 3 --- API (\texttt{apps/api})}
\label{cap:eixo-api}

\begin{fichaeixo}{Ficha do Eixo 3 --- API (\texttt{apps/api}: Backend Único, Jobs, SSE e WebSockets)}
\begin{campos}
  \item[Serviços no código] \arq{apps/api/src/*} (17 módulos de domínio) e job de migração \arq{prisma migrate deploy}.
  \item[Arquivos Terraform] \arq{cloudrun.tf}, \arq{iam.tf}, \arq{network.tf} e \arq{variables.tf}.
  \item[Dimensionamento atual] 2 vCPU, 2 GiB RAM (\cd{sharp} + Prisma Client), \cd{min=1}, \cd{max=5}, concorrência 80, CPU sempre alocada (\cd{cpu\_idle = false}).
  \item[Aderência por nuvem] \sgcp~\nativo \quad \saws~\equivalente\ (ECS Fargate) / \incompativel\ (Lambda) \quad \sazure~\nativo\ (Container Apps)
\end{campos}
\end{fichaeixo}

Por decisão arquitetural deliberada (ADR-01 e regra \emph{``Não criar serviço Cloud Run novo sem necessidade demonstrada''}), a \arq{apps/api} concentra \textbf{todos} os papéis de backend da plataforma em um único artefato implantável: API REST síncrona, receptor de \emph{push} da fila de telemetria, proxy e orquestrador de chamadas de IA, processador de geração assíncrona em segundo plano (após responder \cd{HTTP 202}) e servidor de conexões persistentes (SSE e WebSockets bidirecionais de até 60 minutos).

Essa consolidação reduz drasticamente o custo operacional para uma equipe de uma única pessoa, mas impõe quatro requisitos técnicos severos ao serviço de hospedagem de contêineres da nuvem escolhida --- requisitos que eliminam imediatamente plataformas de funções efêmeras (\emph{Function-as-a-Service} como AWS Lambda ou Azure Functions).

% =============================================================================
\section{Por que a API não cabe em FaaS (Lambda / Functions)}
\label{sec:api-requisitos-computacao}

Antes de comparar os serviços de cada provedor, é essencial documentar os quatro comportamentos reais do código de \arq{apps/api/} que determinam o modelo de computação:

\begin{enumerate}
  \item \textbf{CPU sempre alocada fora do ciclo de requisição HTTP (\cd{cpu\_idle = false}):}
  Em dois fluxos críticos --- a geração de planos de aula com IA (\arq{apps/api/src/planos/}) e a geração de pacotes de missões inéditas (\arq{apps/api/src/geracao/}) --- a rota HTTP grava o registro inicial com status \cd{gerando}, devolve \textbf{\cd{HTTP 202 Accepted}} imediatamente ao navegador e dispara a execução pesada de forma assíncrona no mesmo processo Node.js (\cd{void this.processarEmBackground(...)}). Em ambientes onde a CPU é congelada no instante em que a resposta HTTP é enviada (comportamento padrão do AWS Lambda e do modo \emph{request-based} do Cloud Run), a geração simplesmente para no meio ou leva minutos para avançar poucos milissegundos de CPU.

  \item \textbf{Memória nativa para processamento de imagem (\cd{sharp}) e Prisma Client (2 GiB RAM):}
  Quando a IA gera ilustrações de capa ou de etapas pedagógicas, a API recebe os bytes brutos e executa conversão em memória para \cd{.webp} (qualidade 85) via \cd{sharp} (biblioteca C++ \emph{libvips}) antes de gravar no bucket \cd{-assets}, mantendo simultaneamente o motor Rust/N-API do Prisma Client e buffers de áudio PCM. O consumo em pico exige \textbf{2 vCPU e 2 GiB de RAM} por réplica.

  \item \textbf{Streaming SSE de longa duração (até 300 segundos) sem bufferização:}
  As rotas do Agente de IA do Professor (\cd{POST /agente-professor/conversas/:id/mensagens}) e o endpoint de LLM customizado consumido pelo HeyGen LiveAvatar (\cd{POST /ajuda/llm/chat/completions}) mantêm streams \cd{text/event-stream} abertos por vários minutos. No caso do LiveAvatar, se o primeiro chunk SSE demorar mais de 10 segundos (\emph{Time-To-First-Token}), o provedor de vídeo aborta a sessão --- razão pela qual a API emite uma \cd{FALA\_DE\_ESPERA} em 7 segundos caso o Gemini ainda esteja raciocinando. Se o proxy reverso da nuvem retiver chunks em buffer (comportamento de certos gateways serverless), a fala de espera não chega a tempo e o avatar cai.

  \item \textbf{WebSockets bidirecionais de até 60 minutos (\cd{timeout = 3600s}):}
  A API sustenta dois canais WebSocket de longa duração:
  \begin{itemize}
    \item \textbf{Mascote com Voz Ao Vivo (\cd{/ajuda/mascote/ao-vivo}, ADR-158):} o navegador transmite áudio PCM 16 kHz pelo WebSocket até a API, que mantém um segundo WebSocket persistente com o \textbf{Vertex Gemini Live} (\cd{BidiGenerateContent}) e devolve áudio PCM 24 kHz em tempo real, com orçamento de até 30 minutos contínuos.
    \item \textbf{Unoocraft Multiplayer (\arq{apps/api/src/unoocraft/}):} servidor Socket.IO embutido na API para sincronização de estado das salas multijogador da turma, exigindo timeout de conexão de até 3600 segundos e roteamento com afinidade de sessão (\emph{sticky sessions}) quando há mais de uma réplica ativa.
  \end{itemize}
\end{enumerate}

% =============================================================================
\section{Peças Arquiteturais do Eixo API}
\label{sec:pecas-api}

% -----------------------------------------------------------------------------
\begin{pecaarq}{API-01}{Computação do Backend (\texttt{apps/api}) --- CPU Alocada e Concorrência}{\sgcp~\nativo\enspace\saws~\equivalente\enspace\sazure~\nativo}
\begin{campos}
  \item[Configuração GCP] \cd{google\_cloud\_run\_v2\_service.api} (\arq{cloudrun.tf}): \cd{cpu = "2"}, \cd{memory = "2Gi"}, \cd{cpu\_idle = false}, \cd{min\_instance\_count = 1}, \cd{max\_instance\_count = 5}, \cd{max\_instance\_request\_concurrency = 80}, \cd{timeout = "900s"} (expansível até \cd{"3600s"}), \cd{startup\_cpu\_boost = true}.
  \item[Sondas de saúde] \cd{startup\_probe} e \cd{liveness\_probe} HTTP em \cd{GET /health} (porta 3001).
  \item[Como é hoje (GCP)] O Cloud Run v2 combina a simplicidade de um PaaS serverless (deploy de imagem em 15 segundos, revisões imutáveis, tráfego dividido para rollback imediato) com controle fino do ciclo de vida da CPU. Ao definir \cd{cpu\_idle = false}, as réplicas ativas mantêm as 2 vCPUs 100\% disponíveis tanto durante as 80 requisições simultâneas quanto nos intervalos entre elas, garantindo que as tarefas em segundo plano (\cd{HTTP 202}) terminem em segundos. O faturamento de 1 réplica mínima aquecida 24x7 com CPU sempre alocada fica em torno de US\$ 68/mês em \cd{southamerica-east1}, escalando até 5 réplicas apenas durante o pico escolar das 8h da manhã.
  \item[Alternativas AWS] Na AWS, existem três candidatos com comportamentos muito distintos para o perfil da \arq{apps/api}:
  \begin{itemize}
    \item \textbf{AWS Lambda (\incompativel):} Inviável sem reescrever a arquitetura da API. O limite máximo inegociável de execução é de 15 minutos (quebrando sessões de 30--60 minutos do Mascote Ao Vivo e do Unoocraft), o modelo congela o processo imediatamente após o retorno do \cd{HTTP 202} (travando a geração de planos de aula e missões), e WebSockets exigem o \textbf{API Gateway WebSocket API}, que não mantém conexão TCP direta com o processo Node.js (transforma cada frame WS em uma invocação HTTP separada, incompatível com o proxy stateful do Gemini Live e com o Socket.IO em memória).
    \item \textbf{AWS App Runner (\adaptacao):} Mais próximo da experiência do Cloud Run, mas apresenta limitações severas para a Unooria: o timeout máximo de requisição no ingress Envoy do App Runner historicamente trava em 120 segundos (insuficiente para streams SSE longos e conexões WebSocket de 60 minutos), não possui recurso nativo de \emph{one-off Job} para rodar migrações do banco (\cd{prisma migrate deploy}) e exige configurações extras para coexistir atrás de um mesmo ALB.
    \item \textbf{Amazon ECS on AWS Fargate (\equivalente\ --- Recomendado na AWS):} É a escolha tecnicamente sólida na AWS. Uma \emph{Task Definition} Fargate com \cd{2 vCPU} e \cd{4 GB RAM} (o Fargate exige no mínimo 4 GB para 2 vCPU) tem CPU sempre alocada por natureza, sustenta centenas de conexões concorrentes no mesmo processo Node.js e integra-se nativamente ao Application Load Balancer (ALB). O ponto de atenção é operacional: o escalonamento horizontal do ECS Fargate baseado em métricas do CloudWatch é significativamente mais lento que o do Cloud Run (leva de 90 a 180 segundos para provisionar e registrar uma nova Task no Target Group durante a rampa das 8h da manhã, contra 5 a 10 segundos no Cloud Run).
  \end{itemize}
  \item[Equivalente Azure] O \textbf{Azure Container Apps (ACA)} no perfil de carga \emph{Consumption} (ou \emph{Dedicated}) é o equivalente direto (\nativo) do Cloud Run. Enquanto uma réplica está ativa (\cd{minReplicas = 1}), a CPU permanece alocada ao contêiner sem congelamento pós-resposta HTTP, permitindo que as rotinas \cd{HTTP 202} rodem normalmente. O escalonamento via KEDA responde rapidamente a regras de concorrência HTTP (\cd{concurrentRequests = "80"}), e as sondas \cd{startupProbe} e \cd{livenessProbe} mapeiam 1:1 com a declaração atual.
\end{campos}
\end{pecaarq}

% -----------------------------------------------------------------------------
\begin{pecaarq}{API-02}{Conexões de Longa Duração: SSE, WebSockets (60 min) e Afinidade}{\sgcp~\nativo\enspace\saws~\equivalente\enspace\sazure~\nativo}
\begin{campos}
  \item[Rotas críticas] SSE: \cd{/agente-professor/*}, \cd{/ajuda/llm/chat/completions}. WebSockets: \cd{/ajuda/mascote/ao-vivo} (voz PCM bidirecional) e \cd{/socket.io/*} (Unoocraft).
  \item[Requisitos de rede] Desativação de bufferização de resposta no proxy, timeout de inatividade/conexão $\ge 3600\text{s}$ e \emph{session affinity} para Socket.IO.
  \item[Como é hoje (GCP)] O Cloud Load Balancing Global em conjunto com o Cloud Run suporta upgrade HTTP/1.1 \cd{Connection: Upgrade} (WebSockets) e streaming SSE sem bufferização intermediária (desde que o Cloud CDN esteja desabilitado no Backend Service da API, como já configurado em \arq{loadbalancer.tf}). O timeout do serviço Cloud Run permite até \textbf{3600 segundos} (60 minutos), e a afinidade de sessão garante que os pacotes de reconexão do Socket.IO atinjam a mesma instância.
  \item[Equivalente AWS] No par \textbf{ALB + ECS Fargate}, tanto SSE quanto WebSockets funcionam de maneira nativa e estável, mas exigem três ajustes obrigatórios no Terraform da AWS (sem os quais as conexões caem aos 60 segundos):
  \begin{enumerate}
    \item Elevar o \cd{idle\_timeout} do Application Load Balancer do padrão de \cd{60} segundos para \textbf{\cd{3600}} segundos (o ALB suporta até 4000 segundos).
    \item Habilitar \textbf{Stickiness} (afinidade por cookie gerenciado pelo ALB ou duração de sessão) no \emph{Target Group} da API para que o handshake do Socket.IO e o upgrade WebSocket não caiam em Tasks Fargate diferentes quando \cd{desired\_count > 1}.
    \item Configurar \cd{deregistration\_delay.timeout\_seconds} no Target Group para pelo menos \cd{300} segundos, evitando que o desligamento de uma Task corte abruptamente uma sessão ativa de voz com o mascote ou uma aula com o Agente do Professor.
  \end{enumerate}
  \item[Equivalente Azure] O ingress Envoy do \textbf{Azure Container Apps} suporta nativamente WebSockets e streaming SSE, além de oferecer afinidade de sessão embutida (\cd{stickySessions.affinity = "sticky"}). Quando exposto atrás do \textbf{Azure Application Gateway} (ou Front Door), ajusta-se o \emph{Request Timeout} para \cd{3600} segundos (no Front Door padrão, onde há limite de 240 segundos de inatividade por pacote, os \emph{heartbeats} WS e comentários SSE \cd{: keepalive} enviados a cada 30 segundos pela API mantêm o túnel aberto).
\end{campos}
\end{pecaarq}

% -----------------------------------------------------------------------------
\begin{pecaarq}{API-03}{Job de Migração de Banco (\texttt{prisma migrate deploy})}{\sgcp~\nativo\enspace\saws~\equivalente\enspace\sazure~\nativo}
\begin{campos}
  \item[Configuração GCP] \cd{google\_cloud\_run\_v2\_job.migrate} (\arq{cloudrun.tf}): mesma imagem da API, comando \cd{["pnpm", "--filter", "@edu/api", "prisma", "migrate", "deploy"]}, 1 vCPU / 1 GiB RAM, \cd{max\_retries = 0}, \cd{timeout = "600s"}.
  \item[Como é hoje (GCP)] Para nunca executar \cd{prisma migrate deploy} no \cd{entrypoint} das réplicas web (o que causaria corrida de locks no PostgreSQL quando múltiplas réplicas subissem juntas e atrasaria o \cd{startup\_probe}), a fundação separa a migração em um \textbf{Cloud Run Job} (\cd{\{contrato\}-migrate}). O pipeline de CI/CD atualiza a imagem do Job, dispara \cd{gcloud run jobs execute --wait} dentro da VPC e só promove o tráfego da API se o Job terminar com código \cd{0}.
  \item[Equivalente AWS] Utiliza-se o \textbf{Amazon ECS RunTask} (Standalone Task). Define-se uma \cd{aws\_ecs\_task\_definition} dedicada para migração (1 vCPU, 2 GB RAM, comando sobrescrito para \cd{prisma migrate deploy}, anexada às subnets privadas e ao Security Group do RDS). No deploy, a GitHub Action invoca \cd{aws ecs run-task} seguido de \cd{aws ecs wait tasks-stopped} e verifica o \cd{exitCode} do contêiner antes de atualizar o serviço principal no ECS.
  \item[Equivalente Azure] O \textbf{Azure Container Apps Jobs} (\cd{azurerm\_container\_app\_job} com \cd{trigger\_type = "Manual"}, \cd{replica\_timeout\_in\_seconds = 600} e \cd{replica\_retry\_limit = 0}) é idêntico ao Cloud Run Jobs em semântica, declaração Terraform e disparo via CLI (\cd{az containerapp job start --wait}).
\end{campos}
\end{pecaarq}

% -----------------------------------------------------------------------------
\begin{pecaarq}{API-04}{Processamento em Background (\texttt{HTTP 202}) vs. Filas Dedicadas}{\sgcp~\nativo\enspace\saws~\equivalente\enspace\sazure~\nativo}
\begin{campos}
  \item[Fluxos na API] Planos de aula (\cd{POST /planos/gerar}) e pacotes de missões (\cd{POST /geracao/pacotes}).
  \item[Mecanismo atual] Execução assíncrona no próprio processo com \cd{cpu\_idle = false} + recuperação automática de registros órfãos em \cd{gerando} há mais de 10 minutos.
  \item[Análise nos 3 provedores] Tanto no Cloud Run (\cd{cpu\_idle = false}) quanto no ECS Fargate e no Azure Container Apps, o padrão atual --- disparar a promessa assíncrona no processo da API após responder \cd{HTTP 202} e consultar o status por polling no Hub --- funciona sem necessidade de criar um serviço worker separado. Caso um deploy ou desligamento de réplica interrompa uma geração em andamento, a própria API detecta registros presos no estado \cd{gerando} além do limite de tempo (10 minutos) e permite a retentativa limpa pelo professor. No estágio atual de 10 mil estudantes por contrato, manter tudo no mesmo serviço da API permanece a decisão correta em qualquer uma das três nuvens.
\end{campos}
\end{pecaarq}

% =============================================================================
\section{Síntese Comparativa do Eixo API}

\begin{table}[H]
\centering
\begin{tabularx}{\textwidth}{@{}L{34mm}YYY@{}}
\toprule
\headrow
\thd{Requisito da API} &
\thd{GCP (Cloud Run v2)} &
\thd{AWS (ECS Fargate + ALB)} &
\thd{Azure (Container Apps)} \\
\midrule
\textbf{CPU ativa pós-\cd{HTTP 202}} &
Sim (\cd{cpu\_idle = false}) &
Sim (nativo da Task Fargate) &
Sim (réplica ativa mantém CPU) \\
\zebra
\textbf{Dimensionamento unitário} &
2 vCPU / 2 GiB RAM &
2 vCPU / 4 GB RAM (mínimo Fargate p/ 2 vCPU) &
2 vCPU / 4 GiB RAM (perfil Consumption) \\
\textbf{Concorrência por réplica} &
Até 80 req/instância (configurável até 1000) &
Ilimitada pelo Fargate (governada pelo Node.js e CPU) &
Regra KEDA HTTP (\cd{concurrentRequests = 80}) \\
\zebra
\textbf{Velocidade de scale-out} &
\textbf{5 a 15 segundos} (revisão pré-aquecida no pool Google) &
\textbf{90 a 180 segundos} (pull da imagem + ENI + health check ALB) &
\textbf{15 a 35 segundos} (provisionamento de pod no ACA) \\
\textbf{Timeout máximo (WS / SSE)} &
\textbf{3.600 s} (60 min) nativo no serviço &
\textbf{4.000 s} no ALB (\cd{idle\_timeout} ajustado) &
\textbf{3.600 s+} no ACA / App Gateway (heartbeat p/ Front Door) \\
\zebra
\textbf{Migração (\cd{prisma migrate})} &
Cloud Run Job (\cd{--wait}) &
ECS Standalone Task (\cd{run-task} + \cd{wait}) &
Container Apps Job (\cd{job start --wait}) \\
\bottomrule
\end{tabularx}
\caption{Comparativo técnico de execução do backend monolítico (\texttt{apps/api}) nas três nuvens.}
\label{tab:comparativo-eixo-api}
\end{table}

% =============================================================================
% CAPÍTULO 6 — Eixo 4: Uso de IA (Vertex AI, Agente do Professor, RAG e Voz)
% =============================================================================
\chapter{Eixo 4 --- Uso de IA (Modelos, Agente, RAG e Voz Ao Vivo)}
\label{cap:eixo-ia}

\begin{fichaeixo}{Ficha do Eixo 4 --- Uso de IA (Modelos, Agente do Professor, RAG e Voz Ao Vivo)}
\begin{campos}
  \item[Serviços no código] \arq{apps/api/src/ai/}, \arq{apps/agente-professor/} (Python + Google ADK), \arq{apps/api/src/agente-professor/}, \arq{apps/api/src/ajuda/}, \arq{apps/api/src/planos/} e \arq{apps/api/src/geracao/}.
  \item[Arquivos Terraform] \arq{agente\_professor.tf}, \arq{cloudrun.tf}, \arq{iam.tf} e \arq{apis.tf}.
  \item[Contratos e ADRs] ADR-105 (quotas no banco), ADR-126/127 (Agent Engine + Vertex AI Search), ADR-132 (LiveAvatar), ADR-158 (Mascote Gemini Live) e ADR-161 (painel de custos de IA).
  \item[Aderência por nuvem] \sgcp~\nativo \quad \saws~\adaptacao\ (alto esforço e custo fixo de RAG) \quad \sazure~\adaptacao\ (reescrita do agente + custo fixo de Search)
\end{campos}
\end{fichaeixo}

Se nos eixos de Plataforma, Experiências e API a migração entre GCP, AWS e Azure é essencialmente um exercício de infraestrutura (\emph{Terraform}, rede e IAM) sem tocar no código de aplicação, \textbf{o Eixo 4 (Uso de IA) é onde reside a assimetria real entre as três nuvens}.

A Plataforma Unooria não consome inteligência artificial apenas como um ``chat genérico de texto''. Ela opera quatro subsistemas distintos de IA, cada um explorando capacidades específicas do ecossistema \textbf{Google Vertex AI}:
\begin{enumerate}
  \item \textbf{Camada Pluggable da API (\arq{apps/api/src/ai/}):} geração de texto pedagógico e correção estruturada com controle de orçamento de raciocínio (\cd{thinkingBudget}) e até \cd{65.536} tokens de saída (planos de aula anuais), geração de ilustrações nativas na proporção \cd{4:5} e síntese de áudio multi-locutor (\cd{multiSpeakerMarkup}) em uma única chamada.
  \item \textbf{Agente de IA do Professor (\arq{apps/agente-professor/}):} agente autônomo escrito em Python sobre o \textbf{Google Agent Development Kit (ADK)} e hospedado no \textbf{Vertex AI Agent Engine} (\cd{ReasoningEngine}), que gerencia sessões conversacionais nativas sem banco de dados auxiliar.
  \item \textbf{RAG Pedagógico Duplo (\textbf{Vertex AI Search / Discovery Engine}):} duas bases de conhecimento com semânticas opostas --- uma base estruturada sem conteúdo (\cd{NO\_CONTENT}) para filtro exato + paginação integral das habilidades da BNCC e currículos locais, e uma base não-estruturada (\cd{CONTENT\_REQUIRED}) com \emph{Layout Parser} de PDFs que devolve trechos com numeração exata de página inicial e final (\cd{pageSpan: pageStart..pageEnd}).
  \item \textbf{Ajuda Multimodal em Tempo Real (\arq{apps/api/src/ajuda/}):} Mascote 3D com conversação de voz bidirecional nativa (\emph{speech-to-speech} via WebSocket no \textbf{Vertex Gemini Live}) e Avatar em Vídeo (\textbf{HeyGen LiveAvatar} em modo FULL conectado ao LLM da API com \emph{Implicit Prompt Caching} de uma base de conhecimento de $\sim 65$ mil tokens).
\end{enumerate}

% =============================================================================
\section{Topologia Comparada da Stack de IA}
\label{sec:topologia-ia}

A Figura~\ref{fig:topologia-ia-comparada} contrasta a arquitetura atual de IA na GCP com as composições equivalentes na AWS e na Azure.

\begin{figure}[H]
\centering
\begin{tikzpicture}[
  font=\sffamily\scriptsize,
  node distance=0.35cm and 0.4cm,
  colbox/.style={
    draw=unooGrid, rounded corners=4pt, fill=unooPaper,
    minimum width=5.1cm, minimum height=10.6cm, anchor=north
  },
  colhead/.style={
    rounded corners=3pt, fill=unooNavy, text=white, font=\sffamily\bfseries\footnotesize,
    minimum width=4.9cm, minimum height=0.65cm, align=center
  },
  svc/.style={
    draw=unooGrid, rounded corners=3pt, fill=unooSurface, text=unooInk,
    minimum width=4.6cm, minimum height=1.15cm, align=center, inner sep=3pt
  },
  svcgcp/.style={
    draw=unooBlue, line width=0.8pt, rounded corners=3pt, fill=unooBlueWash, text=unooInk,
    minimum width=4.6cm, minimum height=1.15cm, align=center, inner sep=3pt
  },
  svcwarn/.style={
    draw=unooOrange, line width=0.8pt, rounded corners=3pt, fill=unooOrangeWash, text=unooInk,
    minimum width=4.6cm, minimum height=1.15cm, align=center, inner sep=3pt
  }
]

% Coluna GCP
\node[colbox] (gcpbox) at (0,0) {};
\node[colhead, fill=unooBlue] at (0,-0.15) {GCP (Arquitetura Atual)};

\node[svcgcp] (gcp1) at (0,-1.25) {
  \textbf{1. Modelos na API (\texttt{src/ai})}\\
  Vertex AI \texttt{gemini-3.6/3.7-flash}\\
  Imagem \texttt{4:5}: \texttt{gemini-3.1-flash-image}\\
  TTS Multi-voz: \texttt{gemini-2.5-flash-tts}
};

\node[svcgcp] (gcp2) at (0,-2.85) {
  \textbf{2. Runtime do Agente + Sessão}\\
  \textbf{Vertex AI Agent Engine} (\texttt{BYOC})\\
  Container Python + \textbf{Google ADK}\\
  Sessões persistidas pelo próprio serviço
};

\node[svcgcp] (gcp3) at (0,-4.65) {
  \textbf{3. RAG BNCC + Acervo PDF}\\
  \textbf{Vertex AI Search} (Discovery Engine)\\
  Store 1: \texttt{NO\_CONTENT} (filtro + pag.)\\
  Store 2: \texttt{Layout Parser} (\texttt{pageSpan})\\
  \textbf{Custo ocioso: $\sim$US\$ 0/mês}
};

\node[svcgcp] (gcp4) at (0,-6.45) {
  \textbf{4. Voz Ao Vivo (Mascote 3D)}\\
  \textbf{Vertex Gemini Live} (\texttt{us-central1})\\
  \texttt{gemini-3.8-live} (WSS Bidi PCM 24kHz)\\
  Cache implícito (prefixo 65k tokens)
};

\node[svc] (gcp5) at (0,-8.15) {
  \textbf{5. Governança e Custos}\\
  \texttt{MedidorDeProvider} + \texttt{QuotaService}\\
  Tarifas por modelo em \texttt{ia\_quotas}\\
  Autenticação 100\% via ADC (IAM)
};

\node[svc] (gcp6) at (0,-9.65) {
  \textbf{6. Avatar em Vídeo (Externo)}\\
  HeyGen LiveAvatar (Modo FULL)\\
  Callback SSE \seta{} \texttt{/ajuda/llm/*}
};

% Coluna AWS
\node[colbox] (awsbox) at (5.5,0) {};
\node[colhead, fill=unooOrange!85!black] at (5.5,-0.15) {AWS (Equivalência Proposta)};

\node[svc] (aws1) at (5.5,-1.25) {
  \textbf{1. Modelos na API (\texttt{src/ai})}\\
  Bedrock: \texttt{Claude 3.7/4 Sonnet} / \texttt{Nova Pro}\\
  Imagem: \texttt{Nova Canvas} (\texttt{1024x1280})\\
  TTS: \texttt{Amazon Polly} (concatenar vozes)
};

\node[svcwarn] (aws2) at (5.5,-2.85) {
  \textbf{2. Runtime do Agente + Sessão}\\
  \textbf{Bedrock AgentCore} ou \textbf{ECS Fargate}\\
  Requer portar ADK \seta{} Strands/LangGraph\\
  + \textbf{DynamoDB/RDS} p/ sessões se Fargate
};

\node[svcwarn] (aws3) at (5.5,-4.65) {
  \textbf{3. RAG BNCC + Acervo PDF}\\
  \textbf{Bedrock Knowledge Bases}\\
  Exige banco vetorial subjacente:\\
  OpenSearch Serverless (\textbf{US\$ 350+/mês})\\
  ou \textbf{RDS pgvector} + parser PDF custom
};

\node[svc] (aws4) at (5.5,-6.45) {
  \textbf{4. Voz Ao Vivo (Mascote 3D)}\\
  \textbf{Amazon Nova Sonic} (Bedrock)\\
  Streaming bidirecional HTTP/2\\
  Cache explícito via \texttt{cachePoint}
};

\node[svc] (aws5) at (5.5,-8.15) {
  \textbf{5. Governança e Custos}\\
  Mesmo \texttt{MedidorDeProvider} na API\\
  Novo adapter \texttt{BedrockAiProvider}\\
  Autenticação via IAM Task Role
};

\node[svc] (aws6) at (5.5,-9.65) {
  \textbf{6. Avatar em Vídeo (Externo)}\\
  HeyGen LiveAvatar (Modo FULL)\\
  Callback SSE \seta{} \texttt{/ajuda/llm/*}
};

% Coluna Azure
\node[colbox] (azbox) at (11.0,0) {};
\node[colhead, fill=unooNavy] at (11.0,-0.15) {Azure (Equivalência Proposta)};

\node[svc] (az1) at (11.0,-1.25) {
  \textbf{1. Modelos na API (\texttt{src/ai})}\\
  Azure OpenAI: \texttt{GPT-4.1} / \texttt{GPT-4o}\\
  Imagem: \texttt{gpt-image-1} (+ crop \texttt{4:5})\\
  TTS: \texttt{Azure AI Speech} (SSML multi-voz)
};

\node[svcwarn] (az2) at (11.0,-2.85) {
  \textbf{2. Runtime do Agente + Sessão}\\
  \textbf{Azure AI Agent Service} ou \textbf{ACA}\\
  Requer portar Google ADK \seta{} Semantic\\
  Kernel / LangGraph + \textbf{Cosmos DB/PG}
};

\node[svcwarn] (az3) at (11.0,-4.65) {
  \textbf{3. RAG BNCC + Acervo PDF}\\
  \textbf{Azure AI Search} + \textbf{Doc Intelligence}\\
  Excelente busca híbrida e filtro OData,\\
  mas cobra instância fixa (\textbf{US\$ 75--250/mês})
};

\node[svc] (az4) at (11.0,-6.45) {
  \textbf{4. Voz Ao Vivo (Mascote 3D)}\\
  \textbf{GPT-4o Realtime API}\\
  WebSocket / WebRTC áudio nativo\\
  Cache automático de prefixo ($>1024$ tok)
};

\node[svc] (az5) at (11.0,-8.15) {
  \textbf{5. Governança e Custos}\\
  Mesmo \texttt{MedidorDeProvider} na API\\
  Novo adapter \texttt{AzureOpenAiProvider}\\
  Autenticação via Managed Identity
};

\node[svc] (az6) at (11.0,-9.65) {
  \textbf{6. Avatar em Vídeo (Externo)}\\
  HeyGen LiveAvatar (Modo FULL)\\
  Callback SSE \seta{} \texttt{/ajuda/llm/*}
};

\end{tikzpicture}
\caption[Comparação da pilha de Inteligência Artificial]{Comparação camada a camada dos subsistemas de Inteligência Artificial entre GCP, AWS e Azure. Os blocos destacados em laranja indicam pontos onde há necessidade de refatoração de código ou degrau relevante de custo fixo mensal.}
\label{fig:topologia-ia-comparada}
\end{figure}

% =============================================================================
\section{Peças Arquiteturais do Eixo Uso de IA}
\label{sec:pecas-ia}

% -----------------------------------------------------------------------------
\begin{pecaarq}{IA-01}{Camada Plugável de Modelos na API (\texttt{AiProvider}): Texto, Imagem 4:5 e TTS}{\sgcp~\nativo\enspace\saws~\adaptacao\enspace\sazure~\equivalente}
\begin{campos}
  \item[Código na API] \arq{apps/api/src/ai/ai.provider.ts}, \arq{vertex-gemini.provider.ts}, \arq{medidor.provider.ts} e \arq{quota.service.ts}.
  \item[Modelos atuais (GCP)] \cd{gemini-3.6-flash} (padrão), \cd{gemini-3.7-flash} (planos de aula, até 65.536 tokens de saída), \cd{gemini-3.1-flash-image} (imagens \cd{4:5}) e \cd{gemini-2.5-flash-tts} (áudio MP3 com \cd{multiSpeakerMarkup}).
  \item[Como é hoje (GCP)] Na API, todo consumo de IA passa pela interface abstrata \cd{AiProvider}, decorada pelo \cd{MedidorDeProvider} (que intercepta o \cd{usageMetadata} de cada chamada e grava o consumo exato de tokens, imagens ou caracteres no PostgreSQL para o painel \cd{/unooria/uso-de-ia}, conforme ADR-161). O \cd{VertexGeminiProvider} explora quatro recursos importantes do Vertex AI:
  \begin{enumerate}
    \item \textbf{Orçamento de raciocínio controlável (\cd{thinkingBudget}):} chamadas que exigem resposta rápida (como correções curtas ou interações síncronas de estudantes) enviam \cd{thinkingBudget: 0} (desligando o pensamento oculto para derrubar a latência e o custo de tokens de saída), enquanto a geração de planos de aula e pacotes pedagógicos habilita o raciocínio profundo e janela de saída de até \textbf{65.536 tokens} em uma única geração.
    \item \textbf{Cache implícito de contexto (\emph{Implicit Caching}):} o endpoint da Ajuda (\cd{/ajuda/llm/chat/completions}) injeta a base completa de conhecimento da plataforma ($\sim 65.000$ tokens estáticos) no \cd{systemInstruction}. No Vertex Gemini, o prefixo repetido atinge automaticamente 75\% a 90\% de desconto (\cd{cached\_Tokens}) sem precisar criar ou gerenciar ciclo de vida de objetos de cache.
    \item \textbf{Proporção nativa \cd{4:5} na geração de imagem:} o \cd{gemini-3.1-flash-image} aceita \cd{imageConfig: \{ aspectRatio: '4:5' \}}, entregando a imagem exatamente na proporção exigida pelos cards e etapas do Design System antes da conversão para \cd{.webp}.
    \item \textbf{Diálogo multi-locutor nativo em uma única chamada TTS:} na geração de áudios de missões (ex.: diálogos entre personagens ou postagens narradas na Exp08), o \textbf{Cloud Text-to-Speech} (\cd{gemini-2.5-flash-tts}) recebe um array de turnos (\cd{multiSpeakerMarkup}) com vozes distintas (ex.: \cd{Aoede} e \cd{Puck}) e devolve um único arquivo MP3 já mixado e com prosódia conversacional entre os interlocutores.
  \end{enumerate}
  \item[Equivalente AWS] Como a arquitetura da API já isola o vendor atrás do token de injeção \cd{AI\_PROVIDER}, migrar para a AWS exige apenas escrever um \cd{BedrockAiProvider} sem alterar nenhum controller de negócio. Contudo, há diferenças funcionais nos modelos:
  \begin{itemize}
    \item \textbf{Texto e JSON Estruturado:} a API \emph{Bedrock Converse} oferece suporte a \textbf{Anthropic Claude 3.7 / 4 Sonnet} (com \emph{extended thinking} configurável e saída de até 64k tokens) e à família \textbf{Amazon Nova Pro / Lite}. Para aproveitar o desconto de cache nos 65k tokens da Ajuda, o adapter precisará inserir blocos explícitos \cd{cachePoint} no payload do Bedrock.
    \item \textbf{Geração de Imagens (\cd{4:5}):} os modelos \textbf{Amazon Nova Canvas} e \textbf{Titan Image Generator v2} aceitam resoluções customizadas (ex.: \cd{1024x1280}, que corresponde exatamente à razão \cd{4:5}), permitindo manter a compatibilidade visual.
    \item \textbf{Síntese de Voz Multi-Locutor (\adaptacao):} o \textbf{Amazon Polly} \textbf{não suporta múltiplos locutores em uma única chamada SSML} da mesma forma que o \cd{multiSpeakerMarkup} do Gemini TTS. Para sintetizar um diálogo entre dois personagens no Polly, o \cd{BedrockAiProvider} precisará quebrar o roteiro por turno, disparar $N$ chamadas paralelas ao Polly (uma por fala/voz) e concatenar os buffers de áudio na própria API antes de enviar ao S3.
  \end{itemize}
  \item[Equivalente Azure] No ecossistema Microsoft:
  \begin{itemize}
    \item \textbf{Texto e Cache:} o \textbf{Azure OpenAI Service} (\cd{gpt-4.1}, \cd{gpt-4.1-mini} e série \cd{o3-mini}/\cd{o4-mini}) possui \emph{Prompt Caching} 100\% automático para prefixos acima de 1.024 tokens (mesma ergonomia zero-configuração do Vertex Gemini). Limite de saída: \cd{gpt-4.1} suporta até 32.768 tokens de saída (metade dos 65.536 do Gemini 3.7 Flash).
    \item \textbf{Geração de Imagens (\cd{4:5}):} o modelo \textbf{gpt-image-1} (e \cd{DALL-E 3}) na Azure trabalha apenas com três resoluções fixas (\cd{1024x1024}, \cd{1024x1536} retrato 2:3 e \cd{1536x1024} paisagem 3:2). O adapter \cd{AzureOpenAiProvider} solicita \cd{1024x1536} e usa a própria biblioteca \cd{sharp} (já presente na API) para aplicar um \emph{center crop} automático para \cd{1024x1280} (\cd{4:5}) junto com a conversão para \cd{.webp}.
    \item \textbf{Síntese de Voz Multi-Locutor (\nativo):} o \textbf{Azure AI Speech} suporta múltiplas tags \cd{<voice name="pt-BR-FranciscaNeural">...</voice><voice name="pt-BR-AntonioNeural">...</voice>} intercaladas em um único documento SSML e devolve um único stream MP3 pronto.
  \end{itemize}
\end{campos}
\end{pecaarq}

% -----------------------------------------------------------------------------
\begin{pecaarq}{IA-02}{Agente de IA do Professor: Runtime (\texttt{apps/agente-professor}) e Sessões}{\sgcp~\nativo\enspace\saws~\adaptacao\enspace\sazure~\adaptacao}
\begin{campos}
  \item[Código e stack] \arq{apps/agente-professor/agente\_professor/} (Python 3.12 + \cd{google-adk}) + cliente HTTP na API (\arq{apps/api/src/agente-professor/agent-engine.client.ts}).
  \item[Configuração GCP] \cd{google\_vertex\_ai\_reasoning\_engine.agente\_professor} (\arq{agente\_professor.tf}): modo \cd{BYOC} (imagem Docker no Artifact Registry), \cd{4 vCPU}, \cd{4 GiB RAM}, \cd{min\_instances = 1}, \cd{max\_instances = 5}, \cd{concurrency = 36}.
  \item[Como é hoje (GCP)] O Agente do Professor foi construído em Python utilizando o \textbf{Google Agent Development Kit (ADK)} (\cd{LlmAgent}, ferramentas Python tipadas e callbacks de telemetria) e implantado no \textbf{Vertex AI Agent Engine} (\cd{ReasoningEngine}) em \cd{us-central1} (ADR-126). Essa escolha resolve dois problemas em um único recurso Terraform:
  \begin{enumerate}
    \item \textbf{Hospedagem gerenciada do contêiner Python (\cd{BYOC}):} o Agent Engine executa a imagem do Artifact Registry sem que a equipe precise gerenciar cluster Kubernetes ou balanceador interno. A API NestJS chama os endpoints REST nativos do Vertex AI (\cd{:query} e \cd{:streamQuery}) autenticada via ADC.
    \item \textbf{Estado e memória de sessão embutidos (\emph{Zero-DB Session Store}):} o histórico de turnos e eventos da conversa do professor com o agente \textbf{não é reenviado pela API a cada mensagem}. O Agent Engine persiste nativamente as sessões (\cd{async\_create\_session}, \cd{async\_get\_session}, \cd{async\_list\_sessions}) e permite renovar o TTL de expiração via \cd{PATCH} no recurso da sessão (\cd{expireTime = +365d} a cada nova mensagem). No PostgreSQL da plataforma (\cd{agente\_professor\_conversas}), a API guarda apenas o ponteiro \cd{sessao\_remota\_id} e o título.
  \end{enumerate}
  \item[Equivalente AWS] Este é o ponto de \textbf{maior acoplamento de código (\emph{lock-in})} de todo o monorepo:
  \begin{itemize}
    \item \textbf{Framework do Agente:} o servidor de produção do Google ADK e seu gerenciador de sessões (\cd{VertexAiSessionService}) são acoplados ao protocolo do Vertex AI Agent Engine. Na AWS, o caminho natural é migrar a definição do agente (cerca de 400 linhas de Python em \arq{apps/agente-professor/agente\_professor/}) para o \textbf{Strands Agents SDK} (framework oficial open-source da AWS) ou \textbf{LangGraph}.
    \item \textbf{Hospedagem e Sessões:} se hospedado no \textbf{Amazon Bedrock AgentCore Runtime} (para agentes customizados em contêiner com memória de sessão gerenciada), preserva-se a arquitetura sem banco adicional, mas reescreve-se todo o cliente \arq{agent-engine.client.ts} da API para assinar requisições AWS \emph{SigV4}. Alternativamente, roda-se o contêiner Python como um serviço privado no \textbf{ECS Fargate} persistindo os turnos de sessão no próprio \textbf{RDS PostgreSQL}.
  \end{itemize}
  \item[Equivalente Azure] Na Azure, existem duas abordagens equivalentes:
  \begin{itemize}
    \item \textbf{Abordagem 1 --- Azure Container Apps interno + PostgreSQL (Mais simples):} empacotar \arq{apps/agente-professor/} como uma aplicação FastAPI + LangGraph/Semantic Kernel rodando em um \textbf{Container App com ingress interno} (visível apenas para a API na mesma VNet) e apontando o \emph{checkpointer} de sessão para o próprio \textbf{Azure Database for PostgreSQL} já existente no contrato.
    \item \textbf{Abordagem 2 --- Azure AI Foundry Agent Service:} serviço gerenciado de agentes da Azure que persiste \emph{Threads}, \emph{Runs} e \emph{Messages} automaticamente (análogo às sessões do Vertex Agent Engine), mas exige remodelar o registro das ferramentas Python.
  \end{itemize}
\end{campos}
\end{pecaarq}

% -----------------------------------------------------------------------------
\begin{pecaarq}{IA-03}{RAG Pedagógico Duplo: Habilidades (BNCC) e Acervo PDF com Páginas}{\sgcp~\nativo\enspace\saws~\adaptacao\enspace\sazure~\adaptacao}
\begin{campos}
  \item[Código de busca] \arq{apps/agente-professor/agente\_professor/ferramentas.py} e \arq{ingestao/} (\cd{ingerir\_habilidades.py}, \cd{ingerir\_acervo.py}, \cd{esquema\_habilidades.json}).
  \item[Configuração GCP] 2 Data Stores no \textbf{Vertex AI Search} (\cd{google\_discovery\_engine\_data\_store} em \arq{agente\_professor.tf}): \cd{habilidades} (\cd{NO\_CONTENT}, schema estrito) e \cd{acervo} (\cd{CONTENT\_REQUIRED}, \cd{layout\_parsing\_config} + \cd{chunk\_size = 500} com \cd{include\_ancestor\_headings = true}).
  \item[Como é hoje (GCP)] O Agente do Professor possui duas ferramentas de consulta com comportamentos que raramente coexistem bem em um banco vetorial simples (ADR-127):
  \begin{enumerate}
    \item \textbf{\cd{consultar\_habilidades} (Data Store \cd{habilidades} --- \cd{NO\_CONTENT}):} armazena dezenas de milhares de habilidades da BNCC e dos currículos locais com esquema rígido (\cd{dynamic = "false"}): \cd{codigo} (\cd{EF06MA01}), \cd{anos} (\cd{["6º ano"]}), \cd{componente\_slug}, \cd{etapa}, \cd{unidade\_tematica}, \cd{objeto\_conhecimento} e \cd{descricao}. Quando o professor pede \emph{``liste todas as habilidades de Matemática do 6º ano''}, a ferramenta aplica filtro booleano exato (\cd{anos: ANY("6º ano") AND componente\_slug: ANY("matematica")}) com \cd{query=""} e \textbf{pagina até o fim} (blocos de 100 documentos) para devolver 100\% da matriz curricular sem omitir nenhuma habilidade. Quando busca por tema, combina o mesmo filtro exato com busca híbrida.
    \item \textbf{\cd{consultar\_acervo} (Data Store \cd{acervo} --- \cd{CONTENT\_REQUIRED}):} ingere livros didáticos e diretrizes em PDF. O \textbf{Layout Parser} nativo do Discovery Engine analisa a estrutura visual do PDF, divide em blocos de 500 tokens preservando os títulos ancestrais (\cd{include\_ancestor\_headings = true}) e anota em cada chunk o intervalo exato de páginas (\cd{chunkMetadata.pageSpan.pageStart} e \cd{pageEnd}). Assim, o agente cita o título da obra e as páginas exatas.
  \end{enumerate}
  Além da precisão técnica, o modelo de cobrança do \textbf{Vertex AI Search} é \textbf{100\% serverless por consumo} (poucos dólares por mil consultas e por GB indexado, com \textbf{zero custo fixo de cluster ligado}).
  \item[Equivalente AWS] O \textbf{Amazon Bedrock Knowledge Bases} gerencia a ingestão de PDFs a partir do S3, mas exige um banco vetorial subjacente provisionado:
  \begin{itemize}
    \item \textbf{Opção A --- Amazon OpenSearch Serverless (Padrão AWS, mas caro):} cada coleção exige um mínimo faturado de \textbf{2 a 4 OCUs} funcionando 24x7, gerando um piso fixo de \textbf{US\$ 175 a US\$ 350/mês por contrato} apenas para manter o índice no ar.
    \item \textbf{Opção B --- RDS PostgreSQL (\cd{pgvector}) + Full-Text Search (Recomendada na AWS):} aproveita a instância \textbf{RDS PostgreSQL 16} já existente no contrato (custo adicional zero de infraestrutura). A ferramenta \cd{consultar\_habilidades} torna-se uma query SQL híbrida direta no PostgreSQL (filtros exatos determinísticos + \cd{tsvector} + distância cosseno com embeddings \cd{Titan Text Embeddings v2}), e o \cd{ingerir\_acervo.py} extrai os blocos paginados via \textbf{Amazon Textract} ou \cd{Docling}/\cd{PyMuPDF} gravando \cd{page\_start} e \cd{page\_end} na tabela de chunks.
  \end{itemize}
  \item[Equivalente Azure] O \textbf{Azure AI Search} combinado ao \textbf{Azure AI Document Intelligence} (modelo \cd{prebuilt-layout}) entrega a mesma qualidade técnica do Discovery Engine (filtros OData exatos, paginação determinística, busca híbrida com \emph{Semantic Ranker} e chunking por layout com número de página). Contudo, cobra por instância provisionada: o tier \textbf{Basic} custa cerca de \textbf{US\$ 75/mês fixos} (e o \textbf{Standard S1} custa \textbf{US\$ 250/mês}). Para eliminar esse custo fixo por contrato na Azure, aplica-se a mesma estratégia de usar \textbf{\cd{pgvector}} no Azure Database for PostgreSQL.
\end{campos}
\end{pecaarq}

% -----------------------------------------------------------------------------
\begin{pecaarq}{IA-04}{Ajuda em Tempo Real: Mascote com Voz Ao Vivo e Avatar em Vídeo}{\sgcp~\nativo\enspace\saws~\equivalente\enspace\sazure~\nativo}
\begin{campos}
  \item[Código na API] \arq{apps/api/src/ajuda/} (\cd{mascote-live.gateway.ts}, \cd{ajuda-llm.controller.ts}, \cd{liveavatar.service.ts}, \cd{ajuda.service.ts}).
  \item[Serviços atuais] \textbf{Vertex Gemini Live} (\cd{gemini-3.8-live} via WSS em \cd{us-central1}) para o Mascote 3D e \textbf{HeyGen LiveAvatar} (SaaS externo) + \cd{gemini-3.6-flash} (com cache de 65k tokens) para o Avatar em Vídeo.
  \item[Mascote 3D (Voz)] No modo Mascote Ao Vivo (\cd{VOZ\_AO\_VIVO}, ADR-158), a API abre um túnel WebSocket autenticando na Vertex AI (\cd{BidiGenerateContent}) com o modelo \textbf{\cd{gemini-3.8-live}} (áudio nativo de entrada PCM 16 kHz e saída PCM 24 kHz + transcrição em tempo real e detecção de interrupção/\emph{barge-in}):
  \begin{itemize}
    \item \textbf{Na AWS (\equivalente):} o equivalente nativo é o \textbf{Amazon Nova Sonic} no Bedrock (\cd{amazon.nova-sonic-v1:0}), modelo \emph{speech-to-speech} bidirecional via stream HTTP/2 (\cd{InvokeModelWithBidirectionalStream}), suportando barge-in e transcrição. Exige reescrever o adaptador em \arq{mascote-live.gateway.ts}.
    \item \textbf{Na Azure (\nativo):} o equivalente direto é a \textbf{Azure OpenAI GPT-4o Realtime API} (\cd{gpt-4o-realtime-preview} / \cd{gpt-4o-mini-realtime-preview}), que também opera sobre WebSockets com áudio PCM16 a 24 kHz, detecção automática de atividade de voz (\cd{server\_vad}) e eventos JSON muito similares aos do Gemini Live.
  \end{itemize}
  \item[Avatar em Vídeo] O vídeo do avatar humano (\cd{AVATAR\_AO\_VIVO}, ADR-132) é processado pelo SaaS externo \textbf{HeyGen LiveAvatar} (Modo FULL via WebRTC), sendo \textbf{100\% agnóstico de nuvem}. Os servidores da HeyGen chamam \cd{POST https://api.<dominio>/ajuda/llm/chat/completions} (SSE); basta que o balanceador da AWS ou Azure não faça bufferização dos chunks SSE para que a \cd{FALA\_DE\_ESPERA} de 7 segundos funcione nas três nuvens.
\end{campos}
\end{pecaarq}

% -----------------------------------------------------------------------------
\begin{pecaarq}{IA-05}{Governança de Quotas (\texttt{ia\_quotas}) e Painel de Custos (ADR-105 / ADR-161)}{\sgcp~\nativo\enspace\saws~\nativo\enspace\sazure~\nativo}
\begin{campos}
  \item[Código na API] \arq{apps/api/src/ai/quota.service.ts}, \arq{medidor.provider.ts} e \arq{apps/api/src/unooria/uso-ia.service.ts}.
  \item[Tabelas PostgreSQL] \cd{ia\_quotas} (tetos diários e preços por 1M de tokens configuráveis pelo \cd{admin}), \cd{ia\_uso\_diario}, \cd{ia\_chamadas\_log} e \cd{ajuda\_limites} (minutos de voz/avatar por escola).
  \item[Portabilidade] Por decisão arquitetural (ADR-105 e ADR-161), toda a governança de consumo de IA da Plataforma Unooria foi construída \textbf{dentro da própria aplicação e do PostgreSQL}, e não em serviços proprietários de faturamento da GCP. Ao migrar para AWS (Bedrock) ou Azure (Foundry/OpenAI), \textbf{o painel de custos de IA e as travas de quota por estudante/professor continuam funcionando sem alterar uma única linha de banco ou de frontend}.
\end{campos}
\end{pecaarq}

% =============================================================================
\section{Síntese Comparativa do Eixo Uso de IA}

\begin{table}[H]
\centering
\begin{tabularx}{\textwidth}{@{}L{32mm}YYY@{}}
\toprule
\headrow
\thd{Capacidade de IA} &
\thd{GCP (Vertex AI)} &
\thd{AWS (Amazon Bedrock)} &
\thd{Azure (AI Foundry / OpenAI)} \\
\midrule
\textbf{LLM Rápido + Planos Longos} &
\cd{gemini-3.6-flash} e \cd{3.7-flash} (saída até \textbf{65.536} tok; \cd{thinkingBudget}) &
\cd{Nova Pro/Lite} e \cd{Claude 3.7/4 Sonnet} (saída até \textbf{64.000} tok) &
\cd{GPT-4.1} / \cd{4.1-mini} / \cd{o4-mini} (saída até \textbf{32.768} tok no GPT-4.1) \\
\zebra
\textbf{Cache de Contexto (65k tok Ajuda)} &
\textbf{Implícito e automático} (75--90\% desc. em prefixos repetidos) &
\textbf{Explícito} (requer inserir marcador \cd{cachePoint} no prompt) &
\textbf{Implícito e automático} (50--75\% desc. acima de 1.024 tok) \\
\textbf{Geração de Imagem \cd{4:5}} &
\cd{gemini-3.1-flash-image} (\cd{aspectRatio: '4:5'} nativo) &
\cd{Amazon Nova Canvas} (resolução \cd{1024x1280} nativa \cd{4:5}) &
\cd{gpt-image-1} (\cd{1024x1536} + crop automático \cd{4:5} via \cd{sharp}) \\
\zebra
\textbf{TTS Diálogo Multi-Locutor} &
\cd{gemini-2.5-flash-tts} (\cd{multiSpeakerMarkup} em 1 chamada) &
\cd{Amazon Polly} (exige 1 chamada por fala + concatenação MP3) &
\cd{Azure AI Speech} (múltiplas tags \cd{<voice>} em 1 chamada SSML) \\
\textbf{Runtime do Agente + Sessão} &
\textbf{Vertex AI Agent Engine} (\cd{BYOC} + Google ADK + sessão nativa sem DB) &
\textbf{Bedrock AgentCore} ou ECS Fargate (portar ADK \seta{} Strands/LangGraph) &
\textbf{AI Agent Service} ou ACA interno (portar ADK \seta{} LangGraph/SK) \\
\zebra
\textbf{RAG Híbrido (BNCC + PDF \cd{pageSpan})} &
\textbf{Vertex AI Search} (Layout Parser nativo + \textbf{US\$ 0 fixo/mês}) &
\textbf{Bedrock KB} (OpenSearch custa \textbf{US\$ 175--350/mês}; ou \cd{pgvector} no RDS) &
\textbf{Azure AI Search} + Doc Intelligence (Basic custa \textbf{US\$ 75/mês}; ou \cd{pgvector}) \\
\textbf{Voz Ao Vivo (Mascote 3D)} &
\textbf{Vertex Gemini Live} (\cd{gemini-3.8-live}, WSS PCM 24kHz) &
\textbf{Amazon Nova Sonic} (stream HTTP/2 bidirecional PCM) &
\textbf{GPT-4o Realtime API} (WSS / WebRTC PCM 24kHz) \\
\bottomrule
\end{tabularx}
\caption{Comparativo funcional e econômico dos subsistemas de Inteligência Artificial nas três nuvens.}
\label{tab:comparativo-eixo-ia}
\end{table}

% =============================================================================
% CAPÍTULO 7 — Custos (TCO), Operação Solo e Matriz de Migração (Lock-in)
% =============================================================================
\chapter{Custos (TCO), Operação Solo e Matriz de Migração}
\label{cap:custos-migracao}

\begin{fichaeixo}{Ficha de Síntese Decisória --- TCO, Operação Solo e Esforço de Migração}
\begin{campos}
  \item[Unidade econômica] 1 contrato público (1 stack dedicada no Terraform) atendendo $\sim \num{10000}$ estudantes (pico de \num{2000} a \num{4000} simultâneos).
  \item[Regiões de referência] GCP \cd{southamerica-east1} (SP) + \cd{us-central1} (IA) $\mid$ AWS \cd{sa-east-1} (SP) + \cd{us-east-1} (IA) $\mid$ Azure \cd{brazilsouth} (SP) + \cd{eastus2} (IA).
  \item[Recomendação executiva] \textbf{Permanecer na GCP como nuvem padrão.} Em editais B2G com exigência contratual estrita de AWS ou Azure, adotar a arquitetura com \cd{pgvector} no PostgreSQL existente para evitar o custo fixo de busca vetorial dedicada.
\end{campos}
\end{fichaeixo}

Em contratos B2G educacionais, o custo de infraestrutura por contrato tem dois componentes muito diferentes: o \textbf{piso fixo mensal de prontidão} (aquilo que a nuvem cobra 24x7 apenas para manter a stack isolada do município/estado ligada, com banco HA e réplicas mínimas aquecidas, mesmo durante férias escolares ou domingos) e o \textbf{custo variável de tráfego e IA} (proporcional ao uso efetivo pelos estudantes e professores).

Como cada contrato público ganha uma stack fisicamente isolada (um projeto/conta/resource group próprio), qualquer serviço que cobre uma mensalidade fixa mínima elevada multiplica esse custo pelo número de municípios/estados atendidos.

% =============================================================================
\section{Estimativa Comparada de TCO Mensal por Contrato (Produção)}
\label{sec:tco-mensal}

A Tabela~\ref{tab:tco-comparado} projeta o custo mensal estimado (em dólares americanos, preços públicos sob demanda sem descontos de \emph{Committed Use} / \emph{Savings Plans}) para \textbf{1 contrato em produção} nas regiões de São Paulo (\cd{southamerica-east1}, \cd{sa-east-1} e \cd{brazilsouth}) para computação/dados e EUA para os endpoints de IA generativa.

{\footnotesize
\begin{tabularx}{\textwidth}{@{}L{36mm}R{23mm}R{25mm}R{25mm}Y@{}}
\caption{Comparativo estimado de TCO mensal (em US\$) para 1 contrato B2G de 10.000 estudantes em produção.}
\label{tab:tco-comparado} \\
\toprule
\headrow
\thd{Camada / Recurso (Prd)} &
\thd{GCP (US\$/mês)} &
\thd{AWS (US\$/mês)} &
\thd{Azure (US\$/mês)} &
\thd{Justificativa da Diferença} \\
\midrule
\endfirsthead
\toprule
\headrow
\thd{Camada / Recurso (Prd)} &
\thd{GCP (US\$/mês)} &
\thd{AWS (US\$/mês)} &
\thd{Azure (US\$/mês)} &
\thd{Justificativa da Diferença} \\
\midrule
\endhead
\bottomrule
\endfoot
\textbf{1. Borda, LB Global e WAF} &
$\sim\text{US\$ } 32$ &
$\sim\text{US\$ } 48$ &
$\sim\text{US\$ } 55$ &
AWS exige CloudFront WAF + ALB regional; Azure Front Door Standard (\$\,35 base + WAF policy). \\
\zebra
\textbf{2. Rede Privada e Egress (NAT)} &
$\sim\text{US\$ } 0\text{--}15$ &
$\sim\text{US\$ } 65\text{--}110$ &
$\sim\text{US\$ } 45$ &
GCP Direct VPC Egress dispensa NAT para serviços Google; AWS NAT Gateway cobra por hora por AZ em SP. \\
\textbf{3. Hub Next.js (\cd{min=1} a \cd{5})} &
$\sim\text{US\$ } 38$ &
$\sim\text{US\$ } 52$ &
$\sim\text{US\$ } 42$ &
1 vCPU / 1--2 GiB sempre aquecido + picos diários das 7h30 às 17h30 em dias letivos. \\
\zebra
\textbf{4. API NestJS (\cd{2 vCPU}, \cd{min=1} a \cd{5})} &
$\sim\text{US\$ } 95$ &
$\sim\text{US\$ } 135$ &
$\sim\text{US\$ } 110$ &
CPU 100\% alocada (\cd{cpu\_idle=false}); Fargate exige mínimo de 4 GB RAM para 2 vCPU em \cd{sa-east-1}. \\
\textbf{5. PostgreSQL 16 HA (2 vCPU, 8 GiB)} &
$\sim\text{US\$ } 290$ &
$\sim\text{US\$ } 310$ &
$\sim\text{US\$ } 285$ &
Instância dedicada com failover automático entre zonas (Regional HA / Multi-AZ) + 50 GB SSD + backups 30d. \\
\zebra
\textbf{6. Storage + CDN (4 buckets)} &
$\sim\text{US\$ } 18$ &
$\sim\text{US\$ } 22$ &
$\sim\text{US\$ } 20$ &
Bundles estáticos com cache longo (1 ano) e mídia otimizada em \cd{.webp}. \\
\textbf{7. Fila de Telemetria + Logs} &
$\sim\text{US\$ } 12$ &
$\sim\text{US\$ } 18$ &
$\sim\text{US\$ } 20$ &
Pub/Sub + Cloud Logging (franquia inicial gratuita de 50 GiB/mês por projeto na GCP). \\
\zebra
\textbf{8. Runtime Agente Professor (\cd{min=1})} &
$\sim\text{US\$ } 115$ &
$\sim\text{US\$ } 110$ &
$\sim\text{US\$ } 95$ &
Contêiner Python aquecido (Agent Engine na GCP vs. ECS Fargate na AWS vs. ACA na Azure). \\
\textbf{9. Índice RAG (BNCC + Acervo PDF)} &
\textbf{$\sim\text{US\$ } 2$} &
\textbf{$\sim\text{US\$ } 350$} \newline {\scriptsize (ou US\$ 5 c/ \cd{pgvector})} &
\textbf{$\sim\text{US\$ } 75$} \newline {\scriptsize (ou US\$ 5 c/ \cd{pgvector})} &
Vertex AI Search é 100\% serverless por query/GB; OpenSearch Serverless tem piso de OCUs; AI Search Basic é fixo. \\
\midrule
\rowcolor{unooBlueWash}
\textbf{Subtotal Infraestrutura Fixa/Base} &
\textbf{$\sim\text{US\$ } 602$} &
\textbf{$\sim\text{US\$ } 1.110$} \newline {\scriptsize (\textbf{US\$ 765} c/ \cd{pgvector})} &
\textbf{$\sim\text{US\$ } 747$} \newline {\scriptsize (\textbf{US\$ 677} c/ \cd{pgvector})} &
Na GCP, o piso mensal por contrato é entre \textbf{20\% e 45\% menor} que nas concorrentes. \\
\midrule
\textbf{10. Consumo Variável de IA (10k alunos)} &
$\sim\text{US\$ } 180\text{--}320$ &
$\sim\text{US\$ } 210\text{--}390$ &
$\sim\text{US\$ } 220\text{--}410$ &
Governado pelo teto em \cd{ia\_quotas} e \cd{ajuda\_limites}; Gemini Flash e cache implícito reduzem custo na GCP. \\
\midrule
\headrow
\thd{TCO Total Estimado / Contrato} &
\thd{US\$ 780 -- 920} &
\thd{US\$ 975 -- 1.500} &
\thd{US\$ 895 -- 1.155} &
\thd{Vantagem estrutural da GCP em multi-stack B2G} \\
\end{tabularx}
}

\subsection{Três lições econômicas da tabela de TCO}

\begin{enumerate}
  \item \textbf{O efeito multiplicador em ambientes de ``Desenvolvimento'' e ``Homologação'':}
  No ambiente \cd{dev} da GCP (\arq{envs/dev.tfvars}), o banco usa tier compartilhado (\cd{db-f1-micro}, sem HA), os contêineres permitem \cd{min\_instances = 0} ou \cd{1} leve, não há NAT Gateway pago por hora e o Vertex AI Search custa zero parado. Isso mantém o ambiente \cd{dev} da GCP em cerca de \textbf{US\$ 45 a US\$ 70/mês}. Na AWS, apenas ligar uma VPC com NAT Gateway Multi-AZ, um ALB e uma coleção OpenSearch Serverless em um ambiente de homologação já consumiria mais de \textbf{US\$ 400/mês antes do primeiro login}.

  \item \textbf{Em AWS ou Azure, \cd{pgvector} no PostgreSQL é obrigatório para proteger a margem:}
  Se um edital público exigir contratação em nuvem AWS ou Azure, \textbf{não se deve utilizar o Amazon OpenSearch Serverless} para o Agente do Professor. Ao armazenar os embeddings e metadados de página das habilidades BNCC e dos PDFs em tabelas \cd{pgvector} dentro do próprio RDS PostgreSQL / Azure Database for PostgreSQL já pago pelo contrato, economizam-se de US\$ 900 a US\$ 4.200 por ano por contrato.

  \item \textbf{Direct VPC Egress vs. NAT Gateways:}
  Na GCP, como a API e o Hub se comunicam apenas com o Cloud SQL (IP privado na VPC) e com APIs do Google (Vertex AI, Cloud Storage, Pub/Sub, Secret Manager --- todas alcançáveis via \emph{Private Google Access} gratuito), o custo de gateway NAT é virtualmente zero. Na AWS, a menos que se provisionem múltiplos \emph{VPC Interface Endpoints (PrivateLink)} --- que também cobram por hora e por AZ ---, todo o tráfego da API em subnet privada para o Bedrock, S3, SQS e HeyGen passa pelo \textbf{NAT Gateway} tarifado por hora e por GB processado.
\end{enumerate}

% =============================================================================
\section{Operação por Uma Única Pessoa (Carga Cognitiva)}
\label{sec:operacao-solo}

A premissa organizacional da Plataforma Unooria (\emph{``Prazo: produção em novembro/2026. Uma pessoa desenvolve. Simplicidade $>$ sofisticação''}) torna a complexidade de manutenção tão importante quanto o custo financeiro. A Tabela~\ref{tab:carga-operacional} avalia o esforço diário de operação nos três provedores.

{\footnotesize
\begin{tabularx}{\textwidth}{@{}L{35mm}YYY@{}}
\caption{Comparativo de carga cognitiva e complexidade operacional para um único desenvolvedor.}
\label{tab:carga-operacional} \\
\toprule
\headrow
\thd{Dimensão Operacional} &
\thd{GCP (Atual)} &
\thd{AWS} &
\thd{Azure} \\
\midrule
\endfirsthead
\toprule
\headrow
\thd{Dimensão Operacional} &
\thd{GCP (Atual)} &
\thd{AWS} &
\thd{Azure} \\
\midrule
\endhead
\bottomrule
\endfoot
\textbf{Quantidade de recursos no Terraform por contrato} &
\textbf{$\sim 45$ recursos} (1 LB global, 2 Cloud Run services, 1 Job, 1 SQL, 4 buckets, 1 Pub/Sub, Agent Engine) &
\textbf{$\sim 85$ a $100$ recursos} (VPC, 4+ Subnets, Route Tables, IGW, NAT GW, Security Groups, CloudFront, ALB, Target Groups, ECS Cluster, Task Defs, EventBridge Pipe) &
\textbf{$\sim 55$ a $65$ recursos} (VNet, Subnets, Front Door, ACA Environment, 2 Container Apps, 1 ACA Job, Flexible Server, 4 Containers Blob, Event Grid) \\
\zebra
\textbf{Túnel de manutenção ao banco de dados} &
\textbf{1 comando CLI} (\cd{gcloud compute ssh --tunnel-\B through-iap}) sem gerenciar chaves SSH &
\textbf{1 comando CLI} (\cd{aws ssm start-session} c/ port forwarding) sem chaves SSH &
Exige subir VM de salto + \textbf{Azure Bastion} ou abrir regra temporária de firewall \\
\textbf{Acesso da API às ferramentas de IA e Storage} &
\textbf{1 único mecanismo IAM} (ADC via Service Account da revisão Cloud Run para SQL, GCS, Pub/Sub e Vertex AI) &
\textbf{IAM Task Role} unificada (SigV4 para S3, SQS e Bedrock), mas exige conciliar Security Groups de rede + políticas de bucket &
\textbf{Managed Identity} unificada (Entra ID RBAC para Blob, Event Grid e Azure OpenAI) \\
\zebra
\textbf{Deploy e Rollback de código backend} &
Revisões imutáveis nativas no Cloud Run (rollback de tráfego em \textbf{5 segundos}) &
Nova revisão de Task Definition + rolling update no ECS (\textbf{2 a 4 minutos}) &
Revisões imutáveis nativas no Container Apps (rollback em \textbf{10 segundos}) \\
\end{tabularx}
}

% =============================================================================
\section{Matriz de Portabilidade e Esforço de Migração (Lock-in)}
\label{sec:matriz-migracao}

Para responder com precisão à pergunta \emph{``quanto custaria, em dias de engenharia, portar a Plataforma Unooria hoje para a AWS ou para a Azure?''}, a Tabela~\ref{tab:matriz-portabilidade} classifica cada parte do sistema pelo seu grau real de acoplamento à GCP.

{\footnotesize
\begin{tabularx}{\textwidth}{@{}L{34mm}L{25mm}YR{22mm}@{}}
\caption{Matriz de acoplamento (lock-in) e esforço estimado de engenharia para migrar a plataforma para AWS ou Azure.}
\label{tab:matriz-portabilidade} \\
\toprule
\headrow
\thd{Componente / Repo} &
\thd{Acoplamento} &
\thd{O que precisa ser alterado para rodar em AWS ou Azure} &
\thd{Esforço} \\
\midrule
\endfirsthead
\toprule
\headrow
\thd{Componente / Repo} &
\thd{Acoplamento} &
\thd{O que precisa ser alterado para rodar em AWS ou Azure} &
\thd{Esforço} \\
\midrule
\endhead
\bottomrule
\endfoot
\textbf{Hub Web} \newline (\cd{apps/hub}, Next.js, OIDC/JWKS) &
\textbf{Zero} \newline (100\% portável) &
Nenhuma linha de código. O Dockerfile standalone, a assinatura RS256 dos tokens de experiência e o cliente OIDC funcionam identicamente em ECS Fargate ou Azure Container Apps. &
\textbf{0 dias} \\
\zebra
\textbf{Experiências e Pacotes} \newline (\cd{apps/experiencias/*}, \cd{sdk}, \cd{ui}) &
\textbf{Zero} \newline (100\% portável) &
Nenhuma linha de código nas experiências. Apenas ajustar os scripts \arq{scripts/publicar-\B experiencia.mjs} e \arq{publicar-\B manifesto.mjs} para suportar \cd{aws s3 sync} ou \cd{az storage blob sync} além de \cd{gcloud storage}. &
\textbf{0,5 dia} \\
\textbf{Banco, RLS e Prisma} \newline (\cd{apps/api/prisma}) &
\textbf{Zero} \newline (100\% portável) &
Nenhuma linha de SQL ou schema. O PostgreSQL 16 puro com Row-Level Security (\cd{app.current\_escola\_id}) e \cd{comTenancy} roda sem alterações no Amazon RDS ou no Azure Database for PostgreSQL. &
\textbf{0 dias} \\
\zebra
\textbf{Storage e Telemetria} \newline (\cd{storage}, \cd{telemetry}) &
\textbf{Baixo} \newline (Isolado em serviços) &
Substituir o SDK \cd{@google-cloud/storage} pelo \cd{@aws-sdk/client-s3} (ou \cd{@azure/storage-blob}) nos módulos de planos/geração, e adaptar o publicador de telemetria (\cd{@google-cloud/pubsub} \seta{} SQS ou Event Grid). O endpoint receptor \cd{POST /telemetry/worker} permanece igual. &
\textbf{1,5 a 2 dias} \\
\textbf{Camada de Modelos IA} \newline (\cd{apps/api/src/ai/}) &
\textbf{Baixo a Médio} \newline (Plugável via DI) &
Criar \cd{BedrockAiProvider} ou \cd{AzureOpenAiProvider} implementando a interface existente \cd{AiProvider}. O \cd{MedidorDeProvider} e o \cd{QuotaService} não mudam. Na AWS, adicionar concatenação de turnos para o Polly TTS; na Azure, adicionar crop \cd{4:5} no \cd{sharp} para imagens. &
\textbf{2 a 3 dias} \\
\zebra
\textbf{Mascote Voz Ao Vivo} \newline (\cd{mascote-live.\B gateway.ts}) &
\textbf{Médio} \newline (Protocolo WSS) &
Adaptar o proxy WebSocket em \arq{apps/api/src/ajuda/mascote-live.gateway.ts} para traduzir os pacotes de áudio PCM 16/24 kHz para a \textbf{GPT-4o Realtime API} (muito similar na Azure) ou para o stream bidirecional do \textbf{Amazon Nova Sonic} (na AWS). &
\textbf{2 a 3 dias} \\
\textbf{Infra como Código} \newline (Repo \cd{Fundacao}) &
\textbf{Alto} \newline (Específico da nuvem) &
Escrever um novo módulo raiz de Terraform para AWS ou Azure espelhando as variáveis de \arq{variables.tf} e os ambientes \cd{dev}/\cd{prd}, mais atualização dos workflows de deploy no GitHub Actions (OIDC federado). &
\textbf{5 a 7 dias} \\
\zebra
\textbf{Agente do Professor} \newline (\cd{apps/agente-\B professor}) &
\textbf{Alto} \newline (ADK + Agent Engine + Search) &
Portar a definição do agente do Google ADK para LangGraph / Strands / Semantic Kernel, reescrever \arq{agent-engine.client.ts} na API e migrar a ingestão/consulta das habilidades BNCC e dos PDFs paginados para \textbf{Azure AI Search} ou \textbf{PostgreSQL \cd{pgvector}}. &
\textbf{5 a 8 dias} \\
\midrule
\rowcolor{unooBlueWash}
\textbf{Total para Migração Completa} &
\textbf{Arquitetura desacoplada} &
85\% do código do monorepo (Hub, Experiências, SDK, Design System, Regras de Negócio e Banco RLS) possui \textbf{zero lock-in}. Todo o esforço concentra-se no Terraform, nos adapters de IA/Storage e no Agente do Professor. &
\textbf{16 a 23 dias úteis} \\
\end{tabularx}
}

% =============================================================================
\section{Conclusão e Diretriz Arquitetural}
\label{sec:conclusao-comparativo}

A análise detalhada dos quatro eixos confirma que \textbf{a escolha da Google Cloud Platform (GCP) para a Plataforma Unooria é tecnicamente e economicamente superior para o contexto atual do negócio}:

\begin{enumerate}
  \item \textbf{Aderência nativa aos 4 eixos com menor complexidade operacional:} o par \textbf{Cloud Run v2} (\cd{cpu\_idle = false}, \cd{timeout = 3600s}, scale-out em segundos) + \textbf{Cloud Run Jobs} resolve tanto o Hub quanto o backend monolítico da API sem a sobrecarga de rede e lentidão de escalonamento do ECS Fargate na AWS.
  \item \textbf{Vantagem decisiva no Eixo de IA e RAG:} a combinação de \textbf{Vertex Gemini 3.6/3.7 Flash} (saída de 64k tokens, controle de \cd{thinkingBudget}, cache implícito de 65k tokens, imagens \cd{4:5} nativas e TTS multi-locutor em chamada única), \textbf{Vertex AI Agent Engine} (hospedagem do agente Python com sessão embutida sem banco extra), \textbf{Vertex AI Search} (Layout Parser com \cd{pageSpan} e custo fixo zero) e \textbf{Vertex Gemini Live} entrega capacidades que, na AWS ou na Azure, exigiriam costurar múltiplos serviços distintos e assumir pisos mensais fixos mais altos.
  \item \textbf{Portabilidade preservada como seguro estratégico:} graças às decisões de engenharia já tomadas no monorepo --- PostgreSQL puro com RLS, autenticação OIDC padrão, experiências estáticas desacopladas e interface \cd{AiProvider} plugável ---, cerca de \textbf{85\% da base de código não sabe em qual nuvem está rodando}. Caso um contrato governamental específico exija implantação em \textbf{Azure} (a segunda colocada em aderência operacional, graças ao \textbf{Azure Container Apps} e à \textbf{GPT-4o Realtime API}) ou em \textbf{AWS} (via \textbf{ECS Fargate} + \textbf{RDS PostgreSQL com \cd{pgvector}}), a migração completa é executável em cerca de 3 a 4 semanas de engenharia sem alterar nenhuma tela, experiência ou regra de negócio.
\end{enumerate}


\appendix
% =============================================================================
% APÊNDICE A — Código-Fonte Mermaid dos Diagramas
% =============================================================================
\chapter{Código-Fonte Mermaid dos Diagramas}
\label{ap:mermaid}

Em conformidade com a convenção da documentação técnica da Plataforma Unooria, este apêndice preserva a representação em \textbf{Mermaid} equivalente aos diagramas \textbf{TikZ} renderizados nos capítulos anteriores, facilitando a manutenção futura por agentes de código e ferramentas externas.

% =============================================================================
\section{Topologia Geral Comparada (Figura~\ref{fig:topologia-comparada})}

\begin{lstlisting}[language=mermaid, label={lst:mermaid-topologia}]
flowchart TB
  subgraph GCP["1. GCP (Arquitetura Atual)"]
    G_EDGE["Cloud DNS + Cloud Armor WAF + Global HTTPS LB"]
    G_CDN["Cloud CDN + 4 Buckets GCS (-exp, -assets, -planos, -mapas)"]
    G_RUN["Cloud Run v2: Hub (Next.js) + API (NestJS 2vCPU/2GiB) + Job (migrate)"]
    G_ASYNC["Cloud Pub/Sub (Push OIDC) -> POST /telemetry/worker"]
    G_DB["Cloud SQL PostgreSQL 16 (HA Regional + RLS + Private IP)"]
    G_AI["Vertex AI: Gemini 3.6/3.7 + Agent Engine (ADK) + Vertex AI Search + Gemini Live"]
    G_EDGE --> G_CDN
    G_EDGE --> G_RUN
    G_RUN --> G_ASYNC
    G_RUN --> G_DB
    G_RUN --> G_AI
  end

  subgraph AWS["2. AWS (Equivalência)"]
    A_EDGE["Route 53 + AWS WAF + CloudFront + ALB"]
    A_CDN["CloudFront OAC + 4 Buckets S3"]
    A_RUN["ECS Fargate: Hub + API (2vCPU/4GB) + Standalone Task (migrate)"]
    A_ASYNC["SQS + EventBridge Pipe (API Destination) -> POST /telemetry/worker"]
    A_DB["RDS PostgreSQL 16 Multi-AZ + RDS Proxy + RLS"]
    A_AI["Amazon Bedrock: Claude/Nova + AgentCore + RDS pgvector + Nova Sonic"]
    A_EDGE --> A_CDN
    A_EDGE --> A_RUN
    A_RUN --> A_ASYNC
    A_RUN --> A_DB
    A_RUN --> A_AI
  end

  subgraph AZURE["3. Azure (Equivalência)"]
    Z_EDGE["Azure DNS + Front Door WAF + Application Gateway"]
    Z_CDN["Front Door CDN + 4 Containers Blob Storage"]
    Z_RUN["Azure Container Apps (ACA): Hub + API + ACA Job (migrate)"]
    Z_ASYNC["Event Grid / Storage Queue -> POST /telemetry/worker"]
    Z_DB["Azure Database for PostgreSQL 16 Flexible Server (Zone-Redundant + PgBouncer)"]
    Z_AI["Azure AI Foundry: GPT-4.1 + ACA/Agent Service + AI Search/pgvector + GPT-4o Realtime"]
    Z_EDGE --> Z_CDN
    Z_EDGE --> Z_RUN
    Z_RUN --> Z_ASYNC
    Z_RUN --> Z_DB
    Z_RUN --> Z_AI
  end
\end{lstlisting}

% =============================================================================
\section{Topologia Comparada da Stack de IA (Figura~\ref{fig:topologia-ia-comparada})}

\begin{lstlisting}[language=mermaid]
flowchart LR
  subgraph GCP_IA["GCP - Stack de IA Atual"]
    G1["AiProvider: Vertex Gemini 3.6/3.7 Flash + Flash Image (4:5) + Flash TTS Multi-voz"]
    G2["Agente Professor: Vertex AI Agent Engine (BYOC Python + Google ADK + Sessoes Nativas)"]
    G3["RAG Duplo: Vertex AI Search (NO_CONTENT p/ BNCC + Layout Parser pageSpan p/ PDFs)"]
    G4["Voz Ao Vivo: Vertex Gemini Live (gemini-3.8-live WSS Bidi PCM 24kHz)"]
    G1 --> G2 --> G3 --> G4
  end

  subgraph AWS_IA["AWS - Equivalencia de IA"]
    A1["AiProvider: Bedrock (Claude/Nova) + Nova Canvas (4:5) + Polly (concatenar turnos)"]
    A2["Agente Professor: Bedrock AgentCore ou ECS Fargate (Strands/LangGraph + DynamoDB/RDS)"]
    A3["RAG Duplo: RDS PostgreSQL pgvector (evita custo fixo do OpenSearch Serverless)"]
    A4["Voz Ao Vivo: Amazon Nova Sonic (Bedrock Bidirectional Stream PCM)"]
    A1 --> A2 --> A3 --> A4
  end

  subgraph AZ_IA["Azure - Equivalencia de IA"]
    Z1["AiProvider: Azure OpenAI (GPT-4.1) + gpt-image-1 (crop sharp 4:5) + AI Speech SSML"]
    Z2["Agente Professor: ACA Interno ou AI Agent Service (LangGraph/SK + PostgreSQL)"]
    Z3["RAG Duplo: Azure AI Search + Document Intelligence (Layout) ou PostgreSQL pgvector"]
    Z4["Voz Ao Vivo: Azure OpenAI GPT-4o Realtime API (WSS / WebRTC PCM 24kHz)"]
    Z1 --> Z2 --> Z3 --> Z4
  end
\end{lstlisting}


\end{document}
